\documentclass[11pt]{article}
\usepackage[utf8]{inputenc}
\usepackage[margin=1in]{geometry}
\usepackage{amsmath,amssymb,amsthm,mathtools}
\usepackage{microtype}
\usepackage{xcolor}
\usepackage[unicode,colorlinks=true,linkcolor=blue!45!black,citecolor=blue!45!black,urlcolor=blue!45!black]{hyperref}
\hypersetup{
 pdftitle={QAOA attains the SK ground-state energy in the thermodynamic-first limit},
 pdfauthor={OpenAI},
 pdfsubject={A proof with the thermodynamic limit taken before the QAOA depth limit},
 pdfkeywords={QAOA, Sherrington-Kirkpatrick, Parisi formula, spin glass, quantum optimization}
}
\numberwithin{equation}{section}
\newtheorem{theorem}{Theorem}[section]
\newtheorem{proposition}[theorem]{Proposition}
\newtheorem{lemma}[theorem]{Lemma}
\newtheorem{corollary}[theorem]{Corollary}
\theoremstyle{definition}
\newtheorem{definition}[theorem]{Definition}
\newtheorem{assumption}[theorem]{Assumption}
\theoremstyle{remark}
\newtheorem{remark}[theorem]{Remark}
\newcommand{\E}{\mathbb E}
\newcommand{\R}{\mathbb R}
\newcommand{\C}{\mathbb C}
\newcommand{\N}{\mathbb N}
\newcommand{\1}{\mathbf 1}
\newcommand{\ii}{\mathrm i}
\newcommand{\dd}{\mathrm d}
\newcommand{\vac}{\boldsymbol\Omega}
\newcommand{\cH}{\mathcal H}
\newcommand{\cF}{\mathcal F}
\newcommand{\cE}{\mathcal E}
\newcommand{\cL}{\mathsf L}
\providecommand{\mathscr}[1]{\mathcal{#1}}
\DeclareMathOperator{\supp}{supp}
\DeclareMathOperator{\dist}{dist}
\DeclareMathOperator{\im}{Im}
\DeclareMathOperator{\re}{Re}
\DeclareMathOperator{\Span}{span}
\DeclareMathOperator{\ad}{ad}
\DeclareMathOperator{\Tr}{Tr}
\newcommand{\ket}[1]{\lvert #1\rangle}
\newcommand{\bra}[1]{\langle #1\rvert}
\newcommand{\braket}[2]{\langle #1,#2\rangle}
\newcommand{\norm}[1]{\lVert #1\rVert}
\newcommand{\abs}[1]{\lvert #1\rvert}
\newcommand{\Pstar}{P_*}

\title{QAOA attains the SK ground-state energy\\
in the thermodynamic-first limit}
\author{OpenAI}
\date{September 25, 2026}
\setcounter{tocdepth}{2}
\makeatletter
\renewcommand{\maketitle}{%
  \begin{center}
    {\LARGE\@title\par}
    \vspace{1em}
    {\normalsize\@author\par}
    \vspace{.4em}
    {\normalsize\@date\par}
  \end{center}
  \vspace{.5em}
}
\makeatother

\ifdefined\pdfinfoomitdate\pdfinfoomitdate=1\fi
\ifdefined\pdftrailerid\pdftrailerid{}\fi

\begin{document}
\maketitle

\begin{abstract}
We prove that the Quantum Approximate Optimization Algorithm (QAOA)
approaches the ground-state energy per spin of the Gaussian zero-field
Sherrington--Kirkpatrick model when system size tends to infinity first
and circuit depth then increases. For every accuracy, some finite depth
and deterministic angles, independent of system size and disorder,
achieve that accuracy in the limiting expected energy per spin using
the standard cost Hamiltonian and transverse-field mixer. This proves
the eventual Parisi-optimality conjecture of Basso, Farhi, Marwaha,
Villalonga, and Zhou in its fixed-parameter thermodynamic formulation.
We give no quantitative bound on the required depth or efficient
angle-selection procedure.
\end{abstract}

\tableofcontents

\section{Introduction and statement of the result}\label{sec:introduction}

The Sherrington--Kirkpatrick (SK) model was introduced by Sherrington and
Kirkpatrick~\cite{SK1975}
as a mean-field Ising spin glass with independent Gaussian pair
interactions. Parisi's replica-symmetry-breaking description uses a
hierarchy of order parameters~\cite{Parisi1979}. Guerra's interpolation
established the Parisi upper bound for the pressure~\cite{Guerra2003},
and Talagrand proved equality~\cite{Talagrand2006}.

Farhi, Goldstone, and Gutmann introduced the quantum approximate
optimization algorithm (QAOA) as alternating evolution under an objective
Hamiltonian and a transverse-field mixer, starting from the uniform
superposition~\cite{FGG2014}. For this model, Farhi, Goldstone, Gutmann, and Zhou
developed an exact description of its performance at every fixed depth in
the infinite-size limit and proved concentration at fixed
parameters~\cite{FGGZ2022}. Basso, Farhi, Marwaha, Villalonga,
and Zhou related that limit to QAOA on large-degree trees and conjectured
that increasing depth eventually attains the Parisi optimum
\cite[Section~7, Equation~(7.2)]{BFMVZ2022}.
Boulebnane et al.~\cite{BKL2026} proved a fixed-depth spin--boson
representation and used it to obtain numerical evidence for convergence
towards the Parisi value. Their representation provides a useful
Hilbert-space language for the limit. The construction here must also
realize bounded classical computations and eliminate longitudinal seed
controls before selecting one ordinary word.

Sels and Morone~\cite{SM2026} studied QAOA on the SK model using a
semiclassical truncated-Wigner approximation and compared its performance
with quantum data. Their numerical results and heuristic analysis give
further evidence for convergence to the Parisi value and suggest a depth
scaling.

The order of limits distinguishes this question from the fixed-instance
optimality proved by Farhi, Goldstone, and Gutmann~\cite[Section~VI,
Equation~(10)]{FGG2014}:
for each fixed instance, increasing the depth approaches its optimum.
That result does not provide a finite depth that works as the number
of spins tends to infinity. We instead choose a complete finite word for
each desired accuracy and then let the number of spins tend to
infinity. The quantum construction simulates a finite classical
calculation, using temporary Gaussian seed controls and then
removing them.

There are two ways to supply the classical calculation. Auffinger,
Chen, and Zeng proved that the zero-temperature Parisi order parameter
has infinitely many support points~\cite[Theorem~1]{ACZ2020}; Chen
proved full support on $[0,1)$ and a smooth
density~\cite[Theorem~1.3]{Chen2026}.
Our primary route uses the companion
\emph{Full support of the zero-temperature Sherrington--Kirkpatrick
order parameter}~\cite{SKSupport2026}. Its scalar diffusion yields a
finite function of independent Gaussian variables and a coefficient
formula for its value. We realize that function directly in the
rooted Hilbert space and identify the coefficient formula with its
energy. The imported scalar facts are full support of the zero-temperature
order parameter, the diffusion's consistency and value identities,
and a finite Gaussian coefficient identity.
Section~\ref{sec:classical-optimization} states these inputs before
constructing their rooted realization.

We also give an independent positive-temperature route.
It follows the Parisi-diffusion approach of
Montanari~\cite{Montanari2019}, whose optimization theorem assumes
interval support of the positive-temperature Parisi measure, and the
broader framework of El Alaoui, Montanari, and Sellke~\cite{AMS2020}.
Lopatto's full-support theorem~\cite{Lopatto2026} supplies that
positive-temperature support property. The finite-iteration
Gaussian and edge-energy analysis of El Alaoui, Montanari, and
Sellke~\cite{EAMS2023} then gives a near-optimal descendant formula.
This second route has its own quantitative and finite-degree
statements and does not use the zero-temperature companion.

\subsection{Model and quantifiers}

Let $(J_{ij})_{1\le i<j\le n}$ be independent standard real Gaussian random
variables. For $\sigma\in\{-1,1\}^n$, set
\begin{equation}\label{eq:sk-model}
 h_{n,J}(\sigma)=\frac1{\sqrt n}\sum_{i<j}J_{ij}\sigma_i\sigma_j,
 \qquad
 C_{n,J}=\frac1{\sqrt n}\sum_{i<j}J_{ij}Z_iZ_j,
 \qquad B_n=\sum_{i=1}^nX_i.
\end{equation}
Here $X,Y,Z$ are the Pauli matrices. The positive SK optimum is
\begin{equation}\label{eq:parisi-constant}
 \Pstar=\lim_{n\to\infty}\frac1n\E\max_{\sigma}h_{n,J}(\sigma).
\end{equation}
The ground-state Parisi formula and its attainment are due to
Auffinger and Chen~\cite{AC2017}. We use the covariance convention
$\xi(t)=t^2/2$; its exact relation to the unordered-pair model is
stated with the scalar input in Section~\ref{sec:classical-optimization}.
Appendix~\ref{sec:positive-temperature-alternative} also obtains this
limit from the positive-temperature formula and a log-sum comparison.

For $p\ge1$ and $\gamma,\beta\in\R^p$, define
\begin{equation}\label{eq:qaoa-state}
 \ket{\psi_{n,p,J}(\gamma,\beta)}
 =e^{-\ii\beta_pB_n}e^{-\ii\gamma_pC_{n,J}}
 \cdots e^{-\ii\beta_1B_n}e^{-\ii\gamma_1C_{n,J}}
 \ket{+}^{\otimes n}.
\end{equation}
All parameters are deterministic. Write
\begin{equation}\label{eq:qaoa-value}
 v_p(\gamma,\beta)=\lim_{n\to\infty}\frac1n
 \E\bra{\psi_{n,p,J}(\gamma,\beta)}C_{n,J}
 \ket{\psi_{n,p,J}(\gamma,\beta)},
 \qquad Q_p=\sup_{\gamma,\beta\in\R^p}v_p(\gamma,\beta).
\end{equation}
The fixed-parameter limit exists~\cite{FGGZ2022,BFMVZ2022}.

\begin{theorem}\label{thm:main}
For the zero-field SK model~\eqref{eq:sk-model},
\begin{equation}\label{eq:main-result}
 \lim_{p\to\infty}Q_p=\Pstar.
\end{equation}
Equivalently, for every $\varepsilon>0$ there exist a finite integer $p$ and
deterministic real vectors $\gamma,\beta\in\R^p$, independent of $n$ and $J$,
such that $v_p(\gamma,\beta)\ge\Pstar-\varepsilon$.
\end{theorem}

This proves the conjectured eventual Parisi optimality in the
fixed-parameter thermodynamic formulation of~\cite{BFMVZ2022}.
Measurement in the computational basis immediately gives $Q_p\le\Pstar$.
Appending a zero-angle layer gives $Q_{p+1}\ge Q_p$. Thus the theorem reduces
to constructing a finite word at every prescribed accuracy. The final word
may contain inverse gates, zero angles, and very large angles, all of which
are allowed in~\eqref{eq:qaoa-state}. A finite word in costs and mixers can
always be written in that alternating form by merging consecutive gates of
the same kind and inserting identity gates.

The theorem is stated in expected energy. The argument gives no
quantitative finite-size estimate, depth bound, or efficient procedure
for choosing the angles.

\subsection{How the construction works}

The proof separates a classical choice of target from its quantum
implementation. The intermediate setting is the large-degree
regular tree: each vertex carries an independent Gaussian label,
and a finite rooted calculation only reads labels within a fixed
radius. Normalized sums of centered child calculations become
jointly Gaussian coordinates. Section~\ref{sec:foundations} gives a
Hilbert-space realization of these coordinates and a continuous
quadratic expression for the limiting energy.

The scalar companion provides a bounded odd function of finitely
many independent Gaussians whose value approaches $P_*$. In
Section~\ref{sec:classical-optimization}, each new Gaussian is the
normalized aggregate of a suitable calculation one level below.
An orthogonality induction makes the joint Gaussian law exact.
The one-particle part of the resulting function gives its rooted
energy, so the scalar value formula passes to this realization
without a separate finite-degree edge expansion.

We next approximate all transmitted functions by bounded smooth
ones. A finite label assigns a \emph{type} to each vertex. Splitting
the label space into sufficiently small types lets every
aggregation omit the receiver's type, the types of all its path
ancestors, and a small reserved set of helper types, at arbitrarily
small energy cost. These exclusions prevent a compiled operation
from feeding back into its own parameters. They do not assert
independence of calculations using the same descendant subtree.
Proposition~\ref{prop:classical-near-optimum} gives the resulting
bounded finite formula, with its output written as $\sin(2\theta)$
for a bounded angle calculation $\theta$.

Theorem~\ref{thm:classical-synthesis} implements these angles using
cost and mixer gates together with a longitudinal field formed from
the Gaussian labels. To implement a neighbor sum, commutators combine
interactions along a receiver--sender--helper--sender walk. When the
walk returns to the same sender, the helper spin, held in its initial
state, supplies the desired term. We subdivide the sender type and
retain only walks whose two sender visits lie in the same part. This
keeps every returning walk, while the different-sender contribution
vanishes as the parts become small. The proof also controls the effects inside
the lower calculations, which are called repeatedly within a fixed
surrounding circuit. Helpers have arbitrarily small total probability;
discarding their output is included in the energy estimate.

Theorem~\ref{thm:seed-simulation} then removes the longitudinal
field from the already selected finite circuit. An imperfect cost echo
creates a small component with a large coherent displacement. Selective
spin rotations change that component while canceling their leading
action on the rest of the state; a long cost echo magnifies the change
into a finite longitudinal field. Separating its direction from the
circuit's existing Gaussian coordinates gives a seed that can be reused
at every requested seed gate. A separate one-particle calculation shows
that this Gaussian identification also preserves the rooted energy.
The elementary ensemble-control ideas are related to those of Li and
Khaneja~\cite{LiKhaneja2005}; Section~\ref{sec:seed-simulation}
proves the estimates for one fixed circuit, including its changing
insertion directions and the singular echo.

After all these choices there is one finite ordinary cost--mixer
word. Only at this point do we invoke the fixed-word tree-to-SK
identity of Basso et al.~\cite{BFMVZ2022}. The exact error allocation
and the order of choices are given in Section~\ref{sec:conclusion},
after the three constructions have been established.

\subsection{A regular MaxCut consequence}\label{sec:regular-maxcut}

For a finite simple graph $G=(V,E)$ with $E\ne\varnothing$, put
\[
 H_{\mathrm{MC}}(G)=\frac12\sum_{\{u,v\}\in E}(I-Z_uZ_v),
 \qquad B_G=\sum_{v\in V}X_v.
\]
For $a,b\in\R^p$, let $\langle\cdot\rangle^G_{p,a,b}$ denote expectation
with the gate order and signs of~\eqref{eq:qaoa-state}, using cost
$H_{\mathrm{MC}}(G)$, mixer $B_G$, angles $a,b$, and initial state
$\ket{+}^{\otimes |V|}$.  The expected fraction of edges cut by a
computational-basis measurement is
\[
 q_{p,G}(a,b)=\frac{\langle H_{\mathrm{MC}}(G)\rangle^G_{p,a,b}}{|E|}.
\]
Write $\operatorname{OPT}(G)=\operatorname{MaxCut}(G)/|E|$ for the maximum
cut fraction.

\begin{corollary}[Regular MaxCut in successive limits]
\label{cor:regular-maxcut}
For every $\eta>0$ there are an integer $p\ge1$, deterministic normalized
angles $\gamma,\beta\in\R^p$, and an integer $D_0\ge2$ such that, for every
integer $D\ge D_0$ and every finite simple unweighted $(D+1)$-regular graph
$G$ of girth greater than $2p+1$,
\begin{equation}\label{eq:high-girth-maxcut}
 q_{p,G}(2\gamma/\sqrt D,\beta)
 \ge \frac12+\frac{\Pstar-\eta}{\sqrt D}.
\end{equation}
The word $p,\gamma,\beta$ is independent of $D$, the graph size, and $G$;
the displayed usual MaxCut cost angles, unlike the normalized angles
$\gamma$, are scaled with the degree.

Let $G_{n,d}$ be uniform among simple labelled $d$-regular graphs on $n$
vertices, along admissible $n$ with $nd$ even.  For the same word, put
\[
 \Delta_D(G)=\operatorname{OPT}(G)
             -q_{p,G}(2\gamma/\sqrt D,\beta).
\]
The word may be chosen so that
\begin{equation}\label{eq:random-regular-maxcut}
 \lim_{D\to\infty}
 \limsup_{\substack{n\to\infty\\ n(D+1)\ \mathrm{even}}}
 \Pr_G\!\left[\sqrt D\,\Delta_D(G_{n,D+1})>\eta\right]=0.
\end{equation}
Here the inner limit keeps $D,p,\gamma,\beta$ fixed, and $\Pr_G$ is only
over the graph: $q_{p,G}$ already averages quantum measurement.
\end{corollary}

The complete proof is given in Section~\ref{sec:regular-maxcut-proof},
after the fixed-word tree interface and the SK construction.

The optimality comparison concerns random regular graphs in these
successive limits, not all high-girth graphs or fixed-degree optima.
No joint size--degree limit is asserted. Here $p$ counts QAOA rounds,
not a degree-uniform two-qubit gate depth; the cut fraction above
averages quantum measurement.

\subsection{Reading paths and conventions}

Section~\ref{sec:foundations} develops the fixed-word rooted
calculus. Section~\ref{sec:classical-optimization} imports the scalar
theorems and constructs the bounded excluded formula.
Sections~\ref{sec:classical-synthesis} and~\ref{sec:seed-simulation}
implement it and remove the seeds. Section~\ref{sec:conclusion}
assembles the ordinary word and proves Theorem~\ref{thm:main}.
Section~\ref{sec:regular-maxcut-proof} then proves the regular-MaxCut
consequence from this theorem and the fixed-word tree interface.

For the positive-temperature reading path,
Appendix~\ref{sec:positive-temperature-alternative} replaces the
scalar input and its rooted realization. It derives the same
classical conclusion independently, using the shared bounded
approximation and exclusion lemmas of
Section~\ref{sec:classical-optimization}. The quantum constructions
and final assembly are the same for both paths.

Physical system size is $n$, tree degree is $D+1$, and ordinary
circuit depth is $p$. The scalar mesh count in
Section~\ref{sec:classical-optimization} is $N$; top-generation
occupation operators appearing in the operator estimates are
identified by their operator context. Mixer angles are denoted by
$\beta$, inverse temperature by $\beta_0$, and the scalar Parisi
coefficient $\gamma(t)$ is distinct from the vector of cost angles.
Inner products are antilinear in the first argument.

All constants may depend on a fixed finite formula or word.
Uniformity in any approximation parameter is stated when required.
Degree limits are taken with all words and their angles fixed;
further approximation limits concern the resulting tree values.

\section{The fixed-word tree calculus}
\label{sec:foundations}

This section constructs the common space in which we compare ordinary
QAOA words and classical formulas.  It has three tasks: identify the
fixed-word SK value with a tree limit, express that limit as a bounded
quadratic form, and justify the controlled operations used in synthesis.
The same space will carry the finite Gaussian construction in
Section~\ref{sec:classical-optimization}.

We first isolate the transfer from an ordinary word to the tree.
A word always means a finite word with fixed real coefficients; none
of the statements below makes a uniform assertion in its length or
its coefficients. The initial vector $\vac$ below is the plus product
state on the finite light cone used to define the edge expectation.

\begin{proposition}[Fixed-word transfer]\label{prop:tree-transfer}
Let $W$ be an ordinary QAOA word, and apply the same word to the
$(D+1)$-regular tree with cost
$C_D=D^{-1/2}\sum_{\{i,j\}\in E}Z_iZ_j$, with each unordered edge
counted once.  Its expectation at an edge is
defined by restricting to the finite light cone of that edge.  Then
\begin{equation}\label{eq:tree-transfer}
 v_W=\lim_{D\to\infty}\frac{\sqrt D}{2}
       \bra{W_D\vac}Z_iZ_j\ket{W_D\vac}.
\end{equation}
Here $v_W$ is the disorder-averaged thermodynamic SK energy per vertex,
with the normalization in the introduction.  Both limits exist for
every fixed word and every fixed choice of its angles.
\end{proposition}

\begin{proof}
This is Theorem~1 of \cite{BFMVZ2022}, together with its
normalizations in Equations~(2.3), (2.6), and (6.1)--(6.2).
Its Equation~(3.11) gives the fixed-depth tree value, and
Equation~(6.3) identifies that value with the SK limit.
That paper uses the negative of our tree cost.
Conjugation by $X$ on one bipartition of the tree changes the
sign of every $Z_iZ_j$, preserves every mixer, and fixes the initial
vector.  It therefore changes the sign of the edge correlation as
well, giving precisely \eqref{eq:tree-transfer}.  The SK convention
already agrees with the positive Gaussian Hamiltonian in its
Equation~(6.1).  Consecutive gates of the same kind and inverses are
included by inserting zero-angle gates.
\end{proof}

We first construct the rooted space and identify its energy functional.
This also gives the Gaussian coordinates used by the classical
construction. We then prove the weighted convergence estimates that
realize ordinary words in this space and justify controlled gates inside
a fixed surrounding circuit.

\subsection{Pointed spaces and occupation embeddings}
\label{sec:pointed-spaces}

The number of children will tend to infinity while the number of
circuit layers remains fixed.  We therefore organize a symmetric
collection of child states by how many differ from the initial state.
The limiting space is a symmetric Fock space: its $m$-particle sector
records $m$ such child excitations, not $m$ physical spins.
The bosonic representation is related to the fixed-depth spin--boson
description of Boulebnane et al.~\cite{BKL2026}; the occupation,
insertion, and surrounding-circuit estimates used below are proved here.

Let $(S,\mu,\iota)$ be a probability space with a measure-preserving
involution.  The cases needed below are a one-point space and a
standard Gaussian, with $\iota(s)=-s$.  Extra independent marks may
be included in $S$.  Set
\[
 \mathsf k=\C^2\otimes L^2(S,\mu),\qquad
 \omega=\ket{+}\otimes\1.
\]
For any pointed space below, write $\langle T\rangle_0=
\bra{\vac}T\ket{\vac}$ for the vacuum mean and
$Q=1-|\vac\rangle\langle\vac|$ for its centered projection.
All seed variables act by multiplication.  In particular, a circuit
never changes their values.  The local symmetry is
$X\otimes(f\mapsto f\circ\iota)$.

For a pointed Hilbert space $(H,\vac)$ write
$H^\circ=H\ominus\C\vac$ and
$\Gamma_s(H^\circ)=\bigoplus_{m\geq0}(H^\circ)^{\otimes_s m}$.
At finite height define
\begin{equation}\label{eq:recursive-spaces}
 \cH_0=\mathsf k,\qquad
 \cH_{h+1}=\mathsf k\otimes\Gamma_s(\cH_h^\circ),\qquad
 \vac_{h+1}=\omega\otimes\Omega.
\end{equation}
The embedding from height $h$ to height $h+1$ appends empty
descendants: it is $\xi\mapsto\xi\otimes\Omega$ at height zero,
and thereafter is obtained by applying the symmetric Fock functor
to the preceding embedding.  These embeddings preserve the vacuum.
We identify the spaces with their images and use the completion of
their union when convenient.  Every construction in the proof takes
place at some finite height before this completion is used.

Write $R_\iota f=f\circ\iota$.  The joint spin/seed parity on these
spaces is the vacuum-preserving involution
\[
 \mathscr P_0=X\otimes R_\iota,\qquad
 \mathscr P_{h+1}=(X\otimes R_\iota)\otimes
       \Gamma_s(\mathscr P_h|_{\cH_h^\circ}).
\]
These involutions respect the height embeddings and fix the vacuum,
so every odd observable has zero vacuum mean.  We suppress the height
in $\mathscr P$ when the space is understood.

The analogous invariant space on a rooted $D$-ary tree is
\[
 \cH_{D,0}=\mathsf k,\qquad
 \cH_{D,h+1}=\mathsf k\otimes\operatorname{Sym}^D(\cH_{D,h}).
\]
The symmetry here permutes children together with their entire
subtrees and seed coordinates.  It does not discard the seeds.  For
any pointed space $H$, the orthogonal decomposition by the number
of nonvacuum tensor factors gives an isometry
\begin{equation}\label{eq:occupation-embedding}
 \mathcal I_D:\operatorname{Sym}^D(H)\longrightarrow
       \bigoplus_{m=0}^D(H^\circ)^{\otimes_s m}.
\end{equation}
On a symmetric $m$-particle tensor, the inverse distributes its $m$
nonvacuum factors over the $D$ slots with coefficient
$\binom Dm^{-1/2}$.  Applying this isometry recursively gives
isometries $\mathcal I_{D,h}:\cH_{D,h}\to\cH_h$.  Their range
projections tend strongly to the identity at each fixed height.
Indeed, a finite-particle vector involving finitely many child
vectors is approximated inductively by vectors in these ranges.
The finite-tree parity is the tensor product of the local
involutions over its vertices; the occupation embeddings
intertwine it with $\mathscr P_h$.

Creation is linear in its argument, annihilation is antilinear, and
\[
 q(d)=a^*(d)+a(d),\qquad
 [q(d),q(e)]=2\ii\im\braket{d}{e}.
\]
Conjugation by the recursive parity sends the child field $q(d)$
to $q(\mathscr P_h d)$; an odd child direction therefore gives
an odd field at the next height.
The fields are first defined on finite-particle vectors and then
closed.  Each $q(d)$ is essentially self-adjoint there, as follows
either from the one-mode Schr\"odinger representation or from the
convergence of its exponential series on finite-particle vectors.
We use $\mathcal W(d)=e^{\ii q(d)}$.  In particular,
\begin{equation}\label{eq:weyl-shift}
 \mathcal W(d)^*q(v)\mathcal W(d)
       =q(v)-2\im\braket{v}{d}.
\end{equation}

\subsection{Root insertions and edge energy}
\label{sec:root-energy}

At finite degree, a rooted vector records the response of a root
observable to a circuit. Let $U$ be a circuit on a finite subtree and
$U_j$ its restriction to child subtree $j$. Factoring the child
operations defines the root interaction-picture unitary $u$ by
\[
 U=\left(\bigotimes_{j\text{ child}}U_j\right)u.
\]
For a root Pauli observable $P$, its insertion vector is
\[
 f_P=U^*PU\vac=u^*Pu\vac,
\]
because the child operations commute with $P$. For a word evaluated
on a neighborhood containing both endpoint light cones, the pairing
of its two $Z$ insertions gives the expectation of $Z_iZ_j$.
We will prove convergence of ordinary insertion vectors
after establishing the weighted estimates below. First we identify
the energy of any convergent family of rooted vectors with bounded
expected occupation; this geometric step needs no circuit estimates.

The insertion vector lives at one root, whereas an edge joins two
rooted neighborhoods.  The following operator moves one centered
vector down to a single child while placing the new root in its
initial state.  It will turn the edge pairing into a quadratic form
on one common space.

Define the centered root shift by
\begin{equation}\label{eq:root-shift}
 \cL:\cH^\circ\longrightarrow\cH^\circ,
 \qquad \cL v=\omega\otimes a^*(v)\Omega.
\end{equation}
At height $h$ its range lies at height $h+1$.  The definitions
are compatible with the height embeddings, so $\cL$ is an
isometry on the centered direct limit.  Extend it by
$\cL\vac=0$ when writing $\cL^*$; this extension is a partial
isometry on the whole pointed space.  In particular $\cL^*$
extracts the root-vacuum one-particle component, with a constant
root seed.

\begin{proposition}[Energy as a root shift]\label{prop:energy}
Let $f_D$ be centered, root-invariant, finite-radius vectors on
the $(D+1)$-regular tree.  Root invariance means that the vector
itself is fixed by every rooted permutation of complete child
subtrees, including their seed coordinates, at every vertex of
its finite neighborhood.  Suppose their occupation embeddings
converge strongly to $v$ at a fixed finite height and their
expected top-generation occupations are bounded uniformly in
$D$.  Place the same rooted vector at the endpoints of an edge,
calling the resulting vectors $f_{D,i}$ and $f_{D,j}$.  Then
\begin{equation}\label{eq:energy-shift}
 \lim_{D\to\infty}\frac{\sqrt D}{2}
       \braket{f_{D,i}}{f_{D,j}}
       =\cE(v):=\re\braket{v}{\cL v}.
\end{equation}
For any centered $v,w$,
\begin{equation}\label{eq:energy-continuity}
 |\cE(v)-\cE(w)|
       \leq(\|v\|+\|w\|)\|v-w\|.
\end{equation}
\end{proposition}

\begin{proof}
We give the exact finite-degree decomposition.  Put $M=D+1$,
let $E$ be the pointed child space, and write $N_M$ for the
number of nonvacuum children in
$\mathsf k\otimes\operatorname{Sym}^M(E)$.  Let $R_M$ contract
a distinguished child against its vacuum, and let
\[
 J_M\xi=\omega\otimes\frac1{\sqrt M}
 \sum_{j=1}^M\vac^{\otimes(j-1)}\otimes\xi
                         \otimes\vac^{\otimes(M-j)}.
\]
The latter is an isometry from $E^\circ$.  Under occupation
embeddings, $J_M^*=\cL^*$ and
\begin{equation}\label{eq:vacuum-deletion}
 R_M=\sqrt{1-N/M},\qquad
 \|R_Mf-f\|^2\leq M^{-1}\bra fN_M\ket f.
\end{equation}
Indeed, in the $m$-particle sector the fraction of placements
surviving deletion is
$\binom{M-1}m/\binom Mm=1-m/M$.

Delete the distinguished edge and let $P_A,P_B$ denote the
vacuum projections on its two half-trees.  Decompose $f_{D,i}$
with $P_B+(1-P_B)$ and $f_{D,j}$ with $P_A+(1-P_A)$.
The term with both vacuum projections is zero: the surviving
half-tree vectors are centered.  If just one side is excited,
the other side being vacuum forces the root and all the other
branches of the first vector to be vacuum.  With
$s_D=J_M^*f_D$ and $r_D=R_Mf_D$, the two such terms are
$\braket{s_D}{r_D}/\sqrt M$ and its conjugate.  Heights in
this pairing are matched by adding empty descendants.  Each
remaining excited-side vector has squared norm
$\bra{f_D}N_M\ket{f_D}/M$, by symmetry among branches.  Thus
\begin{equation}\label{eq:exact-edge-decomposition}
 \braket{f_{D,i}}{f_{D,j}}
 =\frac2{\sqrt M}\re\braket{s_D}{r_D}+\varepsilon_D,
 \qquad
 |\varepsilon_D|\leq\frac{\bra{f_D}N_M\ket{f_D}}M.
\end{equation}
After multiplication by $\sqrt D/2$ the error tends to zero.
Equation~\eqref{eq:vacuum-deletion} and strong convergence show
that the first term tends to
$\re\braket{\cL^*v}{v}=\re\braket{v}{\cL v}$.
This proves \eqref{eq:energy-shift}.

Finally, \eqref{eq:energy-continuity} follows by adding and subtracting
$\braket{w}{\cL v}$ and using $\|\cL\|=1$.
\end{proof}

The occupation hypothesis in Proposition~\ref{prop:energy}
is a sufficient hypothesis for the displayed proof; strong
convergence alone does not bound the remainder in
\eqref{eq:exact-edge-decomposition}.  Lemma~\ref{lem:ordinary-updates} and
Theorem~\ref{thm:gate-calculus} will supply that hypothesis for
the quantum gates used here.

\begin{remark}[Classical rotations]\label{rem:classical-rotation-energy}
The classical sector consists of vectors whose spin is $\ket{+}$
at every vertex and whose coefficients depend only on seeds.
If a $Y$ rotation whose angle depends only on seeds produces an insertion
$v=\ket{+}m+\ket{-}n$ with seed-only coefficients in all
descendants, then $\cE(v)=\cE(\ket{+}m)$.  Indeed every
$\cL v$ has root spin $\ket{+}$, while
$\cL(\ket{-}n)$ has a child spin excitation and is orthogonal
to the classical sector.  Seed-only here permits dependence
on descendant seeds; it does not mean dependence only on the
root seed.
\end{remark}

\subsection{Classical Gaussian positions and descendant formulas}
\label{subsec:classical-gaussian-sector}

The root-shift functional also applies to classical vectors, before
they are compiled into quantum gates.  We now give a common Gaussian
realization for both classical constructions in this paper.  This is
a definition on the limiting pointed space; a finite-degree energy
limit still requires the hypotheses of Proposition~\ref{prop:energy}.

Keep every quantum spin in $\ket{+}$ and use real coefficients in the
seed and symmetric Fock factors.  This gives the real classical
subspace.  For a real centered classical child vector $f$, let
$\mathsf G(f)=q(f)$ be the position operator on the root's child
Fock factor.  For any finite list $f_1,\ldots,f_k$, these positions
strongly commute and have joint centered Gaussian vacuum law with
covariance $\braket{f_i}{f_j}$.  Indeed, choose a real orthonormal
basis of their finite-dimensional span.  The position operators of
these orthogonal modes act on separate Fock factors, and
\[
 \left\langle
       e^{\ii\sum_{j=1}^k t_j\mathsf G(f_j)}
 \right\rangle_0
 =\exp\left(-\frac12\left\|\sum_{j=1}^kt_jf_j\right\|^2\right).
\]
In particular, orthonormal directions give independent standard
normal positions: the variance convention is $\|f\|^2$, with no
additional factor of two.  The positions act trivially on the local
seed factor.  Their joint vacuum law is therefore independent of the
receiving root's seed, even when a child direction itself depends on
the seed at its own root.

We identify a scalar function of these commuting variables with the
vector obtained by applying that function to the vacuum.  Under this
identification,
\begin{equation}
 \norm{\bigl(\mathsf G(f)-\mathsf G(g)\bigr)\vac}
      =\norm{f-g}.
 \label{eq:classical-gaussian-isometry}
\end{equation}
All fields are realized on the same pointed space, so this equality
also specifies their coupling.  The empty-descendant embeddings in
\eqref{eq:recursive-spaces} preserve inner products, commute with
passing to the corresponding Fock fields, and keep the local seed
factor fixed.  Thus directions constructed at different finite
heights can first be embedded in one common height without changing
their joint Gaussian law.

\begin{definition}[Finite descendant formula]
\label{def:finite-descendant-formula}
A finite descendant formula is a finite expression built recursively
from a local seed and finitely many Gaussian positions of centered
formulas one level below, using pointwise scalar functions.  The
functions are required to be measurable and to produce square-integrable
vectors whenever the formula is used as a vector.  Additional empty
levels identify formulas of different radii.  Its finite-degree version
replaces every Gaussian position by the normalized sum of the
corresponding independent child variables.
\end{definition}

The definition by itself is a limiting-space construction, not a
finite-degree convergence theorem.  Later approximation arguments
impose boundedness or polynomial-growth conditions on the pointwise
maps as needed.  Because $\cL$ places a constant seed at the new root,
$\cL^*$ extracts precisely the first Gaussian chaos with constant
receiving-root seed.  If $e_1,\ldots,e_k$ are orthonormal real centered
child vectors and $F$ is a square-integrable function of
$\mathsf G(e_1),\ldots,\mathsf G(e_k)$ alone, this gives
\begin{equation}\label{eq:classical-first-chaos}
 \cL^*F
   =\sum_{j=1}^k
       \E\bigl[F\,\mathsf G(e_j)\bigr]e_j.
\end{equation}
Here $F$ may have a nonzero constant component, which $\cL^*$
annihilates.  The identity follows by the orthogonal projection onto
the first chaos in these Gaussian coordinates; the remaining
orthogonal child modes and all higher chaoses have zero projection.
This is the normalization used in the primary classical construction.

The rooted space, Gaussian positions, and energy functional now give
the setting for the classical construction in
Section~\ref{sec:classical-optimization}. The rest of this section
establishes convergence and continuity of the quantum operations that
will implement it.

\subsection{Weights and the topology of convergence}
\label{sec:smooth-topology}

The fields just defined are unbounded.  To pass their products and
parameter derivatives through limits, we need a common domain and
bounds that control both seed moments and excitations at every level.
These recursive weights serve that purpose; they are distinct from
the top-generation number weight used in the seed-removal argument.
If all local-seed coefficients under consideration are uniformly
bounded, take $W_0=1$ on $\mathsf k$.  For coefficients of polynomial
growth in a Gaussian seed, start with multiplication by $1+|s|$
and make it block diagonal at $\omega$:
\[
 W_0=P_0+Q_0(1+|s|)Q_0,
 \qquad P_0=\ket{\omega}\bra{\omega},\quad Q_0=1-P_0.
\]
Compression is understood on its natural closed domain.  This is
a positive self-adjoint operator bounded below by one.  It differs
from multiplication by $1+|s|$ by a bounded finite-rank operator
whose range consists of vectors with all seed moments.  Consequently
their positive-integer power domains coincide, with equivalent
graph norms.  For completeness, this last assertion follows by
induction from $A^m-B^m=\sum_{j=0}^{m-1}A^j(A-B)B^{m-1-j}$ on the
common core; every occurrence of $A-B$ has smooth finite-dimensional
range, and the resulting estimates close in both directions.

Since $W_0\omega=\omega$, we can recursively set
\begin{equation}\label{eq:recursive-weights}
 \Lambda_h=W_h|_{\cH_h^\circ},\qquad
 W_{h+1}=W_0\otimes1+1\otimes\dd\Gamma(\Lambda_h).
\end{equation}
Thus $W_h\vac_h=\vac_h$ and $W_h\geq1$.  On a finite tree use
$W_{D,h+1}=W_0+\sum_{j=1}^D Q_jW_{D,h}^{(j)}Q_j$.
The occupation embeddings intertwine these weights exactly.  This
also shows that their range projections commute with every power
of $W_h$.

Write
\[
 \mathscr S_h=\bigcap_{m\geq0}\operatorname{Dom}(W_h^m),
 \qquad \|\xi\|_m=\|W_h^m\xi\|.
\]
Finite-particle tensors of smooth child vectors and local vectors
with all seed moments are dense in every one of these graph norms.
One proof first cuts off the spectral variables of $W_0$ and each
$\Lambda_h$, then cuts off particle number, and finally approximates
in the resulting bounded spectral subspaces by algebraic tensors.
The finite-degree projections give simultaneous approximants in all
fixed graph norms.

When comparing unitary maps on the varying embedded spaces,
extend each embedded unitary by the identity on the orthogonal
complement of its range projection.  These are genuine
unitaries on the common Hilbert space.  The projections commute
with the weights and converge strongly to one in every fixed
graph norm on smooth vectors, so this convention preserves
all the bounds and limits.  Parameter derivatives of the
extensions vanish on that complement.

We will repeatedly use the following precise convention.  An
operator family $F_D(x)$ has \emph{uniform smooth bounds} if, for
every nonnegative integer $m$ and multi-index $\alpha$, there are
$r,k,C$, independent of $D$ and $x$, such that
\begin{equation}\label{eq:smooth-family-bound}
 \|W_D^m\partial_x^\alpha F_D(x)\xi\|
 \leq C(1+|x|)^k\|W_D^r\xi\|,
 \qquad \xi\in\mathscr S_D.
\end{equation}
The same condition is required of the adjoint family.  Constants
may depend on the fixed expression, its finite list of channels,
the height, $m$, and $\alpha$.  Dependence on a further limiting
parameter is explicitly stated when relevant.  Smooth convergence
of operator families means strong convergence on smooth inputs
of both the families and their adjoints, including all specified
parameter jets, with these uniform bounds.  The limits of the
adjoint jets must be the adjoints of the corresponding limiting
jets on the common domain; bounded adjoint jets alone do not
imply their convergence.  For vector families only the vector
convergence is required.  This definition also permits triangular
arrays of one-copy spaces under the embeddings above.

Two elementary consequences will be useful.  Products of a fixed
number of such families have the same bounds, with a possibly
larger loss of input powers.  Also, strong convergence together
with bounded higher graph norms gives convergence in every fixed
lower graph norm.  The latter follows from the spectral inequality
\begin{equation}\label{eq:weight-interpolation}
 \|W^m\xi\|\leq R^m\|\xi\|
       +R^{m-M}\|W^M\xi\|\qquad(M>m, R\geq1),
\end{equation}
by taking the strong limit first and $R\to\infty$ second.
Consequently it is enough to prove convergence on the dense domain
just described, and then use the uniform estimates.  This is the
meaning of convergence between different finite-degree spaces.

\begin{lemma}[Weyl estimates]\label{lem:weyl-bounds}
Let $N$ be the particle number on $\Gamma_s(H^\circ)$.  For every
integer $m\geq0$ there is $C_m$ such that
\begin{align}
 \|(1+N)^m\mathcal W(d)\xi\|
 &\leq C_m(1+\|d\|)^{2m}\|(1+N)^m\xi\|,
       \label{eq:weyl-number-bound}\\
 \|(\mathcal W(d)-1)\xi\|
 &\leq2\|d\|\|(N+1)^{1/2}\xi\|.
       \label{eq:weyl-small-bound}
\end{align}
For the weight $W=W_0+\dd\Gamma(\Lambda)$, the same conclusion
holds with a polynomial in finitely many $\|\Lambda^j d\|$ in
place of $(1+\|d\|)^{2m}$ and with a finite loss of input powers.
All statements hold for spin-controlled displacements as well.
\end{lemma}

\begin{proof}
On finite-particle vectors,
$\|a(d)\xi\|\leq\|d\|\|N^{1/2}\xi\|$ and
$\|a^*(d)\xi\|\leq\|d\|\|(N+1)^{1/2}\xi\|$.
The fundamental theorem of calculus applied to
$e^{\ii t q(d)}\xi$ proves \eqref{eq:weyl-small-bound}.
Conjugation gives
\[
 \mathcal W(d)^*N\mathcal W(d)=N+q(\ii d)+\|d\|^2.
\]
Commuting $N$ through a field changes that field to another field
with the same norm of its argument.  Expanding the $m$th power
and applying the creation and annihilation bounds therefore proves
\eqref{eq:weyl-number-bound}.  For the weighted version use
\[
 \mathcal W(d)^*\dd\Gamma(\Lambda)\mathcal W(d)
 =\dd\Gamma(\Lambda)+q(\ii\Lambda d)
       +\braket{d}{\Lambda d},
\]
and $[\dd\Gamma(\Lambda),a^*(d)]=a^*(\Lambda d)$.
Every expansion uses only finitely many weighted norms of $d$.
Approximation in graph norms proves the identities and estimates
on the stated domains.  A spin-controlled displacement is the
sum of two ordinary displacements on orthogonal spin subspaces.
\end{proof}

\subsection{Normalized sums and a unitary central limit}
\label{sec:unitary-clt}

A normalized sum creates or removes one child excitation in the
limit.  A product of small unitaries has the corresponding field
exponential as its limit, but may also retain a scalar second-order
term.  We prove both statements with the weighted bounds needed for
later substitution into a fixed circuit.

\begin{lemma}[Sums]\label{lem:centered-sums}
Let $T_D$ and $T_D^*$ have uniform smooth bounds on one-copy
spaces, and suppose their smooth limits exist.  If
$\braket{\vac}{T_D\vac}=0$, then under the occupation embeddings
\begin{equation}\label{eq:centered-sum-limit}
 D^{-1/2}\sum_{j=1}^D T_D^{(j)}
 \longrightarrow a^*(T\vac)+a(T^*\vac)
\end{equation}
smoothly.  Without centering,
$D^{-1}\sum_jT_D^{(j)}\to\braket{\vac}{T\vac}\,1$.
Both families have uniform smooth bounds.
\end{lemma}

\begin{proof}
Put $g=QT_D\vac$, $h=QT_D^*\vac$, and $A=QT_DQ$.
For centered $T_D$ its normalized sum has the exact occupation
representation
\begin{equation}\label{eq:sum-decomposition}
 a^*(g)\sqrt{1-N/D}
 +\sqrt{1-N/D}\,a(h)+D^{-1/2}\dd\Gamma(A),
\end{equation}
where factors beyond the truncation are zero.  On the $m$-particle
sector the two square roots express, respectively, the numbers
$D-m$ of vacant and $D-m+1$ of potentially filled slots.
The formula first holds on algebraic tensors and then on smooth
vectors.  Creation and annihilation have uniform weighted bounds
by the proof of Lemma~\ref{lem:weyl-bounds}.  For the remaining
term, a one-copy bound between powers of $\Lambda$ implies a bound
for its sum between powers of $1+\sum\Lambda_j$: apply the
one-copy bound to each occupied slot, use
$\Lambda_j\leq\sum\Lambda_j$, and pay at most one additional
power for the number of occupied slots.  This argument works at
any desired output power.  It proves the required bounds for
\eqref{eq:sum-decomposition}.  On a fixed-particle smooth tensor
the last term tends to zero and the square roots tend to one;
density and \eqref{eq:weight-interpolation} finish the proof.
For an average the same decomposition has off-diagonal factor
$D^{-1/2}$ and complementary-block factor $D^{-1}$, while its
vacuum block tends to the stated scalar.
\end{proof}

The sum lemma determines first-order fields.  Controlled gates also
require products of $D$ factors whose arguments are of size
$D^{-1/2}$.  The next lemma keeps their second-order vacuum contribution;
it cannot in general be recovered from the first-order field alone.

\begin{lemma}[Products of small unitaries]\label{lem:unitary-clt}
Let $K_D(x)$, $x\in\R^b$, be unitary, satisfy $K_D(0)=1$, and
have uniform smooth bounds for all jets and adjoint jets.  Suppose
these jets converge smoothly and
$\braket{\vac}{\partial_\alpha K_D(0)\vac}=0$ for every first
derivative.  Write
\[
 K_\alpha=\partial_\alpha K(0)=\ii T_\alpha,
 \qquad K_{\alpha\lambda}=\partial_\alpha\partial_\lambda K(0).
\]
The $T_\alpha$ are symmetric on the smooth domain.  For every
fixed $a\in\R^b$,
\begin{equation}\label{eq:unitary-product-limit}
 K_D(a/\sqrt D)^{\otimes D}\longrightarrow
 \exp\left\{\ii\sum_\alpha a_\alpha q(T_\alpha\vac)
 +\frac12\sum_{\alpha,\lambda}a_\alpha a_\lambda
 \left\langle K_{\alpha\lambda}
           +\tfrac12\{T_\alpha,T_\lambda\}\right\rangle_0
 \right\}.
\end{equation}
The scalar in this exponent is imaginary.  The convergence holds
smoothly with all $a$-derivatives, locally uniformly in $a$, and
the product has polynomial smooth bounds in $a$, uniform in $D$.
The assertion also holds with an additional finite list of
unscaled parameters $\theta$: use
$K_D(a/\sqrt D;\theta)^{\otimes D}$, assume
$K_D(0;\theta)=1$ and first-derivative centering for every
$\theta$, and assume the hypotheses jointly for all parameter
jets, uniformly on compact $\theta$-sets.  The conclusion then
includes all $\theta$-jets and the same compact uniformity.
\end{lemma}

\begin{proof}
We first identify the candidate limit from exponential vectors;
then we prove the bounds needed for smooth convergence.  If $x,z\perp\vac$
are smooth, the vectors
$(\vac+x/\sqrt D)^{\otimes D}$ converge to
$\varepsilon(x)=\bigoplus_{m\geq0}x^{\otimes m}/\sqrt{m!}$,
and similarly for $z$.  The coefficient of $D^{-1}$ in the
one-copy matrix element is
\begin{align*}
 c={}&\braket{z}{x}
 +\sum_\alpha a_\alpha
       \bigl(\braket{z}{K_\alpha\vac}
                      +\braket{\vac}{K_\alpha x}\bigr)
 +\frac12\sum_{\alpha,\lambda}a_\alpha a_\lambda
                        \langle K_{\alpha\lambda}\rangle_0.
\end{align*}
Taylor's formula makes the remaining error $o(D^{-1})$, also
for converging triangular arrays of smooth $x,z$.  The full
pairing consequently tends to $e^c$.  For
$t=\sum_\alpha a_\alpha T_\alpha\vac$, the exponential-vector
pairing of $e^{\ii q(t)}$ has an additional term $-\|t\|^2/2$.
Since
$\|t\|^2=\sum_{\alpha,\lambda}a_\alpha a_\lambda
\langle\{T_\alpha,T_\lambda\}/2\rangle_0$,
the correction in \eqref{eq:unitary-product-limit} is exactly
the required one.  Twice differentiating $K(x)^*K(x)=1$ at zero
shows that its real part is zero. The displayed limit is thus unitary.

We next prove uniform bounds for the product
$U_D(a)=K_D(a/\sqrt D)^{\otimes D}$. For one child let $A=Q\Lambda Q$ and
$B_D(a)=K_D(a/\sqrt D)^*A K_D(a/\sqrt D)$.  Taylor's formula
at zero, with the assumed bounds, gives
\begin{align*}
 |\langle B_D(a)\rangle_0|&\leq D^{-1}P(a),\\
 \|\Lambda^m QB_D(a)\vac\|
  +\|\Lambda^m QB_D(a)^*\vac\|&\leq D^{-1/2}P_m(a),
\end{align*}
and the complementary block has polynomial smooth bounds.  Here
and below $P_m$ denotes some fixed nonnegative polynomial in
$1+|a|$, whose degree is allowed to depend on $m$.
The same vacuum-row decomposition as in
\eqref{eq:sum-decomposition}, now applied to the unnormalized
sum of the $B_D(a)$, shows that
\[
 U_D(a)^*WU_D(a)=W_0+\sum_jB_D(a)^{(j)}
\]
has polynomial bounds from any sufficiently high input weight
to any prescribed output weight, uniformly in $D$.
Composing this estimate a fixed number of times bounds its powers.
Since $W^mU_D=U_D(U_D^*WU_D)^m$, this proves the assertion for
$U_D$ itself; the adjoint is treated identically.  All identities
may first be made on smooth tensors, which the product preserves,
and then closed using these very estimates.

For derivatives,
\[
 U_D(a)^*\partial_\alpha U_D(a)
 =D^{-1/2}\sum_j
 \bigl(K_D^*\partial_\alpha K_D\bigr)(a/\sqrt D)^{(j)}.
\]
The one-copy vacuum mean is $O(D^{-1/2}P(a))$, since it vanishes
at zero.  Lemma~\ref{lem:centered-sums}, with this scalar part
retained, gives a uniform bound.  Further differentiation produces
products of finitely many such sums.  A sum carrying $k\geq2$
derivative-scaling factors is bounded by the average estimate
times $D^{1-k/2}\leq1$.  Induction proves all jet bounds.

Weak convergence on the total set of exponential
vectors, together with equality of the limiting norms, proves
strong convergence.  Higher-weight convergence follows from
\eqref{eq:weight-interpolation}; derivative convergence follows
from finite differences and the uniform bounds on one further
derivative.  This also gives local uniformity in $a$.
Since the strong limit is unitary, the adjoints converge
strongly as well: for unitaries $U_D,U$, one has
$\|(U_D^*-U^*)\xi\|=\|\xi-U_DU^*\xi\|$.
The adjoint jet bounds already proved, followed by the same
finite-difference argument, give convergence of adjoint jets.

To verify the unscaled-parameter assertion, the only additional
estimate concerns a logarithmic derivative
$A_\nu(x;\theta)=K(x;\theta)^*\partial_{\theta_\nu}K(x;\theta)$.
It satisfies $A_\nu(0;\theta)=0$.  Moreover
\[
 \left\langle\partial_{x_\alpha}A_\nu(0;\theta)\right\rangle_0
 =\partial_{\theta_\nu}
       \left\langle\partial_{x_\alpha}K(0;\theta)\right\rangle_0
 =0.
\]
Its vacuum entry at $a/\sqrt D$ is therefore
$O(D^{-1}P(a))$, its vacuum row and column are
$O(D^{-1/2}P(a))$, and its complementary block is
$O(D^{-1/2}P(a))$ with smooth bounds.  The unnormalized
sum of its $D$ copies consequently has uniform smooth
bounds by the decomposition used above.  Every further
$\theta$-derivative has the same two vanishings, and mixed
$a$-derivatives are treated by the normalized-sum or average
estimates.  This proves all joint jet bounds.  Pointwise
convergence for fixed $\theta$, finite differences, and the
adjoint argument just given prove the parametric conclusion.
\end{proof}

\begin{remark}\label{rem:unitary-clt-domains}
The anticommutator in \eqref{eq:unitary-product-limit} is a
quadratic form on smooth vectors; its vacuum expectation is
unambiguous because all jets preserve that domain.  Essential
self-adjointness of the one-copy derivative $T_\alpha$ is not
needed: only the vector $T_\alpha\vac$ enters the self-adjoint
Fock field.  This avoids an unnecessary domain hypothesis.
\end{remark}

\subsection{Smooth limits for ordinary root maps}
\label{sec:ordinary-root-calculus}

We now prove that the ordinary root maps introduced in
Section~\ref{sec:root-energy} have smooth limits, and hence that their
limiting insertions have the root-shift energy computed there.
At height $h$, write $u_h$ for the limiting root map and put
$v_{P,h}=u_h^*Pu_h\vac_h$ for a root Pauli observable $P$.
The following induction constructs these maps from the identity.
In the limiting cost update below, the field direction at height $h$ is the
preceding word's insertion on a child subtree at height $h-1$:
\[
 v_{Z,h-1}=u_{h-1}^*Zu_{h-1}\vac_{h-1},\qquad
 u_h\longmapsto e^{-\ii tZq(v_{Z,h-1})}u_h.
\]
Here $q(v_{Z,h-1})$ acts on the child Fock factor of $\cH_h$.
Once the heights include the fixed word's light cone, adding empty
descendants identifies these updates under the compatible embeddings.
We henceforth suppress the height indices, writing $v_P=u^*Pu\vac$;
a direction inside $q$ always comes from the child level.

\begin{lemma}[Ordinary root updates]\label{lem:ordinary-updates}
Ordinary finite words have smooth root limits and smooth adjoint
limits at every sufficient finite height.  Starting at $u=1$,
the limiting left updates for a mixer of angle $\beta$ and a cost of angle $t$
are, respectively,
\begin{equation}\label{eq:ordinary-root-updates}
 u\longmapsto e^{-\ii\beta X}u,
 \qquad
 u\longmapsto e^{-\ii tZq(v_Z)}u.
\end{equation}
The cost update leaves $v_Z$ unchanged.  For a cost conjugated
by a global mixer, replace $Z$ by the same rotated Pauli $P$
both at the root and in $v_P$.  A Gaussian seed gate has the
additional update $u\mapsto e^{-\ii tZs}u$.

All these assertions are unchanged when the top root has $D+1$
children and the descendants have $D$ children, still using
the coefficient $D^{-1/2}$. For an ordinary word $W$, its limiting
insertion $v_Z$ therefore satisfies $v_W=\cE(v_Z)$, with
$\cE$ defined in Proposition~\ref{prop:energy}.
\end{lemma}

\begin{proof}
At finite $D$, factoring out the new child costs gives the exact
left update
\[
 \exp\left(-\ii tZ D^{-1/2}\sum_{j=1}^D
                            U_j^*Z_jU_j\right).
\]
The one-copy observable is centered, because $Z$ is odd under
the joint spin/seed symmetry and the preceding word preserves
that symmetry.  Its vacuum column is $v_Z$.  On each root
$Z$-eigenspace, the update is a product of small one-child
unitaries.  Lemma~\ref{lem:unitary-clt} gives precisely the
cost in \eqref{eq:ordinary-root-updates}; the scalar correction
is zero for an exponential of a fixed observable.  The smooth
estimates follow inductively from the same lemma and
Lemma~\ref{lem:weyl-bounds}.  Mixers are local, and a seed
gate is multiplication by a two-by-two unitary; its input
derivatives have polynomial seed bounds.  Both therefore
preserve the smooth domain.  The update commutes with the
root $Z$, proving the assertion about $v_Z$; conjugation proves
the rotated-axis statement.

For the full-root assertion use $M=D+1$ in the occupation
decomposition while retaining $a/\sqrt D$ in each one-copy
factor.  In the proof of Lemma~\ref{lem:unitary-clt}, every
limiting coefficient is multiplied by $M/D\to1$, and every
bound is unchanged up to a constant.  Induction on the fixed
word proves that full-root and cavity-root maps have the same
limit. Padding the subtree by the finite light-cone radius
justifies all factorizations before limits are taken.

The finite-degree ordinary insertion vectors meet the hypotheses of
Proposition~\ref{prop:energy}: parity makes them centered, the
construction preserves rooted symmetry, and the smooth weight
dominates $1+N$. Their edge pairing is the expectation of $Z_iZ_j$.
The root-shift formula and Proposition~\ref{prop:tree-transfer}
therefore give $v_W=\cE(v_Z)$.
\end{proof}

\subsection{Substitution of commuting controls}
\label{sec:commuting-controls}

The ordinary cost update uses a scalar time multiplying one field.
A controlled operation instead inserts a commuting list of signal
operators into the parameters of a local matrix. The following lemma
makes that substitution legitimate,
including its derivatives and adjoints; no choice of independent
Gaussian coordinates is needed.

\begin{lemma}[Smooth commuting substitution]
\label{lem:commuting-substitution}
Let $S_D=(S_{D,1},\ldots,S_{D,b})$ be a fixed finite list of
strongly commuting self-adjoint operators.  Suppose their
polynomials converge smoothly and their joint exponentials satisfy,
for each $m$, bounds
\begin{equation}\label{eq:joint-exponential-bound}
 \|W_D^m e^{\ii t\cdot S_D}\xi\|
 \leq C_m(1+|t|)^{k_m}\|W_D^{r_m}\xi\|,
\end{equation}
uniformly in $D,t$, and converge strongly on smooth inputs.
Let $F_D(x)$ and $F_D(x)^*$ have uniform smooth bounds and
converge with all their jets on compact $x$-sets, and suppose
$F_D(x)$ commutes with the spectral projections
of every $S_{D,j}$.  Then the substitution $F_D(S_D)$ has uniform
smooth bounds and converges smoothly.  The assertion includes
extra scalar parameters and their jets.  If $F_D(x)$ is unitary
for real $x$, so is $F_D(S_D)$.
The limiting list $S$ consists of the strongly commuting
self-adjoint generators of the limiting joint unitary group.
Their actions and polynomials on the smooth domain agree with
the assumed polynomial limits.

For a real smooth function $f(s,x)$ whose derivatives have
polynomial growth in $x$, uniformly in $s$, the operators
$e^{\ii t f(s,S_D)}$ have polynomial smooth bounds in $t$.
Polynomial growth in a Gaussian $s$ is also permitted with the
seed-moment weight.  In every seed-dependent assertion, the
local seed multiplication algebra is required to strongly
commute with every $S_{D,j}$ and with the operator coefficient
families; thus $f(s,S_D)$ is a joint commuting substitution.
No nondegeneracy assumption on the joint
covariance of the $S_{D,j}$ is needed.
\end{lemma}

\begin{proof}
First identify the limiting signal operators.  Write
$U(t)=\lim_D e^{\ii t\cdot S_D}$ for the strong operator limit.
Strong limits of products of uniformly bounded operators preserve
the group law, and the limits at $t$ and $-t$ preserve inverses;
thus each $U(t)$ is unitary.  For smooth $\xi$,
\[
 \|(e^{\ii t\cdot S_D}-1)\xi\|
       \leq\sum_{\alpha=1}^b|t_\alpha|\|S_{D,\alpha}\xi\|.
\]
The uniform polynomial bounds give strong continuity on this dense
domain, hence on the full space.  The strongly continuous joint
unitary group therefore determines strongly commuting self-adjoint
generators $S_1,\ldots,S_b$.  Passing to the limit in the integrated
one-parameter group identities identifies each generator on smooth
vectors with the corresponding polynomial limit.  The smooth bounds
make this domain invariant, so iteration identifies all polynomial
limits.  This is the realization used in $F(S)$.

We give the estimates, which also specify the substitution on the
smooth domain when $F_D$ is unbounded. For such coefficients, commutation
with the spectral measures is understood in the closed-operator
(or decomposable-operator) sense. On the invariant smooth domain we use
the equivalent identities with the joint unitary group. Sharp spectral
projections are not required to preserve the smooth domain. These
identities pass to the limit because the moving inputs
$e^{\ii t\cdot S_D}\xi$ converge in every required graph norm, as do
the coefficients acting on them. Thus the limiting coefficients retain
the commutation needed for the Fourier products below. Suppress $D$.  Fix an
output weight $m$, and take $r_m,k_m$ from
\eqref{eq:joint-exponential-bound}.  Choose an integer
$q>k_m+b$.  For a sufficiently large integer $N$, the family
\[
 G(x)=(1+|x|^2)^{-N}F(x)
\]
and its first $q$ derivatives are integrable in $x$ as operators
from some fixed input graph norm to the $r_m$th output graph
norm.  Such an $N$ exists because only finitely many polynomial
growth exponents are involved.  With
$\widehat G(t)=(2\pi)^{-b}\int e^{-\ii t\cdot x}G(x)\,\dd x$,
integration by parts gives
\[
 \|W^{r_m}\widehat G(t)\eta\|
       \leq C(1+|t|)^{-q}\|W^r\eta\|.
\]
Therefore
\begin{equation}\label{eq:operator-fourier-substitution}
 F(S)\xi=\int_{\R^b} e^{\ii t\cdot S}\widehat G(t)
                (1+|S|^2)^N\xi\,\dd t
\end{equation}
converges absolutely in the $m$th graph norm and has a bound of
the form \eqref{eq:smooth-family-bound}.  The polynomial on the
right acts first; its bound can be absorbed by increasing the
input weight last.  This order of choices avoids any circular
dependence of exponents.

Integration by parts in $t$, with the same integrable graph-norm
bounds, shows that the result does not depend on $N$: multiplying
a symbol by $1+|x|^2$ corresponds to applying $1-\Delta_t$ to
its Fourier transform and hence to multiplication by
$1+|S|^2$ after integration.  The common smooth domain is
preserved by all factors.  Fubini's theorem and the convolution
formula show on that domain that
\[
              F(S)G(S)=(FG)(S).
\]
In this identity the order of $F$ and $G$ is retained; one only
commutes the Fourier factors and polynomials in $S$ past the
operator coefficients, as permitted by the hypothesis.  It can
first be checked for rapidly decreasing smooth symbols and
then for the divided symbols above, moving their polynomial
factors to the input.  Taking adjoints in pairings of smooth
vectors likewise gives $F(S)^*=F^*(S)$ on the common domain.
For scalar symbols this construction is the usual joint
spectral functional calculus by Fourier inversion.  In
particular the constant symbol $1$ substitutes to the identity.
Thus a pointwise unitary symbol and its adjoint substitute to
mutually inverse isometries on a dense invariant domain.  Their
extensions are unitary on the full Hilbert space.  These facts
also identify the construction with the finite-tree substitution
where the controls already have a simultaneous multiplication
representation.

For convergence, first restrict $x$ to a compact set.  Convergence
of the integrands on smooth inputs and the integrable bound just
proved give convergence of the Fourier integrals.  The polynomial
decay supplied by $(1+|x|^2)^{-N}$ controls the discarded tails
uniformly.  The same proof applies to every scalar-parameter
derivative and to adjoints.  Uniform higher output bounds then
give smooth convergence by \eqref{eq:weight-interpolation}.
For coefficients merely measurable in $s$, use dominated
convergence in $s$; for polynomial seed bounds increase the
seed-moment input weight.

Finally, every $x$-derivative of $e^{\ii t f(s,x)}$ is bounded
by a polynomial in $|t|$ times a fixed polynomial in $|x|$
(and $|s|$ in the moment-weight case).  Applying the same
argument proves the last assertion.  The proof used a joint
spectral representation and Fourier inversion, without choosing
independent field coordinates.
\end{proof}

An important supply of lists satisfying the hypotheses is the
following.  If $M_1,\ldots,M_b$ are commuting centered one-copy
signals and the one-copy joint exponentials have polynomial
smooth bounds, apply Lemma~\ref{lem:unitary-clt} to
$K(x)=\exp(\ii\sum x_\alpha M_\alpha)$.  Its scalar correction
vanishes, and hence
\begin{equation}\label{eq:commuting-signal-limit}
 \left(D^{-1/2}\sum_jM_\alpha^{(j)}\right)_{\alpha=1}^b
 \longrightarrow \bigl(q(M_\alpha\vac)\bigr)_{\alpha=1}^b.
\end{equation}
The limiting fields commute: the imaginary part of their Gram
matrix is the vacuum expectation of a commutator of the original
signals, divided by $2\ii$, and is zero.  Their polynomials
converge by Lemma~\ref{lem:centered-sums}.  The last assertion
of Lemma~\ref{lem:commuting-substitution} bootstraps polynomial
exponential bounds for smooth functions of such a list.  Starting
with local seed functions and a single commuting spin axis, this
proves the same statements for every fixed finite feedforward
calculation of commuting signals.

\subsection{Exclusion calculations and controlled gates}
\label{sec:exclusion-calculus}

We now prove the stronger calculus needed for synthesis.  Its
objects are temporary operations, not additional gates in the
final QAOA word.  The purpose of the calculus is to justify
replacing such operations by finite products inside an already
chosen surrounding circuit.

An even \emph{mask} is a measurable subset of $S$ invariant
under $\iota$.  For masks $I,J$ write
\begin{equation}\label{eq:masked-adjacency}
 (A_{IJ}f)_i=\1_I(s_i)D^{-1/2}
       \sum_{j\sim i}\1_J(s_j)f_j.
\end{equation}
Masks and the coefficients attached to them are part of a fixed
finite expression.  They are not optimized after sampling a
tree or its seeds.

\begin{definition}[Exclusion calculation and gate]
\label{def:exclusion-gate}
A finite exclusion calculation uses local seeds, one prescribed
commuting spin axis at each control site, smooth pointwise
functions of finite lists of messages, and normalized sums
\eqref{eq:masked-adjacency}.  An axis may depend on the site's
mask.  A formula is split into finitely many named roles when
different occurrences need different masks.  The following
conditions are imposed.
\begin{enumerate}
\item Every transmitted field is odd under simultaneous seed
reflection and spin flip.  Every pointwise function is smooth
in its aggregate inputs, with polynomial bounds for every
derivative.  Its local-seed coefficients are measurable and
uniformly bounded, unless the seed-moment variant is specified.
\item In every normalized sum, the receiving mask is disjoint
from every role mask occurring below that sum.  Thus all
ancestor receiving types are excluded from the descendant
calculation.  In particular a message sent across a cut does
not use the site on the other side of that cut.
\item A gate chooses a target mask and applies a two-by-two
unitary $V$ at each target, with parameters computed by such
calculations.  The matrix $V$ itself must be smooth in the
aggregate inputs, with polynomial bounds for every matrix
derivative, uniformly in the local seed (or with the stated
polynomial seed bound).  If external parameters are present,
the bounds include all their mixed jets, locally uniformly
in those parameters.  The target mask is disjoint from all control
masks used across an edge in those parameters.  At the target
itself only its seed, not its spin, is used as a parameter.
\item All parameter operators of this one gate belong to the
same commuting control algebra.  The gate, including its
two-by-two matrix, preserves the joint symmetry.
\end{enumerate}
The global gate is the product of these controlled matrices.
Its factors commute: their targets are distinct and none is a
control for another factor of the same gate.  A finite-range
local observable is therefore affected by only finitely many
of them at finite $D$.
\end{definition}

The exclusions make the information across a removed edge depend on
finitely many boundary inputs.  We next show that the resulting
finite recursion respects the weighted limits already proved.  The
replacement assertion is formulated inside an existing circuit,
which is the form required by the synthesis argument.

\begin{theorem}[Smooth gate calculus]\label{thm:gate-calculus}
Fix a finite sequence of exclusion gates, ordinary mixers and
costs, and local seed gates.  Its root interaction-picture maps,
their adjoints, and their root insertion vectors have smooth
limits under the occupation embeddings.  They have uniform
smooth bounds at every sufficient finite height.  The same
limit is obtained at a full root with $D+1$ children.  Hence
their centered insertion energies are given by
Proposition~\ref{prop:energy}.

An exclusion gate may be replaced inside any such fixed
sequence by a sequence of pointwise unitary matrices with the
same parameter calculations, provided every fixed number of
aggregate-input derivatives converges on compact sets and has
uniform polynomial bounds.  Convergence almost everywhere in
the local seed, or convergence in seed measure followed by a
subsequence, is sufficient for bounded measurable local-seed
coefficients.  The resulting root maps, adjoints, and insertion
vectors converge smoothly, and their energies converge.

If local-seed coefficients are uniformly bounded, these bounds
hold conditionally on the root seed, with constants independent
of its value, using $W_0=1$.  Root maps commute with root-seed
multiplication.  Consequently, if a root column $c$ has
conditional $m$th graph norm at most $C_m$, then
\begin{equation}\label{eq:masked-column-bound}
 \|W^m\1_Jc\|\leq C_m\sqrt{\mu(J)}.
\end{equation}
The same statements hold for any further scalar boundary-input
jets of a fixed exclusion calculation.  They also hold,
including adjoint jets, for a finite list of external unscaled
parameters, locally uniformly in those parameters, when the
matrix and signal hypotheses hold jointly for their mixed jets.
\end{theorem}

\begin{proof}
We give both the finite-tree recursion and the analytic
induction.  This is necessary because a global operator-norm
error proportional to the number of vertices would not imply
the assertion.

\paragraph{Boundary inputs.}
Unroll the fixed expression and index its normalized-sum occurrences
by their named roles, retaining repeated occurrences separately.
These indices form the finite boundary-channel list
$\alpha=1,\ldots,b$.  Cut an edge immediately above a branch root.
For an occurrence $A_{I_\alpha J_\alpha}f_\alpha$, replace the
parent's contribution to the normalized sum by a formal real
coordinate $y_\alpha$, retaining the receiving mask $I_\alpha$.
Write $y=(y_1,\ldots,y_b)$; these are substitution variables,
not additional random seeds.  A sender supplies the coordinate
$m_\alpha(y)=\1_{J_\alpha}(s)f_\alpha(y)$, evaluated at that
sender.  The receiving child uses $m_\alpha(y)/\sqrt D$ on its
role mask and ignores the coordinate when that role is inactive.
All outgoing messages that are
requested by the parent are unchanged by the cut: their
formulas exclude the parent's receiving type.  The parameters
of the gate are computed from the finite list of zero-boundary
child signals and these added inputs.  At the branch root this
gives a target matrix $V(y)$, equal to identity on nontargets,
and the outgoing message list $m(y)$.  The same list is
used for each child with its receiving masks imposed there.
Indeed, a component sent to a child excludes that child's
type from its own computation, so it contains no contribution
from that child.  The mask exclusions are needed for exactly
this assertion.

We first derive the factorization in the physical coordinates, before
conjugating by the preceding circuit. Retain the target operations at
distance at most $d$ from the branch root, but compute all their
parameter formulas on the padded tree. Let $G^{[d]}(y)$ be this
truncated operation, and temporarily write $V_{\rm phys}(y)$ and
$m_{\rm phys}(y)$ for the root matrix and outgoing messages just
defined. For $d=0$ the operation is only $V_{\rm phys}(y)$. For
$d\geq1$, partitioning the targets among the root and its children gives
\[
 G^{[d]}(y)=V_{\rm phys}(y)
       \prod_{j\ \mathrm{child}}
       G_j^{[d-1]}(D^{-1/2}m_{\rm phys}(y)).
\]
The same formula list is evaluated in each child with its receiving
masks imposed there. The exclusion rule ensures that the input sent
to a child does not feed back through that child's receiving type.

The zero-boundary child operation is the baseline to factor out. Put
\[
 B^{[d]}=\prod_jG_j^{[d-1]}(0),\qquad
 P=\bigotimes_jU_j,\qquad U=Pu,
\]
where $U$ is the preceding subtree circuit and $U_j$ its child
circuits. For $d=0$ take $B^{[0]}=1$. The residual operation in a
child, expressed in its preceding frame, is
\[
 K_j^{[d-1]}(a)
   =U_j^*G_j^{[d-1]}(0)^*G_j^{[d-1]}(a)U_j.
\]
At the current branch use the same definition,
\[
 K^{[d]}(y)=U^*G^{[d]}(0)^*G^{[d]}(y)U.
\]
Thus $K^{[d]}(0)=1$. This compares two operations on the same
branch; it does not discard the zero-input operation.

Conjugate the physical signals by $P$, writing
$V(y)=P^*V_{\rm phys}(y)P$ and $m(y)=P^*m_{\rm phys}(y)P$.
Equivalently, conjugate each zero-boundary child signal by its
$U_j$, while leaving the root spin in its present frame. Define
$E^{[d]}(y)=P^*(B^{[d]})^*G^{[d]}(y)P$. Factoring the baseline
from the physical identity gives the exact finite-degree recursion
\begin{align}
 E^{[d]}(y)
 &=V(y)\prod_{j\ \mathrm{child}}
            K_j^{[d-1]}(D^{-1/2}m(y)),
                         \label{eq:exclusion-E-recursion}\\
 K^{[d]}(y)
 &=u^*E^{[d]}(0)^*E^{[d]}(y)u.
                         \label{eq:exclusion-K-recursion}
\end{align}
For $d=0$ the first line is $E^{[0]}(y)=V(y)$. For $d=1$ it
multiplies the root matrix by the ratios of the child target
matrices. Notice that $m(0)$ need not vanish: removing the parent
input leaves the root's own signals. Hence $E^{[d]}(0)$ is generally
nontrivial, although every ratio $K^{[d]}$ equals one at zero.

The second line follows by canceling $B^{[d]}(B^{[d]})^*$ between
the two factors. At zero external input the same factorization gives
\[
 G^{[d]}(0)U=B^{[d]}P E^{[d]}(0)u.
\]
The new child circuits are $G_j^{[d-1]}(0)U_j$, so the new root
map is
\begin{equation}\label{eq:exclusion-root-update}
                    u_{\rm new}=E^{[d]}(0)u.
\end{equation}
Only targets within the finite influence distance can respond to a
cut input. Taking $d$ above that distance gives the actual root
update; all more distant zero-input operations are already factored
into the child circuits.

The factors and their parameter arguments commute before
the preceding circuit is applied: targets are disjoint from
controls, and each site supplies only its prescribed axis.
More explicitly, let $\mathcal A_j$ be the abelian von
Neumann algebra generated by all seeds and prescribed
control axes in child branch $j$, excluding target spin
axes.  For every scalar $a$, each truncated gate
$G_j^{[d]}(a)$ commutes with $\mathcal A_j$, since it
changes only target spins and has coefficients in that
algebra.  Consequently
\[
 K_j^{[d]}(a)\in\bigl(U_j^*\mathcal A_jU_j\bigr)'.
\]
Before the final root conjugation by $u$, every component
of $m(y)$ and every zero-boundary parameter sum is
affiliated with the abelian algebra generated by the
root controls and the child algebras
$U_j^*\mathcal A_jU_j$.  Thus $K_j^{[d]}(a)$ commutes
with every spectral projection of every such component,
even when that component uses signals from branch $j$
itself.  The unused components of a common message list
cause no exception.  This proves the exact commutation
required for substituting the operator $m(y)$ into the
scalar family.  All these algebras also commute with
root-seed multiplication, as do the preceding root
maps, so the seed-dependent version of commuting
substitution has its additional commutation hypothesis.
This also proves
that any two descriptions giving the same finite-tree
operation, including its cut versions, have the same root
map.

\paragraph{Zero-boundary signals.}
Assume by induction that the preceding maps $u,u^*$ have
the asserted smooth bounds and limits.  Compute a signal
in its physical control algebra from the bottom roles up.
At a sum, its zero-boundary child columns have already been
computed; after conjugation by the preceding child maps,
they still commute, are still centered by odd symmetry,
and have the required smooth bounds.  A signal at the new
root is calculated pointwise in the resulting child sums
and the root's chosen axis, and is then conjugated by the
preceding $u$ when its one-copy Heisenberg version is needed.

Here is the induction that also controls exponentials.
Local commuting seed/axis variables plainly have joint
exponentials with polynomial smooth bounds.  Suppose this
is true for the child signal list.  Lemma~\ref{lem:unitary-clt}
gives \eqref{eq:commuting-signal-limit}, including the
polynomial frequency bounds on joint exponentials, and
Lemma~\ref{lem:centered-sums} gives polynomial convergence.
Lemma~\ref{lem:commuting-substitution} then computes each
smooth pointwise function and gives polynomial bounds for
the joint exponentials of every new finite signal list.
Conjugating by $u$ and $u^*$ preserves those bounds and
convergence.  This proves the signal assertions in order
of calculation depth, without any assumption on an unknown
outgoing exponential.

\paragraph{Propagation of a gate across a cut.}
Proceed next in $d$ in
\eqref{eq:exclusion-E-recursion}--\eqref{eq:exclusion-K-recursion}.
For $d=0$, $V(y)$ is a smooth two-by-two matrix function of
the signal lists just treated; commuting substitution gives
all assertions.  Suppose they hold for the child family
$K^{[d-1]}(a)$, including all input jets and adjoints.  It
equals identity at $a=0$.  Treating the inputs as odd under
the symmetry gives
$\mathscr P K^{[d-1]}(a)\mathscr P=K^{[d-1]}(-a)$,
where $\mathscr P\vac=\vac$.  Therefore each first derivative
has vacuum mean zero.

For scalar $a$, apply Lemma~\ref{lem:unitary-clt} to the
child product.  With
$\partial_\alpha K^{[d-1]}(0)=\ii T_\alpha$, its limiting
value is
\begin{equation}\label{eq:exclusion-product-limit}
 \exp\left\{\ii\sum_\alpha a_\alpha q(T_\alpha\vac)
 +\frac12\sum_{\alpha,\lambda}a_\alpha a_\lambda
    \left\langle K_{\alpha\lambda}
       +\tfrac12\{T_\alpha,T_\lambda\}\right\rangle_0
 \right\}.
\end{equation}
The second term is scalar for scalar $a$.  After substitution
it can be an operator in the commuting control algebra; it
must therefore be retained.  The product lemma supplies
polynomial bounds for every scalar $a$-jet.  The components
of $m(y)$ are smooth functions of a finite list of commuting
normalized signals and a root control axis.  The finite-tree
commutations already established imply the hypotheses of
Lemma~\ref{lem:commuting-substitution}, which substitutes
$a=m(y)$.  Use the spectral projections of the fixed root
axis if needed.  Multiplication by $V(y)$ completes the
induction for $E^{[d]}(y)$, and conjugation in
\eqref{eq:exclusion-K-recursion} completes it for $K^{[d]}(y)$.
The argument includes all boundary-input derivatives by
the ordinary chain rule.

This proves \eqref{eq:exclusion-root-update} with smooth
bounds and limits.  It adds one gate to the preceding
sequence.  Ordinary cost and mixer updates may be inserted
between these gates by Lemma~\ref{lem:ordinary-updates};
its proof requires only smooth bounds on the preceding maps
and preservation of symmetry.  Local seed gates are handled
as in that lemma.  Induction on the finite sequence proves
the first assertion.  The same proof with $D+1$ top factors
instead of $D$ gives the same limit, since $(D+1)/D\to1$.
Smoothness supplies the bounded occupation required for
Proposition~\ref{prop:energy}.

External parameters can be carried through this induction
without being scaled.  At the child-product step use the
unscaled-parameter assertion of Lemma~\ref{lem:unitary-clt};
at the substitution step use the mixed-jet version of
Lemma~\ref{lem:commuting-substitution}.  Each unitary strong
limit is unitary, so adjoint convergence follows from the
identity in the proof of Lemma~\ref{lem:unitary-clt}; its
jet version follows by finite differences with the adjoint
jet bounds.  For real commuting signals the adjoint is the
same signal, and compositions use the reversed products of
the already convergent adjoint families.  Thus adjoint
convergence is established at each induction step, rather
than inferred from bounds alone.

\paragraph{Continuity inside a sequence.}
Repeat these two inductions with a further approximation
index.  Compact convergence of a fixed number of matrix
derivatives and uniform polynomial bounds are precisely
the hypotheses of commuting substitution.  The normalized
sum and product lemmas allow the child maps themselves to
vary with this index.  Hence every new root map and its
adjoint converge on smooth vectors with uniform smooth
bounds.  A fixed finite surrounding sequence is handled
by induction and the product rule.  The assertion does
not make a bound uniform in the length of that sequence.
For a pointwise product formula on a single target and a
common commuting control algebra, the entire product is
first regarded as one matrix-valued function; the theorem
then applies to its smooth matrix limit.  This is how
arbitrarily long local commutator or sum formulas can be
used without an estimate proportional to graph size.
For measurable seed limits, dominated convergence in the
root seed and then the same child induction give the
stated extension.  Norm continuity
\eqref{eq:energy-continuity} proves energy convergence.

\paragraph{Conditional bounds and masks.}
Take $W_0=1$.  Every root map commutes with root-seed
multiplication, by inspection of the finite-tree gates
and their factorizations.  Disintegrate the root Hilbert
space over that seed.  All preceding estimates remain
valid in each fiber with the same constants: coefficients
are uniformly bounded there, and the normalized child
columns are averaged over independent child seeds.
Moreover the weight commutes with root-seed multiplication.
Integrating the squared conditional estimate over $J$
proves \eqref{eq:masked-column-bound}.  With polynomial
seed coefficients use instead the seed-moment weight and
integrate the polynomial majorants.  This completes the
proof.
\end{proof}

The gate theorem controls a whole unitary.  In the helper
contractions we also need its first variation with respect to time:
this is a Hamiltonian response to boundary inputs.  The next
corollary derives that response from the same unitary limit, so that
the scalar second-order term and lower-level responses are retained.

\begin{corollary}[Additive cut responses]\label{cor:additive-response}
Fix a preceding finite sequence as in Theorem~\ref{thm:gate-calculus}.
Suppose a temporary Hamiltonian is diagonal in a common
collection of site axes, with coefficients given by an
exclusion calculation.  Require its local Hamiltonian
matrices and every aggregate-input derivative to have
polynomial bounds as in Definition~\ref{def:exclusion-gate}.
In a cut branch let $H_D(a)$ denote
the sum of the Hamiltonian terms affected by scalar
boundary input $a$, truncated beyond its finite influence
distance and conjugated by the preceding subtree circuit.
Assume explicitly that
\begin{equation}\label{eq:additive-response-parity}
                  H_D(-a)=\mathscr P H_D(a)\mathscr P,
                  \qquad \mathscr P\vac=\vac,
\end{equation}
where $\mathscr P$ is the joint spin/seed symmetry.
The first two boundary derivatives have smooth limits
$H_\alpha=\lim_D\partial_\alpha H_D(0)$ and
$H_{\alpha\lambda}=\lim_D\partial_\alpha\partial_\lambda H_D(0)$,
and $H_\alpha$ is centered.  These limits concern the response
derivatives; no limit for the extensive $H_D(0)$ is required.
The change in its $D$ child branches when zero boundary input is
replaced by $a/\sqrt D$ is
\[
 \sum_{j=1}^D
   \bigl[H_D(a/\sqrt D)-H_D(0)\bigr]^{(j)}.
\]
Its smooth limit is
\begin{equation}\label{eq:additive-cut-response}
 \sum_\alpha a_\alpha q(H_\alpha\vac)
 +\frac12\sum_{\alpha,\lambda}a_\alpha a_\lambda
                           \langle H_{\alpha\lambda}\rangle_0.
\end{equation}
The correction is realized as the self-adjoint generator of the
limiting time group of the child cut ratios; the display specifies
its action on the smooth domain.  We use the same realization
after commuting substitution and for the correction generators below.
The same formula applies recursively to lower cut responses
and to their input jets.  It is continuous with the smooth
bounds of Theorem~\ref{thm:gate-calculus}, for a fixed finite
list of channels.  Commuting parameter operators may replace
$a$ under the same exclusions.

In particular, if a sequence of such correction generators
$R_j$ satisfies $\|R_j\xi\|\leq\delta_j\|W^r\xi\|$ with
$\delta_j\to0$, then
\begin{equation}\label{eq:small-generator-error}
 \|(e^{\ii tR_j}-1)\xi\|
       \leq |t|\delta_j\|W^r\xi\|.
\end{equation}
If their exponentials and adjoints also have uniform smooth
bounds, these errors tend to zero in every fixed lower
graph norm and may be inserted in a fixed surrounding
sequence.
\end{corollary}

\begin{proof}
Differentiating \eqref{eq:additive-response-parity} at zero
gives $\mathscr P H_\alpha\mathscr P=-H_\alpha$, proving
centering.  To justify all Hamiltonian-jet bounds without
an estimate on an extensive Hamiltonian, introduce an
unscaled real time $t$ and use the unitary cut ratio
\[
 K_D(a;t)=\exp\bigl(\ii t[H_D(a)-H_D(0)]\bigr).
\]
All current Hamiltonian terms commute before conjugation,
and simultaneous conjugation preserves this property.
In particular $H_D(a)$ and $H_D(0)$ commute.  Thus the
display is exactly the prior-frame ratio of the zero-
boundary gate and the gate with boundary input $a$.
A local matrix $e^{\ii t h(x)}$ and all its mixed $t,x$
derivatives have polynomial bounds from the assumed
bounds on $h$ (use Duhamel's formula for matrix derivatives).
The external-parameter version of
Theorem~\ref{thm:gate-calculus}, which also applies to its
truncated cut recursions, therefore gives uniform smooth
bounds and convergence of $K_D(a;t)$ with all time and
boundary-input jets.  In particular
\[
 H_D(a)-H_D(0)=\frac1\ii\partial_tK_D(a;0),\qquad
 \partial_a^\alpha H_D(a)
   =\frac1\ii\partial_a^\alpha\partial_tK_D(a;0)
       \quad(|\alpha|\geq1).
\]
These are the required bounds for the response jets;
no bound for the possibly extensive $H_D(0)$ is asserted
or used.

For the child-product limit compute the first two boundary
jets at zero:
\[
 K_\alpha(t)=\ii t H_\alpha,\qquad
 K_{\alpha\lambda}(t)=\ii t H_{\alpha\lambda}
                -\frac{t^2}{2}\{H_\alpha,H_\lambda\}.
\]
In Lemma~\ref{lem:unitary-clt} one has $T_\alpha=tH_\alpha$;
the two anticommutator terms cancel exactly.  The product
of the $D$ child ratios therefore tends, with all time
jets, to the exponential of $\ii t$ times
\eqref{eq:additive-cut-response}.  Differentiating at $t=0$
establishes the asserted smooth Hamiltonian limit itself.
The finite-degree child products are unitary groups in $t$.
Their strong limits at positive and negative times retain the
group law and inverses, while convergence of time jets gives
strong continuity.  Their self-adjoint generators consequently
act on smooth vectors by \eqref{eq:additive-cut-response};
no separate choice of closure of the displayed expression is used.
The exclusions and Lemma~\ref{lem:commuting-substitution}
allow operator inputs and preserve the group law and this
generator identification.  Repeating this argument at the
finitely many affected levels proves the recursive assertion.

For $R_j$ with this self-adjoint realization, the spectral inequality
$|e^{\ii tx}-1|\leq|t x|$ proves
\eqref{eq:small-generator-error}.  Uniform smooth
bounds and \eqref{eq:weight-interpolation} then give
higher-norm convergence; applying the same reasoning
to the adjoints and using finite composition proves
the final statement.
\end{proof}

\begin{remark}[Scope of uniformity]\label{rem:calculus-uniformity}
The constants in the calculus are uniform in degree and
in the stated pointwise approximation limits, for a
fixed finite calculation and fixed surrounding sequence.
They are not automatically uniform when the number of
roles or channels tends to infinity.  In a refinement
argument, such as the helper contraction below, that
additional uniformity must be proved from estimates on
the combined columns or kernels.  Once it is proved,
the same smooth-domain and continuity arguments apply.
Scalar phases depending only on the root seed commute
with root observables and do not change insertion
vectors or the child conjugations in the cut recursion.
\end{remark}

\section{A near-optimal bounded classical formula}
\label{sec:classical-optimization}

We now construct a bounded classical output whose rooted energy is
arbitrarily close to $P_*$. The construction has two parts. A scalar
Parisi diffusion first gives a finite function of independent Gaussians
with an explicitly computable value. We realize those Gaussians as
positions in the recursive space of
Section~\ref{subsec:classical-gaussian-sector}. Bounded approximation
and type exclusions then put the formula in the class required for
quantum synthesis. These last two operations will also serve the
independent positive-temperature construction in
Appendix~\ref{sec:positive-temperature-alternative}.

\subsection{The scalar input}

We use the following results from the complete companion
\emph{Full support of the zero-temperature Sherrington--Kirkpatrick
order parameter}~\cite{SKSupport2026}. We state the normalization and
the finite formula explicitly because the rooted realization below
depends on their exact form.

Let $\mathcal U$ consist of the nonnegative, nondecreasing,
right-continuous functions $\gamma:[0,1)\to[0,\infty)$ with
$\int_0^1\gamma(t)\,\dd t<\infty$. Define the Parisi value function
by the controlled heat equation
\[
 \Phi_t=-\tfrac12(\Phi_{xx}+\gamma(t)\Phi_x^2),
 \qquad \Phi(1,x)=|x|,
\]
in its integrable-coefficient control interpretation, and put
\[
 \mathcal P(\gamma)=\Phi_\gamma(0,0)
                 -\tfrac12\int_0^1t\gamma(t)\,\dd t.
\]
The ground-state formula and attainment of Auffinger and
Chen~\cite[Theorem~1 and the proof of Lemma~3]{AC2017} identify
$\min_{\gamma\in\mathcal U}\mathcal P(\gamma)$ with the constant
$P_*$ in~\eqref{eq:parisi-constant}. Fix any such minimizer.
Throughout this section $\gamma(t)$ is this scalar coefficient,
not a vector of cost angles.

For clarity, the covariance of the unordered-pair Hamiltonian is
$\E h_{n,J}(\sigma^1)h_{n,J}(\sigma^2)=nR_{12}^2/2-1/2$, where
$R_{12}=n^{-1}\sum_i\sigma_i^1\sigma_i^2$. Adding a centered
Gaussian of variance $1/2$, constant across configurations, gives
covariance $n\xi(R_{12})$ without changing the expected maximum.
This is the normalization of the cited formula.

Write $u=\Phi_{\gamma,x}$, $a=\Phi_{\gamma,xx}$ and let
\[
 \dd X_t=\gamma(t)u(t,X_t)\,\dd t+\dd B_t,\qquad X_0=0.
\]
The companion proves full relative Stieltjes support on $[0,1)$,
with the convention $\mu_\gamma([0,t])=\gamma(t)$, and gives
\begin{equation}
 \E u(t,X_t)^2=t,\qquad
 \E a(t,X_t)^2=1,\qquad
 P_*=\int_0^1\E a(t,X_t)\,\dd t.
 \label{eq:scalar-consistency-value}
\end{equation}
Here $u$ is odd, $a$ is even and nonnegative, and $|u|\le1$.
All their spatial derivatives are bounded on every fixed strip
ending before one. These are
Theorem~1.1, Lemmas~2.1--2.2 and Corollary~7.1 of
\cite{SKSupport2026}; their complete analytic proofs belong to that
companion, not to the rooted-space argument below. In particular,
Cauchy--Schwarz gives
\begin{equation}
 0\le P_*-\int_0^T\E a(t,X_t)\,\dd t\le1-T,
 \qquad 0<T<1.
 \label{eq:scalar-terminal-loss}
\end{equation}
No smooth density of $\mu_\gamma$ or bound uniform as $T\uparrow1$
is used here.

The use of the Parisi gradient and a normalized Euler martingale
follows the classical constructions of
Montanari~\cite[Section~3.1]{Montanari2019} and El Alaoui,
Montanari, and Sellke~\cite[Proposition~5.3]{AMS2020}.
The precise shifted-coefficient identity needed here is the scalar
companion input stated below.

For fixed $0<T<1$ and an integer $N\ge1$, set $\Delta=T/N$ and
$t_j=j\Delta$. On independent standard normal variables
$Z_1,\ldots,Z_{N+1}$, define
\begin{align}
 Y_0&=0,\qquad
 Y_j=Y_{j-1}+\sqrt\Delta\,Z_j
       +\Delta\gamma(t_{j-1})u(t_{j-1},Y_{j-1}), \label{eq:scalar-euler}\\
 c_j&=\bigl(\E a(t_{j-1},Y_{j-1})^2\bigr)^{-1/2},
 \qquad
 g_j=c_j a(t_{j-1},Y_{j-1})Z_j,\qquad 1\le j\le N.
 \label{eq:scalar-columns}
\end{align}
For all sufficiently large $N$, each $c_j$ is finite and positive.
With $\Phi_{\rm G}$ the standard normal distribution function, put
\begin{align}
 H_N&=u(T,Y_N)+2\Phi_{\rm G}(Z_{N+1})-1,
 &F_N&=\operatorname{sgn}(H_N), \label{eq:scalar-rounding}\\
 A_j&=\E[F_Ng_j]\quad(1\le j\le N),&
 B_i&=\E[F_NZ_i]\quad(1\le i\le N+1),\qquad A_0=0 .
 \label{eq:scalar-coefficients}
\end{align}
Conditional on $Z_1,\ldots,Z_N$, the added variable
$2\Phi_{\rm G}(Z_{N+1})-1$ is uniform on $(-1,1)$. Since $|u|\le1$,
\[
 \E[F_N\mid Z_1,\ldots,Z_N]=u(T,Y_N).
\]
Thus sign rounding preserves the terminal gradient as its conditional
mean. The same uniformity makes the value of the sign at zero
immaterial. Proposition~7.2 of~\cite{SKSupport2026} states that
\begin{equation}
 \lim_{N\to\infty}\sum_{j=0}^N A_jB_{j+1}
       =\int_0^T\E a(t,X_t)\,\dd t.
 \label{eq:scalar-coefficient-limit}
\end{equation}
The mesh limit in this assertion is taken at fixed $T$; only
afterwards is $T$ allowed to approach one. This finite scalar
identity, rather than a finite-degree message-passing theorem, is
the input to the next construction.

\subsection{Exact realization in the rooted Gaussian space}

Recall that a real centered child vector $f$ defines a Gaussian
position $\mathsf G(f)=q(f)$ at its parent, with covariance
$\E[\mathsf G(f)\mathsf G(g)]=\braket{f}{g}$.
The parent seed is independent of these child-Fock coordinates.
These two facts allow the scalar variables above to be constructed
successively at finite height.

\begin{lemma}[A finite Gaussian formula and its rooted energy]
\label{lem:gaussian-rooted-realization}
Fix $T,N$ and the scalar data in
\eqref{eq:scalar-euler}--\eqref{eq:scalar-coefficients}, with all
$c_j$ finite and positive. There is a real odd finite descendant
formula $F_N$ of height at most $N+1$ in the classical rooted space
whose Gaussian coordinates have exactly the stated joint law and
whose energy satisfies
\begin{equation}
 \cE(F_N)=\sum_{j=0}^N A_j B_{j+1}.
 \label{eq:scalar-rooted-energy}
\end{equation}
The formula is bounded by one, although its transmitted columns
need not be bounded.
\end{lemma}

\begin{proof}
Let $S$ be a standard Gaussian in the local root-seed factor and
start with $g_0=S$. Recursively set
\[
 Z_{j+1}=\mathsf G(g_j),\qquad 0\le j\le N,
\]
and use~\eqref{eq:scalar-euler}--\eqref{eq:scalar-columns} to form
$Y_j,g_j$ as soon as $Z_j$ is available. The constants $c_j$ here
are those defined under the independent Gaussian law above.
We will verify that the recursively constructed variables have
precisely that law.

At each step identify all earlier vectors by the compatible
empty-descendant embeddings of Section~\ref{sec:foundations}.
These embeddings preserve the earlier Gaussian positions:
lifting an embedded child vector gives the embedded copy of its
previous lift. Thus the recursion uses the same earlier variables
at the next height. Inductively $g_j$ has height at most $j$ and
$Z_{j+1}$ has height at most $j+1$. All vectors can be compared
in the same finite-height space whenever an inner product is taken.

We show inductively that $g_0,\ldots,g_j$ are real centered
orthonormal vectors. The claim holds for $g_0$. If it holds through
$g_{j-1}$, the positions $Z_1,\ldots,Z_j$ are jointly centered
Gaussian with identity covariance, hence independent standard
normals. For $j\ge1$, the coefficient
$c_j a(t_{j-1},Y_{j-1})$ is measurable in
$Z_1,\ldots,Z_{j-1}$. Conditional mean zero of $Z_j$ proves that
$g_j$ is centered and orthogonal to all earlier $g_i$ with $i\ge1$.
The normalization in~\eqref{eq:scalar-columns} gives
$\norm{g_j}=1$. The receiving-root seed $S$ belongs to a different
tensor factor from every child-Fock position, so it is independent
of all these positions. In particular
$\braket{g_0}{g_j}=0$.
The child directions themselves may contain seeds at their own
roots; this does not change the independence at the receiving root.
This closes the induction and also proves that
$Z_1,\ldots,Z_{N+1}$ have the required joint law. The coefficients
are bounded at the fixed grid times, so every constructed column
has moments of every finite order.

The map $u(t,\cdot)$ is odd and $a(t,\cdot)$ is even. Consequently
each $Y_j$ and $g_j$ is odd under simultaneous seed reversal.
Since $\Phi_{\rm G}(-z)=1-\Phi_{\rm G}(z)$, the same is true of
$F_N$ outside its null zero set. Hence $F_N$ is centered.

It remains to compute its energy, not just its distribution.
The vector $F_N$ is constant in the root-seed factor and depends
only on the finite orthonormal family of Gaussian modes
$\mathsf G(g_0),\ldots,\mathsf G(g_N)$.
The adjoint root shift extracts precisely its constant-root-seed
one-particle component. Since $q(g)=a^*(g)+a(g)$ has variance
$\norm g^2$, this component is
\[
 \cL^*F_N=\sum_{j=0}^N\E[F_NZ_{j+1}]\,g_j
        =\sum_{j=0}^N B_{j+1}g_j.
\]
There are no other first-chaos coefficients: $F_N$ is independent
of Gaussian modes orthogonal to the displayed family.
Moreover $F_N$ is independent of the receiving-root seed, so
$\braket{F_N}{g_0}=0=A_0$. Inner products with the other $g_j$
are exactly the scalar coefficients $A_j$. All vectors are real,
and therefore
\[
 \cE(F_N)=\re\braket{F_N}{\cL F_N}
        =\re\braket{\cL^*F_N}{F_N}
        =\sum_{j=0}^N A_jB_{j+1}.
\]
This proves the claimed identity, with the Gaussian variance and
root-shift normalizations used throughout the paper.
\end{proof}

Combining~\eqref{eq:scalar-terminal-loss},
\eqref{eq:scalar-coefficient-limit} and
\eqref{eq:scalar-rooted-energy} produces a bounded near-optimal
rooted vector by first choosing $T$ and then a finite $N$.
We next make its output smooth. At these fixed parameters set
$F_{N,k}=\tanh(kH_N)$. The null-set observation after
\eqref{eq:scalar-coefficients} and dominated convergence give
$F_{N,k}\to F_N$ in $L^2$. The root-shift continuity estimate is
\begin{equation}
 |\cE(f)-\cE(g)|
 \le(\norm f+\norm g)\norm{f-g}.
 \label{eq:energy-l2-continuity}
\end{equation}
Both outputs have norm at most one. Thus for every
$\varepsilon>0$ we may choose finite $T,N,k$ so that
\begin{equation}
 z=F_{N,k},\qquad |z|\le1,\qquad
 \cE(z)\ge P_*-\varepsilon.
 \label{eq:scalar-near-optimal-smooth}
\end{equation}

For the already fixed $T,N,k$, every transmitted-column map and
the final output map is jointly odd and smooth, with all
derivatives of polynomial growth in its finite list of odd
inputs. Indeed the Euler maps use finitely many fixed-time
bounded spatial derivatives of $u,a$, followed in a column by
multiplication by one Gaussian coordinate. Local intermediate
operations may be composed into those maps. The final sign
smoothing has bounded derivatives of every order at fixed $k$.

Two further operations are still necessary: bounding every
transmitted column and imposing type exclusions. We perform
them directly in the limiting Hilbert space. No finite-degree
edge-energy convergence is asserted for the discontinuous raw
$F_N$.

\subsection{Bounded trigonometric approximation}

For use with either classical construction, let
$\pi(x)=\max(-1,\min(x,1))$. The next lemma applies to any finite
descendant formula in its stated class. In the present application
$\pi(z)=z$ by~\eqref{eq:scalar-near-optimal-smooth}.

\begin{lemma}[Bounded trigonometric approximation]
\label{lem:classical-trigonometric-density}
Let $m^*=\pi(z)$, where $z$ is a finite descendant formula with
Gaussian leaf seeds.  Assume that every transmitted child column
is odd, and that each local map producing such a column or the
output is a jointly odd smooth function of its local seed and
finite list of odd aggregate inputs, with polynomial-growth
derivatives.  Purely local intermediate operations may be composed
into these maps before the lemma is applied.
For every $\eta>0$, there is an odd bounded angle formula $\theta$
such that
\begin{equation}
 \norm{\sin(2\theta)-m^*}<\eta.
 \label{eq:angle-density}
\end{equation}
Every sent column in the formula for $\theta$ can be chosen bounded
and odd.  Its pointwise maps can be chosen to be finite real
trigonometric sums in their finite lists of odd inputs, with odd
joint parity for odd outputs.  Its local seed inputs can be chosen
bounded and odd.  All these functions have bounded derivatives of
every order in their continuous arguments.
\end{lemma}

\begin{proof}
Replace a Gaussian local seed $s$ by a smooth bounded odd truncation
$s_R$ satisfying $|s_R|\le|s|$ and $s_R\to s$ pointwise.  It converges
in every finite moment.  At a fixed finite depth this changes the
original formula by an amount tending to zero in every required
finite moment: use \eqref{eq:classical-gaussian-isometry} at each
aggregation, Gaussian moment bounds for its finite list of fields,
and the polynomial bounds for the derivatives of the pointwise maps.

Fix $0<\kappa<1$.  The desired final angle
\[
 \theta_\kappa=\tfrac12\arcsin((1-\kappa)m^*)
\]
is bounded and odd, and is a Lipschitz function of the clipped output.
It satisfies $\sin(2\theta_\kappa)=(1-\kappa)m^*$.
We approximate this angle and all its subformulas beginning at the
output and working toward the leaves.  For the finite-dimensional
ideal input law at a pointwise operation, a continuous polynomial-growth
function has an arbitrarily accurate $L^p$ approximation, for any
fixed finite $p$, by a bounded smooth function with bounded
derivatives.  This follows by a smooth cutoff outside a large box
and mollification; the uniform moments control the cutoff error.
Symmetrizing by $F(x)\mapsto(F(x)-F(-x))/2$ preserves oddness and
does not increase its error under the symmetric input law.

Once an outer map has been replaced by a fixed bounded Lipschitz
map, perturbations of its inputs propagate continuously in $L^2$.
For its aggregate inputs this is exactly
\eqref{eq:classical-gaussian-isometry}, jointly for the finite list
of child columns.  We can therefore choose the approximations at
the next lower level sufficiently accurately to preserve the
already chosen output accuracy.  There are finitely many operations,
so induction supplies a formula in which every transmitted scalar
function is bounded, smooth, and odd.  No uniform bound on Lipschitz
constants over the sequence of approximation choices is needed;
each outer choice is fixed before its inputs are approximated.

Finally a bounded smooth cutoff supported strictly inside a box can
be extended periodically to a larger box.  Fourier approximation of
this smooth periodic function gives finite trigonometric sums
converging uniformly, with any prescribed finite number of
derivatives, on that box.  To control the complement, first fix the
box so that its probability is small and choose the Fourier
approximation uniformly on the entire period; both the cutoff and
its approximant then have a common finite sup norm.  The same
output-to-leaf continuity argument applies.  Odd symmetrization
leaves only sine terms $\sin(\lambda\cdot x)$ for an odd scalar
output.  This constructs the asserted trigonometric formula.
Choose its angle error smaller than $\eta/4$ and $\kappa<\eta/2$.
Since $x\mapsto\sin(2x)$ is $2$-Lipschitz, the resulting output
satisfies \eqref{eq:angle-density}.
\end{proof}

\subsection{The synthesis interface}

A \emph{mask} is an event in the local seed space invariant under seed
sign reversal.  Masks at different vertices are evaluated on their
respective seeds.  We use a finite partition into masks and write
\(p_R\) for the probability of a mask \(R\).  Probabilities used as
denominators are positive.  For disjoint masks \(R,S\), let
\[
 (A_{RS}f)_i=\1_R(i)D^{-1/2}
       \sum_{j\sim i}\1_S(j)f_j .
\]
This notation first denotes a finite-tree operation.  Its rooted limit
is provided by Theorem~\ref{thm:gate-calculus}.

The scalar formulas considered here have finite expression trees.
Their leaves are bounded measurable functions of the local seed.
An internal operation is either a pointwise function of finitely many
already constructed quantities or an aggregation \(A_{RS}\).
Pointwise functions are smooth in their aggregate arguments, with
polynomial bounds on every derivative, uniformly in the local seed.
Local seed dependence itself need not be smooth.  All quantities sent
across edges are odd under simultaneous reversal of seeds and odd spin
axes.  In the classical formulas there are no spin arguments, so this
is simply joint seed parity.  Sent nonlinear functions, and the output
angles, are bounded.  The aggregate fields need not be bounded.

\begin{definition}[Exclusion formulas]
\label{def:exclusion-formulas}
An aggregation from \(S\) to \(R\) is \emph{excluded} if the formula
sent from \(S\) uses no vertex in \(R\) anywhere below that
aggregation.  The formulas used in this section satisfy this
requirement at every aggregation and retain all ancestor exclusions:
the masks along any branch of an expression tree are distinct.
Different occurrences of a formula with different exclusion lists are
regarded as different formulas.  A gate target is excluded from all
cross-edge dependencies of its parameters.
\end{definition}

The distinction between an expression tree and the underlying regular
tree is useful.  A formula may reuse a mask in different expression
branches.  The exclusion rule forbids its reuse along a single active
branch.  It therefore forbids an immediate return across a cut, without
requiring independent values for two formulas that happen to use the
same descendant subtree.

\subsection{Colors, exclusions, and helpers}

A descendant formula omits the parent \emph{vertex} in its finite-tree
interpretation. We now impose the stronger type exclusions needed
for synthesis. We estimate their cost in the limiting classical space,
where \eqref{eq:classical-gaussian-isometry} makes the cost of deleting
small-probability types explicit.

Let $s_{\mathrm{raw}}$ denote the single standard Gaussian available
as the local control seed. Its sign is an independent fair sign, and
the distribution function of its absolute value is uniform on $(0,1)$.
Splitting the
binary digits of that uniform variable into disjoint subsequences
produces independent uniform variables, apart from a null set of
dyadic ambiguities.  Use one for a finite color variable $C$, and
another, together with the original sign, to generate a fresh
standard Gaussian $S$. This defines a measurable map
\begin{equation}
 \begin{gathered}
 s_{\mathrm{raw}}\longmapsto
       \bigl(C(s_{\mathrm{raw}}),S(s_{\mathrm{raw}})\bigr),\\
 C(-s_{\mathrm{raw}})=C(s_{\mathrm{raw}}),\qquad
 S(-s_{\mathrm{raw}})=-S(s_{\mathrm{raw}}),
 \end{gathered}
 \label{eq:raw-seed-reencoding}
\end{equation}
whose image has the product law of the chosen color and a standard
Gaussian. Every colored formula below is evaluated at this pair at
each vertex: its Gaussian input is $S(s_{\mathrm{raw}})$.
The necessary color masks are bounded measurable even functions of
the original local seed.  Smoothness below refers to the continuous
message arguments conditional on a color; no smoothness of a mask
as a function of the original seed is asserted or required.

One may first adjoin an independent color to the color-free seed
space and then pull back by \eqref{eq:raw-seed-reencoding}.
The local pullback $f(C,S)\mapsto
f(C(s_{\mathrm{raw}}),S(s_{\mathrm{raw}}))$ preserves the product
law, inner products, and the constant seed vector. Tensoring it at
the root and applying the symmetric Fock functor to its centered
child restriction extends it to every recursive height. These
isometries respect the empty-descendant embeddings, preserve parity,
and intertwine $\cL$, which inserts a constant seed at the new root.
Thus both adjoining colors and pulling the entire formula back to
the raw seeds preserve its rooted energy. Partition the colors
into nonhelper types $1,\ldots,M$ and helper types.  Write
\[
 \max_{1\le c\le M}\Pr(C=c)\le\alpha,
 \qquad \Pr(C\text{ is a helper})=p_H.
\]
Every helper type may have a prescribed positive probability, subject
to the total being sufficiently small.  The number of helper types
needed later is finite, at most one per nested propagation depth.

\begin{lemma}[Type exclusion at small cost]
\label{lem:classical-color-exclusion}
Fix a descendant formula in which every transmitted column and the
output are given by jointly odd bounded smooth functions of the
local odd seed and the odd aggregate inputs, with bounded first
derivatives.  Let its
propagation depth be $h$.  It has a finite realization by colored
formulas such that every aggregation excludes the current color,
all colors of the path ancestors, and all helper colors.
If $f$ is its original output and $f^{\mathrm{col}}$ its colored
output, set to zero on helper roots, then
\begin{equation}
 \norm{f^{\mathrm{col}}-f}^2\le C(\alpha+p_H),
 \label{eq:color-exclusion-bound}
\end{equation}
where $C$ depends only on the fixed color-free formula, not on the
number or individual probabilities of its nonhelper colors.
Conditional bounds for the unchanged nonlinear pointwise functions
and their derivatives in the original grouped aggregate coordinates
are uniform over these colors.  No bound independent of $M$ is
asserted for the Euclidean derivative norm after flattening all
source-color fields into separate coordinates.
\end{lemma}

\begin{proof}
Unroll the fixed expression into a finite expression tree.  For
each occurrence, each possible current nonhelper color $c$, and each
possible set $A$ of ancestor colors, make a copy of the lower
formula.  Its aggregates retain only children whose colors avoid
$A\cup\{c\}$ and all helpers.  At a retained child of color $d$,
the recursive copy uses ancestor set $A\cup\{c\}$ and current color
$d$.  Thus its own aggregates exclude $A\cup\{c,d\}$.  All sets
encountered have size at most $h+1$.  There are finitely many copies
because the original depth and the color partition are finite.
The aggregation at color $c$ is the sum of its source-color
aggregations.  Linearity of $\mathsf G$ identifies this finite sum
with the lift of the corresponding masked child column.

These copies also define an orientation-free all-neighbor rule.
When a recursive call crosses from a vertex $i$ of color $c$ to a
vertex $j$ of color $d$, the ancestor list at $j$ contains $c$.
Every aggregate at $j$ therefore gives the neighbor $i$ a zero
mask. Its sum over all neighbors equals its sum over neighbors other
than $i$, with the same normalization $D^{-1/2}$. Induction over the
expression depth identifies the all-neighbor calculation with the
descendant calculation on each active branch. The outer root has no
parent and uses all its neighbors. In particular the finite rule uses
the aggregations $A_{RS}$ of the synthesis interface without requiring
a parent orientation. A full root has $D+1$ neighbors; its limiting
formula agrees with the $D$-child convention by the full-root assertion
of Theorem~\ref{thm:gate-calculus}.

Here is a uniform error estimate.  Give the original formula the
extra independent colors but ignore them in its computation.
At one aggregate, let $g$ be its ideal bounded child column, with
$|g|\le K$, and let $g^{A,d}$ be the recursively modified column on
a retained color $d$.  Couple the two fields using
\eqref{eq:classical-gaussian-isometry}.  The modified aggregate has
child column
\[
 \widetilde g
 =\sum_{d\notin A\cup\{c\},\ d\ {\mathrm{nonhelper}}}
            \1_{\{C=d\}}g^{A\cup\{c\},d}.
\]
The color masks are orthogonal, and $g$ is independent of the colors
at its root.  Hence
\begin{align*}
 \norm{\widetilde g-g}^2
 &\le K^2\bigl(p_H+(|A|+1)\alpha\bigr)\\
 &\quad+\max_{d\notin A\cup\{c\}}
       \E\bigl[|g^{A\cup\{c\},d}-g|^2\mid C=d\bigr].
\end{align*}
Every masked column is centered: color masks are even and the
conditional formula is odd in the Gaussian seeds.  Thus the
Gaussian-isometry identity applies without a mean correction.
For a list of $k$ aggregates the squared Euclidean input error is
the sum of their squared errors.  A pointwise map with Lipschitz
constant $K'$ multiplies that sum by at most $(K')^2$.
Inducting over the finite expression depth, with the unchanged
bounded local seed as the base case, proves a conditional squared
error bounded by $C(\alpha+p_H)$ uniformly over current and
ancestor colors.  Setting the helper-root output to zero adds at
most $\norm f_\infty^2p_H$.

This estimate has the fan-in of the original formula, not the number
of color copies: disjoint child masks were combined by squared
norms before applying a Lipschitz bound.  Each nonlinear pointwise
map still receives the original finite list of grouped fields: if
it originally had $k$ field arguments, its colored version is
$F(S,Y_1^{\mathrm{col}},\ldots,Y_k^{\mathrm{col}})$, where
\[
 Y_r^{\mathrm{col}}=\sum_d\mathsf G(\1_{\{C=d\}}g_{r,d})
       =\mathsf G\left(\sum_d\1_{\{C=d\}}g_{r,d}\right).
\]
The conditional derivative bounds in these $k$ arguments are
unchanged.  The finite sums defining grouped fields are additional
linear operations and can be unbounded; boundedness is required of
the transmitted columns and final angles, not of Gaussian fields.
For a trigonometric formula, substitute each grouped linear sum into
the next pointwise map.  A trigonometric term in linear combinations
of the source fields is still a trigonometric term, so this substitution
preserves the finite trigonometric form without retaining separate
linear pointwise maps.
For the synthesis interface, a rotation by a grouped field is
exactly the product of the source-field rotations on the same axis,
since these classical fields commute.  If a compilation instead
regards all the source fields as separate coordinates, its
derivative constants may depend on the fixed number $M$.  This is
harmless: $M$ and all these finite lists are fixed before any helper
refinement or quantum approximation is taken.  Only the classical
error constant in \eqref{eq:color-exclusion-bound} and the bounds in
the original grouped coordinates are asserted to be independent
of $M$.
\end{proof}

\begin{proposition}[Classical near-optimality in synthesis form]
\label{prop:classical-near-optimum}
For every $\varepsilon>0$ and every $\eta_H>0$, there are a finite
color partition with helper probability $0<p_H<\eta_H$, a finite
number of helper types, and finite bounded odd angle formulas
$\theta_c$ on the nonhelper colors such that the following hold.
Each transmitted column is bounded and odd, every pointwise map
may be chosen as a finite trigonometric sum in its odd inputs, and
every aggregation excludes its current color, all path-ancestor
colors, and all helpers.  Conditional on the finite labels, all
continuous derivatives of the pointwise maps are bounded.  The
output
\begin{equation}
 m^{\mathrm{col}}
  =\sum_{c=1}^M\1_{\{C=c\}}\sin(2\theta_c)
 \label{eq:final-classical-formula}
\end{equation}
is centered, satisfies $|m^{\mathrm{col}}|\le1$, vanishes on helper
roots, and obeys
\begin{equation}
 \cE(m^{\mathrm{col}})\ge P_*-\varepsilon.
 \label{eq:final-classical-value}
\end{equation}
After the uncolored finite formula has been fixed, the total helper
probability may be taken arbitrarily small, with all individual
helper probabilities positive.  All depths, coefficients, label laws,
and functions in the conclusion are fixed independently of degree,
system size, and disorder.
\end{proposition}

\begin{proof}
Use \eqref{eq:scalar-near-optimal-smooth} with error
$\varepsilon/3$ and write $m^*=z$.  Apply Lemma~\ref{lem:classical-trigonometric-density}
to obtain an angle $\theta$ whose output $m=\sin(2\theta)$ has
$\norm{m-m^*}<\varepsilon/12$.  Both outputs are bounded by one,
so \eqref{eq:energy-l2-continuity} loses less than
$\varepsilon/6$.  Fix this entire bounded trigonometric formula,
including its finite depth and Lipschitz constants.

Reserve a distinct helper type for each propagation depth.
Choose their positive probabilities with total
$p_H<\eta_H$ sufficiently small, and partition the remaining
probability into finitely many colors of maximum probability
$\alpha$ sufficiently small.  Apply
Lemma~\ref{lem:classical-color-exclusion} to the angle calculation
and use the $2$-Lipschitz property of sine.  Taking $\alpha,p_H$
smaller if necessary gives
$\norm{m^{\mathrm{col}}-m}<\varepsilon/6$.
Another application of \eqref{eq:energy-l2-continuity} loses less
than $\varepsilon/3$.  The total loss is less than $\varepsilon$.
Oddness gives centering, and the remaining claims follow directly
from the construction. Every choice has been made for a fixed finite
calculation in the limiting rooted space; no uniform rate in
the number of colors, circuit depth, or helper probabilities has
been used.
\end{proof}

\section{Compiling classical formulas with seeded controls}
\label{sec:classical-synthesis}

The preceding section supplies a bounded classical formula whose rooted
energy is close to the SK optimum.  We now implement any formula in that
class by quantum controls with independent Gaussian longitudinal seeds.
A \emph{seeded circuit} is a finite word in the cost, the uniform mixer,
and the gates
\[
 \exp\!\left(-\ii t\sum_i s_iZ_i\right),\qquad t\in\R,
\]
where the same independent standard Gaussian family \((s_i)\) is used
throughout the word and averaged in its value.  Here
\(s_i=s_{\mathrm{raw},i}\) is the raw seed of
\eqref{eq:raw-seed-reencoding}: every colored target formula is evaluated
at \(C(s_i),S(s_i)\).  Both the color and the scalar Gaussian input
are therefore functions of the seed driving this single longitudinal
control.  All angles are deterministic.
Section~\ref{sec:seed-simulation} will remove this extra control after
the present construction has selected one finite word.

We use the masks, aggregations \(A_{RS}\), and seed-only exclusion
formulas of Definition~\ref{def:exclusion-formulas}.  In particular,
every transmitted nonlinear function and every output angle is bounded,
although aggregate fields need not be bounded.  A mask may recur in
different branches of an expression tree, but not along a single active
branch.  We retain this distinction throughout the proof: formulas that
share descendants are not assumed independent.

All approximations in this section take place in the infinite-degree
rooted representation, inside a fixed surrounding circuit.  We first
construct elementary masked controls, then use unused helper spins to
implement one aggregation, and finally compile a finite expression tree
from its outermost operations inward.  The helper contraction is the
main step.  Its proof controls both the visible aggregation and the
changes it induces inside every lower formula.

\subsection{The synthesis statement}

Fix a finite collection of such formulas \(\theta_R\), odd and
seed-only, one for each output mask \(R\).  Their ideal operation is
\begin{equation}
 U_{\mathrm{cl}}=\prod_i \exp(\ii\theta_iY_i),
 \qquad m_i=\sin(2\theta_i).
 \label{eq:ideal-classical-rotations}
\end{equation}
The product is locally defined: an output at a vertex uses only a
fixed-radius neighborhood.  The factors commute because the angles
are seed-only.  Reserve disjoint helper masks, absent from all these
formulas, and set the output angle to zero on them.  One helper mask
per aggregation depth is sufficient.  Let \(H\) denote their union
and let \(p_H\) be its probability.

\begin{theorem}[Classical synthesis]
\label{thm:classical-synthesis}
For every finite collection of bounded smooth odd seed-only exclusion
formulas of Definition~\ref{def:exclusion-formulas}, with disjoint
positive-probability helper masks, one for each positive aggregation
depth, absent from all their dependencies and with output angle zero on
their union \(H\), and every \(\eta>0\), there is a finite seeded
circuit whose insertion vector \(v\) satisfies
\begin{equation}
 \norm{\1_{H^c}(v-v_{\mathrm{cl}})}<\eta,
 \label{eq:synthesis-insertion}
\end{equation}
where \(v_{\mathrm{cl}}\) is the insertion vector of
\eqref{eq:ideal-classical-rotations}.  In particular its value obeys
\begin{equation}
 \abs{\cE(v)-\cE(m)}\le 2\eta+2\sqrt{p_H}.
 \label{eq:synthesis-energy}
\end{equation}
The circuit and all its angles are chosen independently of the degree,
system size, and disorder.  Helper probabilities can be arbitrarily
small but are fixed and positive when the construction is performed.
\end{theorem}

The proof will establish convergence in a fixed surrounding circuit,
not only convergence on a freshly initialized input.  This is needed
because a formula is called repeatedly inside commutator products.

We use the smooth spaces of Section~\ref{sec:foundations}.  At the
root let
\[
 W=1+\dd\Gamma(\Lambda),\qquad
 \norm{T}_{a\leftarrow b}=\norm{W^aTW^{-b}},\qquad
 \norm{d}_{a}=\norm{\Lambda^a d}.
\]
The bounded-seed version of the weights is used until the last
compilation step.  Thus the weights do not differentiate or multiply
the root seed.  In particular they commute with root masks on the odd
subspace.  A family of operators is \emph{uniformly smooth} if for
each \(a\) there are \(b,C\), independent of the member of the
family, with \(\norm{T}_{a\leftarrow b}\le C\); the same convention
applies to adjoints and to parameter derivatives.  For a parameter
\(y\), a polynomial factor in \(1+|y|\) is allowed.  A statement
\(T=O_{\mathrm{sm}}(\varepsilon)\) means that these bounds hold
with \(C\varepsilon\) for every fixed output power and every fixed
jet.  Conditional versions mean the essential supremum in the root
seed.  The constants may depend on the fixed formulas, time intervals,
and positive helper probabilities.

The conditional bounds in Theorem~\ref{thm:gate-calculus} imply the
following useful rule.  If \(T\vac\) has bounded conditional smooth
norms and \(R\) is a root mask, then
\begin{equation}
 \norm{W^a\1_R T\vac}\le C_a\sqrt{p_R}.
 \label{eq:synthesis-masked-column-bound}
\end{equation}
For odd columns, the removal of the vacuum in the definition of
\(\Lambda\) does not change this estimate or the orthogonality of
columns supported on disjoint root masks.

\subsection{Elementary masked controls and matrix products}
\label{sec:elementary-masked-controls}

We first identify the elementary gates from which the helper
operation will be built.  These gates need only the original Gaussian
longitudinal control, parity-preserving local rotations, and the cost.

We first treat the following gates as exact:
\begin{enumerate}
\item local \(X\) rotations with bounded even seed coefficients,
      and local \(Y\) or \(Z\) rotations with bounded odd seed
      coefficients;
\item interactions between two disjoint masks in any pair of odd axes
      in the \(YZ\) plane.
\end{enumerate}
We justify at the end of this subsection that they lie in the closure
of finite seeded circuits.  Working with exact elementary masked gates
in the interim keeps unused helpers exactly pinned.

We will repeatedly use pointwise matrix sums and commutators.  Here
is the topology of that use.  If \(F(x),G(x)\) are Hermitian
two-by-two matrices with smooth polynomially bounded derivatives,
the usual product and group-commutator formulas converge on compact
sets with every fixed number of derivatives.  Their derivative bounds
are polynomial in \(x\), uniformly in the subdivision.  To verify the
last, write the group-commutator loop as
\[
 V_h(x)=e^{\ii hF(x)}e^{\ii hG(x)}
                e^{-\ii hF(x)}e^{-\ii hG(x)}.
\]
Twice integrating its mixed derivative in the two loop times shows
\[
 \norm{\partial_x^\alpha(V_h-1)}
     \le C_\alpha h^2(1+|x|)^{r_\alpha},\qquad |h|\le1.
\]
For a differentiated product of \(N\) such loops, each differentiated
factor contributes \(O(h^2)\); at a fixed derivative order the sum
is a fixed polynomial in \(Nh^2\).  Undifferentiated factors are
unitary.  Thus \(h=N^{-1/2}\) gives the claimed uniform bounds.
The Taylor error of a loop is \(O(h^3)\), with all fixed jets on a
compact set, so telescoping \(N\) factors proves convergence there.
The sum formula is identical with \(Nh\) in place of \(Nh^2\).
Theorem~\ref{thm:gate-calculus} therefore permits these limits inside
any fixed surrounding circuit.

When such a matrix formula is used simultaneously at many vertices,
the coefficient algebra must be common and commuting, and exclude
the pivot mask.  Under that condition the entire finite-tree gate is
a product of exactly those two-by-two formulas.  Its reinterpretation
at a new pivot is permitted only after an exact finite-tree identity
has been established.  These identities are also valid on every cut
tree, so the rooted maps agree.

To obtain the local elementary controls from the raw seeded controls,
write \(s=s_{\mathrm{raw}}\) and conjugate \(X\) by
\(e^{\ii t sZ}\).  Every local coefficient in the colored formulas,
including every mask, is a measurable function of this same \(s\)
by \eqref{eq:raw-seed-reencoding}.  Taking sums and differences
gives \(\cos(2ts)X\) and \(\sin(2ts)Y\); a constant \(X\)
rotation changes \(Y\) to \(Z\).  Finite even and odd trigonometric
sums are thus available, together with their exponentials.  Bounded
measurable even or odd seed functions are limits in Gaussian measure
of uniformly bounded continuous functions with the same parity.
After taking a subsequence, convergence is almost everywhere.  A box
cutoff, periodic extension, and Fejer approximation give trigonometric
approximants with the same uniform bound and parity.  Local seed
derivatives are not part of the smooth topology, so bounded dominated
convergence and Theorem~\ref{thm:gate-calculus} apply in a fixed
surrounding circuit.  If intermediate expressions use the raw linear
seed, its polynomial moments are controlled by the seed-weighted
version of that theorem.

We isolate interactions by sign averaging, a standard refocusing
device for Ising couplings; see Jones and
Knill~\cite[Section~II]{JonesKnill1999}.  Refine the masks to a finite
partition \(R_1,\ldots,R_M\).  For a sign vector
\(\epsilon\in\{-1,1\}^M\), masked \(X\) flips conjugate the
ordinary cost to
\[
 C_\epsilon=D^{-1/2}\sum_{\{i,j\}}
       \epsilon_{R(i)}\epsilon_{R(j)}Z_iZ_j.
\]
For distinct \(a,b\), sign orthogonality gives the finite-tree
identity
\[
 C_{R_aR_b}=2^{-M}\sum_\epsilon\epsilon_a\epsilon_b C_\epsilon.
\]
All the terms on the right commute, so their exponential implements
the desired subcost exactly when the masked flips are treated as
exact.  Separate masked \(X\) rotations put its two endpoints in any
specified axes in the \(YZ\) plane.  This proves the claimed
elementary compilation without introducing interactions within a
mask.

\subsection{A contraction through unused spins}

The elementary local rotations constructed so far depend only on the
seed at the rotated vertex.  The next construction extends them to a
neighbor sum of a fixed lower formula.  It uses a disjoint set of
spins kept in their \(X=1\) state, and its error tends to zero when the sender
mask is divided into sufficiently small parts.

Let \(I,J,K\) be disjoint masks, with \(K\) an unused helper of
probability \(p_K>0\).  Let \(f_j\), on \(J\), be a bounded odd
seed-only exclusion formula whose cross-edge dependencies avoid
\(I,J,K\).  Other reserved helpers are avoided as well.  Temporarily
allow exact rotations by \(f\) on \(J\), including their restrictions
to further local seed masks.  These ideal calls will remain in the
finite approximants constructed here.  At the end of the section we
will remove them by compiling the full formula from its outermost
operations inward.  We shall obtain
\begin{equation}
 H_{I,J,f}=\sum_{i\in I}P_i(A_{IJ}f)_i,
 \qquad P\in\{Y,Z\}.
 \label{eq:desired-propagation}
\end{equation}

Partition \(J\) into finitely many even seed masks \(J_l\), with
probabilities \(p_l\).  The formula \(f\) is kept fixed when this
partition is refined.  The commutators below sum paths of the form
\begin{equation}
 \begin{gathered}
 \begin{array}{ccccccc}
 i\in I&-&j\in J_l&-&k\in K&-&j'\in J_l\\
 &&&&&&\downarrow\\[-2pt]
 &&&&&&f_{j'}
 \end{array}\\
 j'\ne j:\ \text{simple path},\qquad
 j'=j:\ i-j-k-j\ \text{(backtrack)}.
 \end{gathered}
 \label{eq:helper-path-schematic}
\end{equation}
The downward arrow denotes the lower calculation of \(f\), whose
cross-edge dependencies avoid all three displayed masks.  In the
backtrack, \(f\) is attached at \(j'=j\); the two sender spin
factors cancel and the pinned helper supplies its polarization.
The simple paths contribute an error that
disappears from nonhelper insertions as the partition is refined.
We must check their effects at every root cut, including cuts
inside the lower formula \(f\).

Write \(A_{ij}=D^{-1/2}\1_{i\sim j}\) and
\(b_j=\sum_{i\in I}A_{ji}P_i\).  Pointwise commutators on \(J_l\)
give
\begin{align}
 H_l&=\sum_{i\in I,j\in J_l,k\in K}
       A_{ij}A_{jk}P_iX_jZ_k,\\
 G_l&=\sum_{j\in J_l,k\in K}A_{jk}X_jf_jY_k.
 \label{eq:two-helper-generators}
\end{align}
For example, \((2\ii)^{-1}[\sum A_{ij}P_iY_j,
\sum A_{jk}Z_jZ_k]=H_l\), and
\((2\ii)^{-1}[\sum f_jY_j,\sum A_{jk}Z_jY_k]=G_l\).
These are exact identities: the only noncommuting factors meet on
\(J_l\), and the seed-only coefficients commute with every spin.

At the next pivot, \(K\), the two coefficients are
\[
 B_{k,l}=\sum_{j\in J_l}A_{kj}X_jb_j,
 \qquad F_{k,l}=\sum_{j\in J_l}A_{kj}X_jf_j.
\]
They commute with each other and with all spins on \(K\).  Their
calculation paths avoid \(K\), as required by the gate calculus.
Consequently \(-(2\ii)^{-1}[H_l,G_l]
=\sum_{k\in K}X_k B_{k,l}F_{k,l}\).  Summing and rescaling gives
\begin{equation}
 \widetilde H_{\mathcal P}
  =\frac1{p_K}\sum_l\sum_{
       i\in I,\,j,j'\in J_l,\,k\in K}
       P_i A_{ij}X_j A_{jk}X_k A_{kj'}X_{j'}f_{j'}.
 \label{eq:helper-contraction-generator}
\end{equation}
All its summands commute.  At every fixed partition
\(\mathcal P\), its exponential is therefore in the closure of
finite products of elementary masked gates and calls to \(f\).
This statement holds on the full spin space and in every fixed
surrounding circuit; it has not yet used helper pinning.

We now specify the space on which pinning is used.  On a finite
light cone with vertex set \(V\), define the orthogonal projection
\begin{equation}
 \Pi_{K,V}=\prod_{v\in V}
       \left[1-\1_K(s_v)\frac{1-X_v}{2}\right].
 \label{eq:helper-pinning-projection}
\end{equation}
Its range fixes every helper-labeled spin in \(\ket+\), at the root
and throughout its descendants, while retaining all seed coordinates.
For the degree limit, let
\[
 \begin{gathered}
 \mathsf k_K=\operatorname{Ran}\!\left(1-\1_K(s)\frac{1-X}{2}\right),\\
 \cH_{K,0}=\mathsf k_K,\qquad
 \cH_{K,h+1}=\mathsf k_K\otimes\Gamma_s(\cH_{K,h}^{\circ}),
 \end{gathered}
\]
where the centered subspace is taken relative to the same vacuum
\(\vac_h\) as before.  These spaces are subspaces of \(\cH_h\).
The occupation and empty-descendant embeddings restrict to them,
as do the bounded-seed weights, so their smooth norms are inherited
from the original recursion.  We call this the \emph{compressed space}.
Helper-labeled descendant spaces remain present; only their helper
spins are pinned, and the seed law is unchanged.

The ideal propagations, the exact generators
\eqref{eq:helper-contraction-generator}, and surrounding gates avoiding
\(K\) preserve this range.  The observable \(\1_{K^c}(s_o)Z_o\)
at the root \(o\) also preserves it.  Thus nonhelper insertions may
be computed in the compressed space.  This restriction is different
from multiplication of an insertion vector by \(\1_{K^c}\), which
selects only the root label and imposes no condition on descendant
spins.  The individual commutator factors need not preserve pinning;
their finite products will be chosen after the exact contraction.

\begin{proposition}[Simultaneous helper contraction]
\label{prop:helper-contraction}
Consider a fixed finite sequence in the parity-preserving,
root-seed-commuting gate-calculus class, containing finitely many
propagations \eqref{eq:desired-propagation}, possibly with different
\(I,J,f,P\), and otherwise fixed gates.  All its gates and classical
parameters avoid a common helper \(K\).  Start with its spins pinned
at \(X=1\).  Replace every indicated propagation by
\eqref{eq:helper-contraction-generator}, and refine their sender
partitions so that their largest part probability tends to zero.
Then all insertion vectors off \(K\) converge to those of the
original sequence.  The root maps and their adjoints have uniform
conditional smooth bounds on the compressed space, and converge
strongly there up to scalar phases depending only on the root label.
The same estimates hold on cut subtrees.  Responses through the fixed
lower formulas for \(f\) have the uniform jet estimates proved below;
no convergence of arbitrary boundary jets in the growing list of
partition-dependent top channels is asserted.
\end{proposition}

The contraction must control every root cut, including cuts inside
\(f\), uniformly as the sender partition is refined.  The next two
estimates provide the analytic bounds for its proof.  Most cuts use a
fixed number of fields.  At a \(K\)-root, however, each part
\(J_l\) supplies a pair of incoming fields, and the proof will
subtract their same-neighbor covariance.  The second estimate controls
this growing sum by the orthogonality of the masks.

\subsection{Estimates and proof of the contraction}

\begin{lemma}[A fixed number of commuting fields]
\label{lem:fixed-field-exponentials}
Let \(Q_j=q(d_j)\), \(1\le j\le k\), commute, where \(k\) is
fixed and \(\norm{d_j}_a\le C_a\) for every \(a\).  Let
\(F_s:\R^k\to\R\) have polynomial bounds on every derivative,
uniformly in the local seed \(s\).  For bounded \(t\), the operators
\(\exp(\ii tF_s(Q))\), their adjoints, and their time derivatives
are uniformly smooth.  Auxiliary parameter jets obey the same conclusion
when the weighted column jets and the joint derivatives of \(F\)
have the corresponding uniform polynomial bounds.  The constants
require no lower bound on the
Gram matrix of the \(d_j\)'s.  The same assertion holds with a fixed
root control axis, by conditioning on its two eigenvalues.
\end{lemma}

\begin{proof}
The Weyl estimates in Lemma~\ref{lem:weyl-bounds}, applied to
\(\sum_j\xi_jd_j\), give, for every \(a\),
\[
 \norm{W^a e^{\ii\xi\cdot Q}\zeta}
 \le C_a(1+|\xi|)^{r_a}\norm{W^{b_a}\zeta}.
\]
Only upper bounds for weighted column norms enter this inequality.
For an integer \(N\), put
\(g_{s,t}(x)=e^{\ii tF_s(x)}(1+|x|^2)^{-N}\).
Fix the desired output power.  First choose enough derivatives that
integration by parts gives an integrable Fourier bound with weight
\((1+|\xi|)^{r_a}\), and then choose \(N\) large enough that those
derivatives are integrable, uniformly in \(s,t\).  The polynomial
derivative assumptions permit this order of choices.  Functional
calculus gives, on smooth vectors,
\[
 e^{\ii tF_s(Q)}
  =\int_{\R^k}\widehat g_{s,t}(\xi)
        e^{\ii\xi\cdot Q}\,\dd\xi\,(1+|Q|^2)^N,
\]
with the chosen Fourier normalization absorbed in \(\widehat g\).
The polynomial on the right has a fixed smooth power bound by the
creation and annihilation estimates.  Fourier integration therefore
proves the assertion.  Differentiating in any additional parameters
also differentiates the fields and the polynomial factor.  If
\(v(y)=\sum_j\xi_jd_j(y)\), the central commutator relation gives
\begin{equation}
 \partial_y e^{\ii q(v(y))}
   =e^{\ii q(v(y))}
       \left(\ii q(\partial_yv)
                 +\ii\im\braket{v}{\partial_yv}\right).
 \label{eq:moving-column-weyl-derivative}
\end{equation}
Repeated derivatives are the same Weyl operator times finite ordered
products of derivative fields and scalar Gram pairings.  Each field
insertion has a smooth bound polynomial in \(|\xi|\), and every
scalar pairing is quadratic in \(\xi\), by the weighted column-jet
hypotheses.  The Leibniz rule gives the same finite-power bounds for
derivatives of \((1+|Q|^2)^N\); its differentiated fields need not
commute with the original list.  Increasing the Fourier decay and
input weight powers just chosen proves the jet assertions.  No
change of field coordinates, and hence no inverse Gram matrix,
occurs.
\end{proof}

\begin{lemma}[Quadratics with many orthogonal columns]
\label{lem:orthogonal-quadratics}
Let \(R_l\) be finitely many pairwise disjoint root masks, with
\(p_l=\Pr(R_l)\). Suppose \(d_l,e_l\) are supported on \(R_l\)
and are real-commuting, meaning that all fields \(q(d_l),q(e_l)\)
commute pairwise. Suppose also that for every \(a\)
\[
 \norm{d_l}_a+\norm{e_l}_a\le C_a\sqrt{p_l}.
\]
Set
\[
 Q=\sum_l\left(q(d_l)q(e_l)-\braket{d_l}{e_l}\right),
 \qquad \rho=\left(\sum_l p_l^2\right)^{1/2}.
\]
Then \(Q=O_{\mathrm{sm}}(\rho)\).  Its exponentials and adjoints
are uniformly smooth for bounded times, independently of the number
of columns.  They converge to the identity strongly on smooth vectors
as \(\rho\to0\), and in every fixed lower smooth norm on families
with all higher smooth norms bounded.
\end{lemma}

\begin{proof}
Normal ordering gives a pair-creation kernel, its adjoint, and a
one-particle kernel:
\[
 Q=a^*(\textstyle\sum_l d_l\otimes e_l)
    +a(\textstyle\sum_l e_l\otimes d_l)
    +\dd\Gamma\!\left(\textstyle\sum_l
          |d_l\rangle\langle e_l|+|e_l\rangle\langle d_l|\right),
\]
where the first kernel is understood symmetrically.  Orthogonality of
different masks gives
\[
 \left\|\sum_l\Lambda^a d_l\otimes\Lambda^b e_l\right\|_{\rm HS}
 +\left\|\sum_l
       |\Lambda^a d_l\rangle\langle\Lambda^b e_l|\right\|_{\rm HS}
 \le C_{a,b}\rho.
\]
The same inequalities hold with \(d,e\) interchanged.  Pair creation
by a Hilbert--Schmidt kernel has norm at most
\(\sqrt{(n+1)(n+2)}\) times its kernel norm on the \(n\)-particle
sector; the corresponding estimates for annihilation and a
one-particle kernel are bounded by a constant times \(n+1\).
Commuting powers of \(W\) through these operators inserts powers of
\(\Lambda\) on their kernel indices.  This proves the asserted
smooth smallness.

For completeness, this smallness of the generator alone is not used
to assert an exponential power bound.  Every iterated commutator
\(\ad_W^j(Q)\) is another quadratic with uniformly bounded
Hilbert--Schmidt kernels of the displayed form.  Consequently
\[
 \norm{\ad_W^j(Q)\zeta}\le C_j\norm{W\zeta}.
\]
Writing \([W^a,Q]\) with these commutators on the left yields
\(\norm{[W^a,Q]\zeta}\le C_a\norm{W^a\zeta}\) for positive
integer \(a\).  On smooth vectors, differentiation of
\(\norm{W^ae^{\ii tQ}\zeta}^2\) and Gronwall's inequality now give
a bound independent of the column count.  At a fixed finite column
count these differentiations are justified by
Lemma~\ref{lem:fixed-field-exponentials}; its constants are not used
in Gronwall's bound.  Equivalently one may first truncate the weights
and pass to the limit using these same inequalities.

Finally, spectral calculus gives
\[
 \norm{(e^{\ii tQ}-1)\zeta}\le |t|\norm{Q\zeta}
       \le C|t|\rho\norm{W^b\zeta}.
\]
The uniform higher-power bounds and interpolation prove the last
assertion.  The same conclusions hold for \(U^*QU\) when \(U\) is
unitary and both \(U\) and \(U^*\) have uniform smooth bounds.  They
are uniform for families of unitaries with common such bounds: indeed,
\(e^{\ii tU^*QU}=U^*e^{\ii tQ}U\), and the smooth power losses
compose.
\end{proof}

\begin{proof}[Proof of Proposition~\ref{prop:helper-contraction}]
We prove a single replacement with arbitrary preceding root maps
from the gate calculus: they preserve the joint parity, commute with
root seed multiplication, and have uniform conditional smooth bounds.
These hypotheses are preserved by all the gates used here.  For the
sender partition put
\[
 \rho=\left(\sum_l p_l^2\right)^{1/2},\qquad
 \varepsilon=\max_l\sqrt{p_l}.
\]
We will estimate the Hamiltonian correction at each root cut, prove
uniform smooth bounds for its exponential, and then iterate the
replacement.  At helper roots a scalar phase remains; we will check
that it cancels from the child signals before the next replacement.
The partition changes only the terminal coefficients of the lower
calculation: the formula \(f\) and its finite channel list stay fixed.

On the range of \eqref{eq:helper-pinning-projection}, and hence in its
recursive occupation realization, we may set \(X_k=1\) at every
\(k\in K\).  The preservation facts just established apply to the
exact replacement and all its surrounding gates.

\paragraph{The relative root operation.}
Fix the partition and introduce an unscaled real time \(t\), keeping
the preceding circuit fixed as \(t\) varies.  Write \(H_D\) for the desired Hamiltonian and \(\widetilde H_D\) for the
complete contraction Hamiltonian.  On a padded rooted tree \(T\),
retain target terms through a depth \(d\) larger than the influence
distance.  On each child component \(T_a\), use the same formula with
zero parent input and target depth \(d-1\).  All parameter formulas
are computed on the padded tree.  For either Hamiltonian \(F_D\), put
\[
 L_{F,D}=F_D(T)-\sum_a F_D(T_a).
\]
All terms in the two Hamiltonians and their branch baselines commute:
they are diagonal in the same physical site axes.  Factor the preceding
circuit as \(U_D=\mathcal U_Du_D\), where
\(\mathcal U_D=\bigotimes_aU_{a,D}\).  The new child circuits are
\(e^{\ii tF_D(T_a)}U_{a,D}\).  Factoring their product from
\(e^{\ii tF_D(T)}U_D\), as in
\eqref{eq:exclusion-root-update}, gives exactly
\[
 u_{F,D}(t)=\mathcal U_D^*e^{\ii tL_{F,D}}\mathcal U_Du_D.
\]
Consequently the relative root maps satisfy
\begin{equation}\label{eq:helper-relative-root-group}
 \mathcal T_{\mathcal P,D}(t)
 :=u_{H,D}(t)^*u_{\widetilde H,D}(t)
 =U_D^*e^{\ii t(L_{\widetilde H,D}-L_{H,D})}U_D.
\end{equation}
This is a unitary group in \(t\), whose exponent is a root cut
difference.  The zero-input branch Hamiltonians have been subtracted;
no limit of the extensive Hamiltonian difference is being asserted.
The identity holds on the compressed space as well.

At this fixed partition, Theorem~\ref{thm:gate-calculus} applies to
both complete operations with \(t\) as an external parameter.  The
complete contraction is the gate at the \(K\) pivot, with coefficient
\(p_K^{-1}\sum_l B_{k,l}F_{k,l}\).  Their local matrices have polynomial bounds for all mixed time and input
derivatives, locally uniformly in \(t\).  The root maps, their adjoints,
and their time derivatives therefore have smooth degree limits.
Passing to the limit in \eqref{eq:helper-relative-root-group} gives a
unitary group \(\mathcal T_{\mathcal P}(t)\): its group law and
inverses follow from the finite-degree identities, and time-derivative
convergence gives strong continuity.  Its self-adjoint generator
\(A_{\mathcal P}\) is defined by this group.  Differentiation at zero
shows that every smooth vector is in its domain and that, under the
occupation identifications,
\begin{equation}\label{eq:helper-total-generator-action}
 A_{\mathcal P}\zeta
 =\lim_{D\to\infty}
 U_D^*(L_{\widetilde H,D}-L_{H,D})U_D\zeta.
\end{equation}

We now compute this derivative by separating diagonal paths, simple
paths, and lower-formula responses.  This is an algebraic decomposition
of the finite-degree cut difference, not an application of the gate
calculus to each separated Hamiltonian.  All contributions to
\eqref{eq:helper-total-generator-action} use the preceding full frame.
Thus top-path expressions first computed in the child frame receive
one final conjugation by \(u\); the lower-response formulas already
display it.  The smooth bounds are preserved by this conjugation.
The diagonal average and its helper-root scalar are included before
any phase is removed.

\paragraph{The diagonal paths.}
The terms \(j=j'\) in
\eqref{eq:helper-contraction-generator} equal
\begin{equation}
 \sum_{i\in I,j\in J}A_{ij}P_i f_j
       \frac{\#\{k\sim j:k\in K\}}{Dp_K}.
 \label{eq:diagonal-helper-paths}
\end{equation}
Away from the incidence lost at a \(K\)-root cut, the last factor
tends to one, with all smooth bounds required in the gate calculus:
it is a normalized average of bounded local seed masks.  This remains
true after conjugation by preceding maps with uniform smooth bounds.
One can see the rate on the symmetric embedding by decomposing that
average into its vacuum scalar, vacuum row and column, and occupied
block.  The centered row and column are \(O(D^{-1/2})\), and the
occupied block is \(O(D^{-1})\) with a loss of one number power.

If the removed root lies in \(K\), the lost incidence contributes
\((Dp_K)^{-1}\sum_{j\sim o}\1_J(j)b_jf_j\).
The average has a scalar vacuum limit.  Thus it contributes only a
phase conditional on the root's membership in \(K\).  Below we also
recover this scalar as the same-neighbor contraction of the full
\(K\)-root expression.  No estimate discards an uncentered
normalized sum.

\paragraph{The nonbacktracking paths at top masks.}
Remove the terms \(j=j'\).  A remaining path
\(i,j,k,j'\) is simple on a tree.  At a terminal \(j'\in J_l\),
looking away from its parent, its extra coefficient is
\begin{equation}
 R_{j',l}=\frac1{p_K}
   \sum_{k\sim j',\,k\in K} A_{j'k}
   \sum_{j\sim k,\,j\in J_l,\,j\ne j'}
         A_{kj}X_jb_j.
 \label{eq:extra-terminal-coefficient}
\end{equation}
Here all inner fields are computed in descendant trees.  Let
\(d_l\) be the child vacuum column of \(\1_{J_l}Xb\), in its
preceding Heisenberg frame.  It is centered, since it is odd, and
\eqref{eq:synthesis-masked-column-bound} gives
\(\norm{d_l}_a\le C_a\sqrt{p_l}\).
The first normalized lift is \(q(d_l)\).  Masking its next root on
\(K\), conjugating by the preceding map, and lifting again gives
\begin{equation}
 R_{j',l}=O_{\mathrm{sm}}(\sqrt{p_l})
 \quad\hbox{conditionally at }j'\in J_l.
 \label{eq:terminal-smallness}
\end{equation}
The constants include \(p_K^{-1}\), which is fixed.  The reversed
coefficient, with \(f\) in place of \(b\), has the same bound.
This reasoning uses centered lifts at both steps; the parity of the
transmitted fields supplies their centering.

At a root in \(I\), the final incoming column is the combination
over \(J_l\) of the small coefficients just estimated.  Its squared
weighted norm is bounded by
\(C\sum_l p_l^2\), since conditioning on \(J_l\) contributes
one factor \(p_l\) and \eqref{eq:terminal-smallness} contributes
another.  The extra root Hamiltonian is its normalized lift times
the root axis \(P\), and is \(O_{\mathrm{sm}}(\rho)\).

At a root in \(J_l\), the two possible occurrences of that root,
as \(j\) or \(j'\), give a bounded smooth factor \(f\) or the
field \(b\) times one of the small coefficients in
\eqref{eq:terminal-smallness}.  The two factors use distinct
first-step masks.  Thus there is no further same-neighbor term to
subtract.  The root correction is
\(O_{\mathrm{sm}}(\sqrt{p_l})\), uniformly conditionally on the
root seed in this part.  It is a smooth function of a fixed number of
commuting fields with uniformly bounded weighted columns.

At a root in \(K\), let \(d_l,e_l\) be the incoming child columns
of \(\1_{J_l}Xb\) and \(\1_{J_l}Xf\).  All their fields commute.
The root contribution of the full top-path sum is
\(p_K^{-1}\sum_l q(d_l)q(e_l)\).  The terms using the same neighbor
are the limit of the average
\[
 \frac1{Dp_K}\sum_{j\sim o}\1_J(j)b_jf_j
       \longrightarrow
       \kappa=\frac1{p_K}\sum_l\braket{d_l}{e_l}.
\]
In particular, this is an average on one branch, not the product of
two independent lifts.  After subtracting it, the correction is
\begin{equation}
 Q_K=\frac1{p_K}\sum_l
       \left(q(d_l)q(e_l)-\braket{d_l}{e_l}\right).
 \label{eq:helper-root-quadratic}
\end{equation}
The weighted columns belonging to different parts are orthogonal.
Lemma~\ref{lem:orthogonal-quadratics} gives
\(Q_K=O_{\mathrm{sm}}(\rho)\) and uniform exponential bounds
independent of the number of parts.  The scalar \(\kappa\) is the
harmless local phase identified from the diagonal paths.

This list accounts for all top-path cuts.  If the root belongs to
\(I,J,K\), a top path not containing it cannot change after the cut:
the dependencies of \(f\) exclude all three masks.  At a \(J\)
root, its own \(f\) uses only descendants.  A root outside both the
top masks and the lower calculation has no correction.  It remains
to check cuts inside the lower computation of \(f\).

\paragraph{Responses through the lower formula.}
We have accounted for cuts through the four-vertex helper paths.  A cut
can also change the value of \(f\) at their terminal vertices.  The
following recursion keeps those changes, including mixed responses
between channels, in the estimate.

First work in the physical frame.  Write
\(H_{\rm off}(T)=\sum_{l,j'\in J_l}\omega_{j',l}f_{j'}^T\),
where \(\omega_{j',l}=X_{j'}R_{j',l}\) initially denotes its full
top-path coefficient.  If a cut root \(v\) has a lower mask, it
cannot lie on any of the top paths.  Each coefficient is consequently
unchanged and lies wholly in one cut component.  For the child
components \(T_a\), the exact finite-tree identity is
\begin{equation}
 H_{\rm off}(T)-\sum_aH_{\rm off}(T_a)
   =\sum_{a,l,j'\in T_a\cap J_l}\omega_{j',l}
          \bigl(f_{j'}^T-f_{j'}^{T_a}\bigr).
 \label{eq:physical-lower-cut-identity}
\end{equation}
Only terminals within the fixed influence radius contribute.  At an
active terminal \(j'\in J_l\), a boundary input into its lower
computation therefore changes the Hamiltonian by
\begin{equation}
 X_{j'}R_{j',l}\bigl(f_{j'}(y)-f_{j'}(0)\bigr).
 \label{eq:terminal-response}
\end{equation}
Whenever this difference is active, its parent lies in a lower mask,
outside \(I,J,K\).  Consequently every first step of
\(R_{j',l}\) goes into a descendant branch.  That coefficient has
no boundary response of its own on this active path.  If the parent
is in a top mask, the boundary input into \(f\) is absent, and
\eqref{eq:terminal-response} is zero.  This observation ensures that
using the downward expression
\eqref{eq:extra-terminal-coefficient} loses no term.

Here is the precise difference recursion.  At a root \(v\) of the
lower calculation let \(m(y)\) be its fixed finite list of outgoing
channel values.  Let \(\mathcal H_v^{[d]}(y)\) be the sum of the
weighted terminal terms with terminal distance at most \(d\),
evaluated with boundary input \(y\), and put
\(\Delta_v^{[d]}(y)=\mathcal H_v^{[d]}(y)-\mathcal H_v^{[d]}(0)\).
The coefficients are computed downward, as just justified; the
terminal formulas themselves use all required sites.  In the
physical frame the recursion is a local term
\eqref{eq:terminal-response}, when the root is terminal, plus
\begin{equation}
 \sum_{a\text{ child}}
 \left[
   \Delta_a^{[d-1]}(D^{-1/2}m(y))
   -\Delta_a^{[d-1]}(D^{-1/2}m(0))
 \right].
 \label{eq:additive-response-recursion}
\end{equation}
Only channels of the fixed \(f\)-calculation are included.  At the
actual output root, where there is no parent input, the root
Hamiltonian correction is obtained from the local term and the
child sum evaluated at \(D^{-1/2}m(0)\), without this subtraction.
The definition counts distance to the terminal vertex; the top
coefficient beyond that vertex has already been computed downward.
Taking \(d\) beyond the finite influence range gives exactly the
Hamiltonian difference across the cut.  Indeed the identical
recursion with \(\mathcal H_a\) first holds by partitioning
terminals among children, and replacing
\(\mathcal H_a(z)=\Delta_a(z)+\mathcal H_a(0)\) cancels the
constant terms.  This also explains the argument \(m(0)\) in the
second summand.

Now factor the preceding circuit as
\(U_v=(\bigotimes_aU_a)u_v\).  These preceding maps are fixed in
the boundary variable of the new Hamiltonian.  Simultaneous
conjugation of the physical identity gives
\begin{equation}
 \widehat\Delta_v(y)=u_v^*\left[T_v(y)+\sum_a
   \left\{\widehat\Delta_a(D^{-1/2}m(y))
             -\widehat\Delta_a(D^{-1/2}m(0))\right\}\right]u_v,
 \label{eq:conjugated-lower-cut-identity}
\end{equation}
where \(\widehat\Delta_a=U_a^*\Delta_aU_a\), and \(T_v(y)\) is
the local terminal term \eqref{eq:terminal-response}, zero at a
nonterminal root. The direct term and messages are interpreted in
the same child frames.  There is no
derivative of \(U_a\) or \(u_v\).  Although a preceding Heisenberg
image of the top coefficient may appear to have larger support, its
boundary independence follows from the physical identity before
this fixed conjugation.  Repeated lower masks in different expression
branches cause no additional terms: these are identities for the
whole nonlinear function of all incoming channels, so their mixed
responses are retained.

The required parity covariance is explicit: if \(\Pi\) is the joint
parity, then \(\Pi\omega\Pi=-\omega\) and
\(\Pi f(y)\Pi=-f(-y)\).  Thus
\(\Pi\{\omega[f(y)-f(0)]\}\Pi=\omega[f(-y)-f(0)]\), and
the recursion preserves this covariance.  We pass the algebraic
recursion \eqref{eq:additive-response-recursion} to the limit by the
centered normalized-sum calculation used in the proof of
Corollary~\ref{cor:additive-response}.  At a fixed partition, the
terminal coefficients and their jets are finite products and smooth
compositions of centered odd normalized sums, together with the
separately estimated bounded-mask averages normalized by \(D^{-1}\).
The normalized-sum bounds and the conditional bounds for preceding
maps give uniform degree bounds in every required higher smooth norm;
these constants may depend on the fixed partition.  They propagate
through the fixed-depth difference recursion.  Taylor expansion in the
boundary inputs therefore retains the centered first response and the
scalar second response, while its third-order remainder carries
\(D\cdot D^{-3/2}\).  Commuting substitution handles the outgoing
message arguments.  These are termwise limits in the derivative of
the complete relative family, not a unitary calculus for the separated
Hamiltonians.  If \(D_\alpha,D_{\alpha\beta}\) are the first two input derivatives
of a child response at zero, set
\[
 r_\alpha=D_\alpha\vac,\qquad
 c_{\alpha\beta}=\bra{\vac}D_{\alpha\beta}\ket{\vac}.
\]
Its contribution is
\begin{align}
 &\sum_\alpha\bigl(m_\alpha(y)-m_\alpha(0)\bigr)q(r_\alpha)
 \nonumber\\[-2pt]
 &\quad+\frac12\sum_{\alpha,\beta}
 \bigl(m_\alpha(y)m_\beta(y)
        -m_\alpha(0)m_\beta(0)\bigr)c_{\alpha\beta},
 \label{eq:limiting-additive-response}
\end{align}
followed by the preceding root conjugation.  At a nonterminal lower
root with no parent input, the Hamiltonian correction instead uses
the unsubtracted child contribution
\[
 u_v^*\left[
 \sum_\alpha m_\alpha(0)q(r_\alpha)
 +\frac12\sum_{\alpha,\beta}
       m_\alpha(0)m_\beta(0)c_{\alpha\beta}
 \right]u_v.
\]
This is the root correction specified after
\eqref{eq:additive-response-recursion}, rather than the value at
\(y=0\) of the boundary difference, which is zero.  Both expressions
use the same response columns and scalar second responses.
First derivatives are odd and hence centered.  The same calculation
applies to every fixed input jet.  The partition estimates below are
made after this degree limit.

At depth zero, \eqref{eq:terminal-smallness} and the derivative
bounds of \(f\) make every fixed response jet
\(O_{\mathrm{sm}}(\varepsilon)\), conditionally in the local
seed.  More explicitly, the sum of terminal jets over \(l\) is a
direct integral on the disjoint events \(J_l\), so its conditional
bound is the maximum of their \(C\sqrt{p_l}\) bounds, not their
sum.  Suppose the induction holds at depth \(d-1\).  All weighted norms of
\(r_\alpha\) and all \(|c_{\alpha\beta}|\) are then bounded by
\(C\varepsilon\).  The functions \(m\) and their derivatives have
fixed polynomial bounds.  Multiplication, differentiation, and the
preceding root conjugation in
\eqref{eq:limiting-additive-response} therefore give the same
smallness at depth \(d\), with a possible finite loss of powers.
This induction has fixed finite depth and a fixed number of channels,
both independent of the number of parts \(J_l\).  It requires no
independence between different lower channels.  In particular it
applies when two expression branches reuse some lower masks.  At a
nonterminal lower root, if there are \(k\) outgoing channels and \(s\) ordinary Gaussian
inputs used to compute \(m\), the root correction uses at most
\(s+k\) field directions.  Each \(r_\alpha\) is a single child
vector containing the whole lower calculation; its Hilbert-space
support need not be finite dimensional.  The scalar second
responses contribute no new directions.

The diagonal degree-factor errors have the same response recursion
with terminal coefficient
\[
 b_{j'}\left(\frac{\#\{k\sim j':k\in K\}}{Dp_K}-1\right)
\]
in place of \(X_{j'}R_{j',l}\).  On an active lower response path
the parent is not in \(K\) or \(I\), so this is again a downward
coefficient.  Its limiting smooth bounds vanish by the average
estimate already proved.  Thus all its lower responses vanish as
well.

\paragraph{Uniform topology and iteration.}
After the degree limit, the diagonal degree-factor errors vanish.
The preceding computations identify the total action in
\eqref{eq:helper-total-generator-action}.  At a helper root its scalar
part is \(\sigma_{\mathcal P}(s)=\kappa\1_K(s)\), with
\[
 |\kappa|\le p_K^{-1}\sum_l\|d_l\|\,\|e_l\|
             \le C p_J/p_K.
\]
This bounded root-seed multiplication commutes with every root map.
Hence
\[
 \mathcal S_{\mathcal P}(t)
 =e^{-\ii t\sigma_{\mathcal P}}\mathcal T_{\mathcal P}(t)
\]
is a strongly continuous unitary group.  Its self-adjoint generator
\(R_{\mathcal P}\) acts on smooth vectors by the total remaining
correction.  The four cut cases, in the common preceding frame, give
\[
 \begin{array}{c|l}
 \text{cut-root mask}&\text{remaining correction}\\ \hline
 I&O_{\mathrm{sm}}(\rho)\quad\text{(combined column)}\\
 J_l&O_{\mathrm{sm}}(\sqrt{p_l})\quad\text{(conditional on }J_l\text{)}\\
 K&u^*Q_Ku=O_{\mathrm{sm}}(\rho)\quad\text{(after its scalar phase)}\\
 \text{lower mask}&O_{\mathrm{sm}}(\varepsilon)
                  \quad\text{(fixed lower channels)}
 \end{array}
\]
Since the parts partition the fixed sender mask,
\[
 \rho^2=\sum_l p_l^2\le p_J\max_l p_l,
 \qquad \varepsilon^2=\max_l p_l.
\]
Both quantities tend to zero under the refinement hypothesis, with
\(f\) and \(p_K>0\) fixed.

To obtain uniform smooth bounds for this actual group, consider the
full field expression for the remaining correction at each root.  At
an \(I\)-root it uses a bounded number of combined columns.  At a
\(J_l\)-root include the fields for both possible sender occurrences.
At a lower root use all ordinary channel fields and response columns
in the unsubtracted root formula.  These current-generator fields
commute: before preceding conjugation they belong to the same physical
seed and fixed-axis algebra.  Thus the total expression is a smooth
function of a fixed number of commuting fields and, when present, the
fixed root axis, with the uniform polynomial derivative bounds already
proved.  At a \(K\)-root it is the complete centered quadratic
\(Q_K\).  Define \(V_{\mathcal P}(t)\) by the joint spectral
calculus of these expressions, as a direct integral over the root seed,
with the final conjugation by \(u\) wherever it is not already
displayed.  Lemma~\ref{lem:fixed-field-exponentials} gives its uniform
conditional smooth bounds in the fixed-field cases;
Lemma~\ref{lem:orthogonal-quadratics} gives the count-uniform bound at
\(K\).  These bounds concern the total expression at each root.

We verify that \(V_{\mathcal P}\) is the actual relative group
\(\mathcal S_{\mathcal P}\).  At fixed partition the field-calculus
group preserves the smooth domain and has smooth time derivatives
there, with
\(\partial_tV_{\mathcal P}(t)\zeta
 =\ii R_{\mathcal P}V_{\mathcal P}(t)\zeta\)
by the action identified in \eqref{eq:helper-total-generator-action}.
For a smooth vector
\(\zeta\), therefore, \(V_{\mathcal P}(s)\zeta\) remains in the
domain of \(R_{\mathcal P}\), and its image under that generator
varies continuously by the smooth bounds.  Differentiating gives
\[
 \frac{d}{ds}\left[
 \mathcal S_{\mathcal P}(t-s)V_{\mathcal P}(s)\zeta
 \right]=0.
\]
Evaluation at \(s=0,t\) proves
\(\mathcal S_{\mathcal P}(t)\zeta=V_{\mathcal P}(t)\zeta\), and
density extends the identity to the Hilbert space.  Thus the uniform
smooth bounds just proved apply to the relative root group itself;
no separate closure of the formal sum of corrections is chosen.

Spectral calculus for its generator now yields
\[
 \|[\mathcal S_{\mathcal P}(t)-1]\zeta\|
 \le |t|\|R_{\mathcal P}\zeta\|\longrightarrow0.
\]
Uniform higher smooth bounds upgrade this to convergence in every
fixed lower graph norm; the same holds for adjoints.  The relative
factor multiplies the desired root map on the right, so
\[
 u_{\widetilde H,\mathcal P}(t)
 =u_H(t)\mathcal T_{\mathcal P}(t)
 =e^{\ii t\sigma_{\mathcal P}}u_H(t)\mathcal S_{\mathcal P}(t).
\]
This proves convergence of the actual root maps up to their root-seed
phase, with uniform conditional smooth bounds.  The same argument
applies on cut subtrees, and the fixed lower response jets retain the
estimates proved above.

A root-seed phase cancels in each descendant Heisenberg conjugation
before its vacuum column is lifted; helper-labeled descendant spaces
need not be removed.  Starting from the fixed initial preceding maps,
apply this single-replacement argument successively through the fixed
finite sequence.  The continuity part of
Theorem~\ref{thm:gate-calculus} applies at each step with the conditional
smooth bounds just established.  It gives insertion convergence off
\(K\) and proves the simultaneous contraction statement.
\end{proof}

\subsection{Pointwise functions and the order of compilation}

Proposition~\ref{prop:helper-contraction} implements an aggregation
whenever its lower formula is still available as an ideal call.  To
remove all ideal calls without losing unused helper spins, we combine
it with pointwise synthesis and process the formula from its outermost
aggregations inward.

Propagation and elementary local controls generate the required
pointwise nonlinearities.  Suppose rotations by odd fields
\(h_1,\ldots,h_k\) are available on the same target mask.  Same-axis
rotations give their real linear combinations.  A constant local
\(X\) rotation interchanges the two odd axes.  For an even local
coefficient \(a(s)\), conjugating \(a(s)X\) by a rotation with odd
angle \(h\), and taking the difference at opposite angles, gives
\(a(s)\sin(2h)Y\), up to a sign convention.  For an odd local
coefficient \(c(s)\), conjugating \(c(s)Y\) by the same \(Z\)
rotation and taking the sum gives \(c(s)\cos(2h)Y\).
Thus every finite trigonometric expression
\begin{equation}
 \sum_\nu a_\nu(s)\sin(\xi_\nu\cdot h)
       +c_\nu(s)\cos(\xi_\nu\cdot h),
 \qquad a_\nu\text{ even},\quad c_\nu\text{ odd},
 \label{eq:parity-trigonometric-synthesis}
\end{equation}
can be implemented as an odd-axis Hamiltonian by pointwise matrix
sums.  This covers joint parity in the local seed and aggregate
arguments; it does not require multiplying arbitrary even fields.

The same construction applies to bounded smooth odd functions of a
finite list of odd aggregate inputs.  To see the required closure
precisely, cut off on larger symmetric boxes, periodically extend,
and approximate the smooth extension by its Fourier polynomials.
The cutoff can be chosen even, and parity projection gives exactly
the coefficient parities in
\eqref{eq:parity-trigonometric-synthesis}.  Uniform bounds for
derivatives in the aggregate variables are preserved on each fixed
compact set.  Choose successive Fourier truncations accurate for
increasing derivative orders.  The cutoffs and periodic extensions
can be chosen with polynomial derivative bounds independent of the
box: inside the box these follow from the original bounds, and
outside its fundamental cell the absolute value of the reduced
coordinate is no larger than that of the original coordinate.
Fourier errors in the finitely many required derivative norms may
be made at most one.  A diagonal choice gives the polynomial
domination and local all-jet convergence used by
Theorem~\ref{thm:gate-calculus}.  Alternatively, if only the terminal
insertion is needed, first truncate the finitely many input laws and
use their smooth moment bounds.  Local seed coefficients remain
bounded measurable functions throughout.

We finish by specifying the nesting of all limits.  This is also
where helper pinning is used.  Define expression rank by
\[
 \operatorname{rank}(g(s))=0,\qquad
 \operatorname{rank}(F(h_1,\ldots,h_k))=\max_j\operatorname{rank}(h_j),
 \qquad
 \operatorname{rank}(A_{RS}f)=1+\operatorname{rank}(f).
\]
Local masking does not change rank.  Reserve a distinct helper
\(K_r\) for each positive rank.  Every original formula avoids
every helper.  Before processing rank \(r\), all unresolved calls
are classical formulas with this property; every already produced
elementary gate acts only on original masks and previously used
helpers \(K_q\), \(q>r\).  Thus the current surrounding word has
no target, spin control, or nonzero emitted seed channel on \(K_r\).
Its spins are exactly pinned, even if rotations on original output
masks have already occurred.  The estimates above allow arbitrary
preceding states in the gate-calculus class on those original masks.

\begin{enumerate}
\item At the largest remaining rank, replace each pointwise
nonlinearity by the pointwise synthesis just described.  Take these
limits successively in the fixed surrounding circuit.  After fixing
their finite products to any desired accuracy, the remaining calls
at rank \(r\) are rotations by explicit propagated odd inputs
\(A_{IJ}f\), where \(\operatorname{rank}(f)\le r-1\).

\item Apply Proposition~\ref{prop:helper-contraction} to all those
propagations simultaneously, using the reserved helper \(K_r\).
It is unused by every other exact gate at this stage and is
therefore pinned.  Choose finite sender subdivisions sufficiently
fine for the desired accuracy.  Sender subdivisions change only
the local calling mask; they do not refine the lower formulas or
their channel lists.

\item With that subdivision fixed, implement all the commutators
and sums defining \eqref{eq:helper-contraction-generator} by finite
pointwise products.  Their convergence holds on the full space, so
these finite products may temporarily rotate the helper away from
\(X=1\).  Pinning was used in the preceding contraction limit,
not during this inner approximation.

\item The resulting finite word contains only lower-rank formula
calls and elementary masked gates.  Repeat the same procedure at
the next rank, using its new helper.  The previously produced
finite word may act on higher helpers, but it avoids this new one
exactly.  Thus the pinning hypothesis holds anew.  Source formulas
retain all ancestor exclusions, including those introduced by their
particular calling masks. A sender subdivision only restricts the
local calling mask: exclusion of the original sender mask implies
exclusion of each submask. The commutators create further occurrences
of the same lower formula, without adding a cross-edge dependency or
introducing a control on any still-unused helper.
\end{enumerate}

This induction terminates because the formula rank is finite.
At each stage there are only finitely many occurrences in the
already fixed surrounding word.  Accuracies can therefore be chosen
successively from the outside inward.  No uniform estimate in the
number of commutator factors, the number of formula occurrences
created by an earlier stage, or inverse helper density is asserted
or needed.  The only uniformity in a growing sender subdivision is
the one explicitly proved in
Proposition~\ref{prop:helper-contraction}.

At termination, replace the elementary masked gates by their finite
seeded-circuit approximations from
Section~\ref{sec:elementary-masked-controls}.  All
mask-controlled operations are compiled at this final stage.
Theorem~\ref{thm:gate-calculus} permits the successive replacements
inside the fixed finite remaining word, including the raw Gaussian
seed control with its moment weight.  Exact helper pinning is no
longer required after the contraction limits have been fixed.
This constructs a finite seeded word with arbitrarily small
insertion error off the helpers.

\begin{proof}[Proof of Theorem~\ref{thm:classical-synthesis}]
The finite nested construction proves
\eqref{eq:synthesis-insertion}.  To identify the target value,
the ideal seed-only rotations give
\[
 v_{\mathrm{cl}}
  =m\ket{+}+\cos(2\theta)\ket{-},
\]
where both coefficients are classical functions of seeds; all
descendant spins remain in their pointed spin vectors.  In
\(\braket{v_{\mathrm{cl}}}{\cL v_{\mathrm{cl}}}\), the second
summand on the left is orthogonal to the range of \(\cL\).  Its
shifted version on the right has a child spin excitation and is
orthogonal to the first, purely classical summand on the left.
Consequently \(\cE(v_{\mathrm{cl}})=\cE(m)\).  The same identity
holds after omitting helper insertions, because \(m=0\) there.

For any circuit preserving root seed labels,
\(\norm{\1_Hv}=\sqrt{p_H}\): the mask commutes with the circuit
and with \(Z\), and \(Z\) is unitary.  Since \(\cL\) is an
isometry, Proposition~\ref{prop:energy} gives
\[
 \abs{\cE(v)-\cE(\1_{H^c}v)}\le2\sqrt{p_H}.
\]
Both projected insertion vectors have norm at most one, so
\eqref{eq:synthesis-insertion} changes their energy by at most
\(2\eta\).  This proves \eqref{eq:synthesis-energy}.  All
approximation limits above compare already-defined fixed-word tree
values.  After choosing finite approximants, the resulting seeded
word and all its parameters are independent of the degree, system
size, and disorder.
\end{proof}

\section{Removing the longitudinal seed controls}
\label{sec:seed-simulation}

Theorem~\ref{thm:classical-synthesis} produces a finite word that uses a Gaussian
longitudinal field in addition to the cost and mixer.  We now remove
that extra control while preserving its limiting tree energy.  This is
the last compilation step: the fixed-word transfer then applies to the
resulting ordinary QAOA word.

A \emph{seeded word} is a finite word in the cost, the mixer, and
\(\exp(-\ii\sigma\sum_i Z_i s_i)\), where the \(s_i\) are independent
standard real Gaussians.  The same Gaussian family is used at every
seed gate.  The Gaussian variables are averaged in evaluating the word;
equivalently, they are multiplication operators on independent pointed
Gaussian spaces.  All the angles, including \(\sigma\), are deterministic.
We fix this entire seeded word before choosing any of its approximants.
Its limiting tree energy is \(\cE(v_Z)\), where \(v_Z\) is its
root insertion vector. The seeded root update in
Lemma~\ref{lem:ordinary-updates} gives this vector and the uniform
occupation bound, so Proposition~\ref{prop:energy} identifies its
edge-energy limit with \(\cE(v_Z)\). This uses the seeded tree
calculus, not an extension of the ordinary-word SK transfer.

\begin{theorem}[Removal of seed controls]
\label{thm:seed-simulation}
For every finite seeded word and every \(\epsilon>0\), there is a finite
ordinary cost--mixer word whose limiting tree energy differs from that of
the seeded word by less than \(\epsilon\).  Its angles are independent of the
degree, the system size, and the disorder.  In particular, the approximation
uses the thermodynamic limit at a fixed word.
\end{theorem}

The proof uses two parameters, chosen in order.  We first choose a small
fixed \(b>0\).  With \(b\) fixed, \(r\downarrow0\) creates a component of
amplitude \(r\) with a Weyl displacement of order \(r^{-1}\).  Short cost pulses can
selectively rotate this displaced component while their first-order action
cancels on the component of order-one amplitude.  A cost echo of duration
\(r^{-1}\) converts the change of the small component into a finite
longitudinal field.  All approximation statements below are statements in
the already established fixed-word tree limit.

We first establish estimates for displaced Fock vectors and prepare
an exact two-component state.  Finite selective spin rotations then let
us propagate its small displaced component through the requested word.
The echo calculation produces the seed field; a separate root-shift
calculation identifies its energy.  Finally, a finite generic
perturbation supplies all the nonzero projections required to choose
the selective rotations.

\subsection{Uniform estimates for packets}

The large displacements will not have bounded occupation moments.
We therefore keep each displacement separate from a tail vector
whose moments are controlled.  Write \(N\) for the number operator in
the \emph{top} child Fock space and put
\[
 \mathcal S=\bigcap_{k\geq0}\operatorname{Dom}(1+N)^k,
 \qquad \norm{\eta}_k=\norm{(1+N)^k\eta}.
\]
The spin factor is included in this notation.  No weight is imposed inside
a child-mode direction.  A bounded family of inputs means a family
\(\eta_r\in\mathcal S\) such that
\(\sup_r\norm{\eta_r}_k<\infty\) for every \(k\).  Uniform errors below are
uniform on every such family.  More explicitly, every estimate used below
depends on finitely many of these seminorms; its constants may depend on
the fixed finite pulse list and on \(b\), but not on \(r\) or on the input.

Let
\[
 W(d)=e^{\ii q(d)},\qquad
 W_P(d)=e^{\ii Pq(d)},\qquad P^2=\1,
\]
where \(P\) is a root Pauli matrix.  The Weyl identities are
\begin{equation}
\label{eq:seed-weyl}
 W(d)W(e)=e^{-\ii\im\braket d e}W(d+e),
 \qquad
 W(d/r)^*q(v)W(d/r)
 =q(v)-\frac{2}{r}\im\braket v d.
\end{equation}
A packet with scaled center \(d_r\) is \(W(d_r/r)\eta_r\), with
\(\eta_r\) a bounded family.  A finite sum of packets is allowed, as are
bounded root spin matrices in their tails.  Scalar phases of modulus one
may depend on \(r\) and need not converge.

\begin{samepage}
\begin{lemma}[Packet estimates]
\label{lem:seed-packet-estimates}
The following estimates are uniform for bounded families of tails.
\begin{enumerate}
\item Displacements with bounded directions preserve all the seminorms
\(\norm\cdot_k\), with bounds depending only on the direction norm and
\(k\).  For \(\norm d\leq C\),
\[
 W(rd)\eta=\eta+\ii r q(d)\eta+O(r^2)\norm\eta_1
 \quad\hbox{in norm}.
\]
If \(d_r-e_r\to0\) in norm and both directions are bounded, then
\((W(d_r)-W(e_r))\eta_r\to0\) in norm.  These assertions also hold
for \(W_P\).
\item If \(\norm{d_r-e_r}\geq\kappa>0\), then for every integer \(M\),
\begin{equation}
\label{eq:seed-packet-separation}
 \abs{\braket{W(d_r/r)\eta_r}{W(e_r/r)\xi_r}}=O(r^M).
\end{equation}
The conclusion is unchanged by inserting a bounded root spin matrix.
\item For bounded \(d_r,v_r\) and fixed \(\tau\), conjugating the short
probe \(e^{-\ii r\tau Pq(v_r)}\) through a packet gives
\begin{equation}
\label{eq:seed-packet-probe}
 W(d_r/r)^*e^{-\ii r\tau Pq(v_r)}W(d_r/r)
 =e^{-\ii\tau a_rP}e^{-\ii r\tau Pq(v_r)},
 \qquad a_r=-2\im\braket{v_r}{d_r}.
\end{equation}
Consequently a fixed finite list of short probes preserves its scaled
centers and bounded centered-tail moments, without any smallness
assumption on \(a_r\).
\end{enumerate}
\end{lemma}
\end{samepage}

\begin{proof}
The first assertion follows from the creation--annihilation estimates in
Lemma~\ref{lem:weyl-bounds}, or by differentiating a Weyl displacement on
\(\mathcal S\).  In the Taylor remainder one uses
\(\norm{q(d)^2\eta}\leq C\norm d^2\norm\eta_1\).
The Weyl product identity then gives the assertion for two nearby
directions.  Resolving \(P\) into its two spectral projections proves the
spin-controlled statements.

For the second assertion, the Weyl product reduces the question to a
matrix coefficient of \(W(f_r/r)\), where \(\norm{f_r}\geq\kappa\).
Choose the unit quadrature
\(Q_r=q(\ii f_r/\norm{f_r})\).  Equation~\eqref{eq:seed-weyl}
translates \(Q_r\) by \(2\norm{f_r}/r\).  Split each vector using the
spectral projection of \(Q_r\) onto
\([-\norm{f_r}/(2r),\norm{f_r}/(2r)]\).  The two translated central
intervals are disjoint.  The complementary pieces have norm bounded by
\(C_Mr^M\norm\eta_{k_M}\), by the spectral theorem and the uniform
estimate \(\norm{Q_r^M\eta}\leq C_M\norm\eta_{k_M}\).
This proves~\eqref{eq:seed-packet-separation}.  Root spin matrices commute
with these projections.  Finally,~\eqref{eq:seed-packet-probe} is exactly
the second Weyl identity; its right-hand side is a spin rotation followed
by a displacement of size \(O(r)\).  Both preserve the required tail
bounds.  Notice that this last argument precedes, and does not use, any
claim that the probe acts almost trivially on an ordinary packet.
\end{proof}

We use exact local unitaries and exact insertion directions throughout.
At a given stage, if the local unitary is \(u_r\), put
\[
 v_{P,r}=u_r^*Pu_r\vac.
\]
All the \(v_{P,r}\) have norm one.  A toggled cost probe about an axis
\(P\) in the \(ZY\)-plane is
\(e^{-\ii r\tau Pq(v_{P,r})}\).  It is an ordinary word: conjugate a
cost pulse by two opposite mixer pulses.  The direction is the insertion
at the start of that probe; it is constant during the probe itself.

\subsection{Initialization and the exact two-component decomposition}

We next prepare the two components that the selective probes will distinguish.
Their decomposition is exact, even though they need not be orthogonal.
This exactness will keep the special component's amplitude equal to
\(r\) through all later unitary updates.

Fix once and for all
\(\beta_\star=\pi/4\), and put \(R_\star=e^{-\ii\beta_\star X}\).
This choice makes the two eventual external Gaussian directions
orthogonal.  For \(0<r<1\), set
\[
 \delta_r=\arcsin r,\qquad \gamma_r=b/r,\qquad e=Z\vac,
 \qquad
 \ell_r=Y e^{-2\ii\gamma_r Zq(e)}\vac.
\]
Apply, in order, cost \(\gamma_r\), mixer \(\delta_r\), cost
\(-\gamma_r\), and mixer \(\beta_\star\).  Immediately before the
negative cost the insertion is
\[
 v_{2,r}=\cos(2\delta_r)e+\sin(2\delta_r)\ell_r.
\]
Define \(d_{0,r}=\gamma_r(v_{2,r}-e)\).  The resulting local unitary
has the \emph{exact} decomposition
\begin{equation}
\label{eq:seed-initial-split}
 u_r=a_{0,r}U_{O,r}+a_{1,r}U_{S,r},\qquad
 U_{O,r}=R_\star W_Z(d_{0,r}),\qquad
 U_{S,r}=R_\star XW_Z(-2\gamma_re-d_{0,r}),
\end{equation}
where
\[
 a_{0,r}=\cos\delta_r\,e^{\ii\vartheta_r},\quad
 a_{1,r}=-\ii r e^{-\ii\vartheta_r},\quad
 \vartheta_r=\gamma_r^2\im\braket{v_{2,r}}e.
\]
Indeed, this is obtained by expanding \(e^{-\ii\delta_rX}\) and applying
the Weyl product identity.  Since
\[
 \braket e{\ell_r}=-\ii e^{-2\gamma_r^2},\qquad
 d_{0,r}=-2br\,e+2b\sqrt{1-r^2}\,\ell_r,
\]
we have \(\vartheta_r\to0\) faster than every power of \(r\), and
\(\norm{d_{0,r}-2b\ell_r}\to0\).

At every later stage retain the exact decomposition obtained by applying
all subsequent local updates to the two operators in
\eqref{eq:seed-initial-split}.  Denote its actions on an input \(\eta\)
by \(\Psi_{O,r}(\eta)\) and \(\Psi_{S,r}(\eta)\).  In particular,
\begin{equation}
\label{eq:seed-exact-special-norm}
 u_r\eta=\Psi_{O,r}(\eta)+\Psi_{S,r}(\eta),\qquad
 \norm{\Psi_{S,r}(\eta)}=r\norm\eta.
\end{equation}
No orthogonality is asserted or needed in this exact split.
The echo will depend on the change in the root \(Z\)-insertion when
the special component is flipped.  We therefore record the insertion
carried by that component, divided by its exact amplitude \(r\).
Let \(\psi_{S,r}=\Psi_{S,r}(\vac)\), and define at initialization
\begin{equation}
\label{eq:seed-address-direction}
 w_r=u_r^*Z\psi_{S,r}/r.
\end{equation}
Thus \(\norm{w_r}=1\) exactly.  The next lemma separates this
direction and the initialization direction \(\ell_r\) from all
bounded families of top-number-smooth vectors.

\begin{lemma}[The initial high directions]
\label{lem:seed-high-directions}
The vectors \(\ell_r,w_r\) are centered, have norm one, and satisfy
\[
 \braket{\ell_r}{w_r}\longrightarrow0,
 \qquad
 \norm{\1_{N\leq M}\ell_r}+
 \norm{\1_{N\leq M}w_r}\longrightarrow0
 \quad\hbox{for every fixed }M.
\]
They consequently have vanishing inner products with every bounded family
of top-number-smooth vectors.
\end{lemma}

\begin{proof}
Centering follows from the global spin-flip symmetry; every initialization
gate is even and the inserted \(Y\) or \(Z\) is odd.  The norm statements
follow from unitarity.  Because
\(R_\star^*ZR_\star=Y\) and \(YX=-\ii Z\), expansion of
\eqref{eq:seed-address-direction} gives
\[
 w_r=-Z W_Z(-2\gamma_re-2d_{0,r})\vac+o(1),
\]
where the omitted special--special term has norm \(r\), and the remaining
scalar phase tends to one.  The directions \(d_{0,r}\) are bounded.
Resolving the root spin shows that both displayed high vectors are sums
of coherent packets whose displacements have norm tending to infinity.
Their projections onto a fixed number sector therefore vanish, directly
from the coherent-state expansion.  At a fixed output \(Z\)-spin the
centers in \(\ell_r\) and in the displayed expression for \(w_r\) differ
by \(4\gamma_re+O(1)\); their overlap tends to zero by
Lemma~\ref{lem:seed-packet-estimates} (or the vacuum Weyl formula).
Finally, split a bounded family at \(N\leq M\).  Its complementary norm
is uniformly \(O(M^{-k})\), while its fixed-sector pairing tends to zero.
First take \(r\downarrow0\) and then \(M\to\infty\).
\end{proof}

The special part of~\eqref{eq:seed-initial-split}, on \emph{any} bounded
input family, has four packets.  Index them by the spin \(z=\pm1\)
before \(X\) and by the spin \(y=\pm1\) after \(R_\star\).  Their
scaled centers are \(-2bz e\), and their output spins are \(y\).
The bounded displacement \(-z d_{0,r}\) is part of the tail.  In
particular, the center list is independent of the input; only the tails
depend on it.

\subsection{A finite ensemble of spin rotations}

On the vacuum, the exact split has an ordinary component with bounded top-number
moments and a special component of norm \(r\) carried by four displaced
packets.  A short cost pulse acts on each displaced vector mainly as a
spin rotation.  We need a finite list of such pulses that gives the same
rotation at every special center, while its first-order action on an
undisplaced vector cancels.  The two zero sums in the next lemma impose
that cancellation separately for the two available spin axes.

The selective-rotation argument uses the Lie-bracket and polynomial-
approximation approach to ensemble control of Li and
Khaneja~\cite[Section~II, Examples~1--2]{LiKhaneja2005}. The finite-family
and exact zero-clock requirements needed here are proved below.

\begin{lemma}[Selective rotations with zero clocks]
\label{lem:seed-ensemble-control}
Let \(h_1,\ldots,h_m\) be a finite list of nonzero vectors in \(\R^2\).
There are two perpendicular axes \(P_1,P_2\) in the \(ZY\)-plane with
the following property.  Set \(c_{j,a}=h_j\cdot P_a\), using the
coordinates \((Z,Y)\).  Given any \(\theta\in\R\), there is a finite
list \((a_k,\tau_k)\), with \(a_k\in\{1,2\}\), such that
\begin{equation}
\label{eq:seed-zero-clocks}
 \sum_{k:a_k=a}\tau_k=0\quad(a=1,2),
 \qquad
 \prod_k e^{-\ii\tau_k c_{j,a_k}P_{a_k}}
 =e^{-\ii\theta X}\quad(1\leq j\leq m).
\end{equation}
The product is in chronological order, with later factors on the left.
The assertion is exact, not merely an assertion about closure.
\end{lemma}

\begin{proof}
Choose the axes so that every \(c_{j,a}\neq0\), and so that the pairs
\((c_{j,1}^2,c_{j,2}^2)\) are distinct unless the original vectors are
equal or opposite.  Such axes exist: each prohibited equality or zero is
a nontrivial trigonometric equation in the axis angle, apart from the
stated equal-or-opposite exception.  Delete duplicate square pairs; on
each deleted pair the product \(c_{j,1}c_{j,2}\) is the same.

In \(G=SU(2)^m\times\R^2\), take the two generators having spin
components \(-\ii c_{j,a}P_a\) and clock component the \(a\)-th
coordinate vector.  Their brackets have zero clocks.  A first bracket
has \(\ii X\)-coefficient equal, up to a common nonzero real constant,
to \(c_{j,1}c_{j,2}\).  Repeated double brackets with the first or
second generator multiply this coefficient by \(c_{j,1}^2\) or
\(c_{j,2}^2\), again up to nonzero common constants.  Thus the generated
Lie algebra contains the clock-zero vector with coefficients
\[
 c_{j,1}c_{j,2}F(c_{j,1}^2,c_{j,2}^2)\,\ii X
\]
for every real polynomial \(F\).  Polynomial interpolation on the finite
set of square pairs makes all these coefficients equal to one.

For completeness, exponentials of this Lie-algebra vector are reachable
by finite products of the two original one-parameter groups.  Let
\(\mathcal R\) denote the subgroup of such finite products, and let
\(\mathfrak g_0\) be the Lie algebra generated by the two generators.
The span of their conjugates by elements of \(\mathcal R\) is invariant
under each of the two adjoint one-parameter groups.  Differentiation
shows that it contains all iterated brackets and hence equals
\(\mathfrak g_0\).  Choose finitely many of these conjugates forming a
basis.  The product map of their one-parameter groups has surjective
differential at zero, so its image contains a neighborhood of the
identity in the connected analytic subgroup with Lie algebra
\(\mathfrak g_0\).  Each factor is itself a reachable conjugate.
Subdivide any one-parameter path into finitely many pieces in this
neighborhood.  The whole exponential is therefore in \(\mathcal R\).
Applying this to the common \(X\)-generator with zero clocks gives
\eqref{eq:seed-zero-clocks}.
\end{proof}

The preceding lemma treats fixed real coefficients.  In the circuit,
the coefficients come from actual insertion vectors and change during
the probes.  The next estimate controls that change without replacing
the insertion vectors by approximants before a singular rescaling.

\begin{lemma}[Accuracy of selective probes]
\label{lem:seed-selective-accuracy}
Consider a stage whose vacuum state \(u_r\vac\) has a finite packet
expansion, with norm-\(o(r)\) remainders and uniformly bounded scaled
centers.  Suppose the ordinary component has order-one amplitude and
the special component has amplitude \(O(r)\).  All designated tails have
bounded top-number moments.  Suppose also that one of the following
ordinary geometries holds.
\begin{enumerate}
\item At a standard stage all ordinary centers are zero.
\item At a huge-cost stage the ordinary part is
\(\psi_++\psi_-\), with \(Z\psi_y=y\psi_y\) and centers
\(-tyv_r\), where \(t\neq0\) and \(v_r=u_r^*Zu_r\vac\).
Every special center \(d_{j,r}\) obeys
\(\lim\im\braket{v_r}{d_{j,r}}\neq0\).
\end{enumerate}
Assume the limiting special projection pairs
\[
 h_j=\lim_{r\downarrow0}-2
 \bigl(\im\braket{v_{Z,r}}{d_{j,r}},
       \im\braket{v_{Y,r}}{d_{j,r}}\bigr)
\]
exist and are nonzero.  Use the axes and finite pulse list of
Lemma~\ref{lem:seed-ensemble-control} for these pairs, replacing its
spin generators by actual toggled probes of angles \(r\tau_k\).
Then the net probe action is the identity up to \(o(r)\) on every
ordinary packet, and is the prescribed \(X\)-rotation up to \(o(1)\)
on every normalized special packet.  The ordinary assertion holds also
after any root spin matrix has been inserted in its tail.  These
statements hold uniformly for bounded families of tails sharing the
stated center list, with the probes chosen solely from the vacuum data.
\end{lemma}

\begin{proof}
Let \(V_j\) be a prefix of the short probes, and write
\(v_{P,r}^{(j)}=(V_ju_r)^*P(V_ju_r)\vac\).  Center preservation and
all centered-tail bounds follow first from
Lemma~\ref{lem:seed-packet-estimates}, without estimating the spin
phases.  At a standard stage every ordinary probe has zero center phase.

At a huge-cost stage there is the exact identity
\begin{equation}
\label{eq:seed-hermitian-diagonal}
 \braket{v_{P,r}^{(j)}}{v_r}
 =\braket{PV_j\psi}{V_jZ\psi},\qquad \psi=u_r\vac.
\end{equation}
Group the ordinary packets by their original sign \(y\).  The diagonal
term of sign \(y\) on the right-hand side is
\(y\braket{PV_j\psi_y}{V_j\psi_y}\), which is real.  This remains true
although the probes may rotate the spin within that group.  The two
ordinary centers are separated by \(2|t|\), since \(\norm{v_r}=1\).
An ordinary center and a special center are also separated: their
imaginary projections against \(v_r\) differ by a quantity bounded away
from zero.  Lemma~\ref{lem:seed-packet-estimates} makes all these cross
terms smaller than any fixed power of \(r\).  Terms containing only
special packets are \(O(r^2)\), and the remainder terms are \(o(r)\).
Consequently, at \emph{every} intermediate probe prefix,
\begin{equation}
\label{eq:seed-huge-imaginary}
 \im\braket{v_{P,r}^{(j)}}{v_r}=o(r).
\end{equation}
More quantitatively, let \(\varepsilon_r=o(r)\) bound the norm of the
discarded vacuum remainder at the start of the probe sequence.  For
every fixed \(M\), the preceding argument gives
\begin{equation}
\label{eq:seed-huge-imaginary-modulus}
 \max_j\abs{\im\braket{v_{P,r}^{(j)}}{v_r}}
 \leq C_M\bigl(r^2+\varepsilon_r+r^M\bigr).
\end{equation}
The constant depends on the finite pulse list, a positive lower bound
on center separations, and finitely many designated vacuum-tail
seminorms.  The exact probe prefixes preserve the norm of the discarded
remainder; no number-operator estimate is imposed on it.
This derivation used only center preservation, so there is no circular
small-phase assumption.

By~\eqref{eq:seed-packet-probe}, an ordinary center \(-tyv_r\) therefore
has spin phase \(o(r)\) in every probe.  In either geometry a probe acts
as identity plus \(O(r)\) on the ordinary part and on that part with any
spin matrix inserted.  On the special part the norm is \(O(r)\), so
unitarity suffices.  Using
\[
 v_{P,r}^{(j+1)}-v_{P,r}^{(j)}
 =(V_ju_r)^*\bigl(V^*PV-P\bigr)(V_ju_r)\vac
\]
for the next probe \(V\), or equivalently estimating
\((PV-VP)(V_ju_r)\vac\), gives
\begin{equation}
\label{eq:seed-probe-direction-change}
 \norm{v_{P,r}^{(j)}-v_{P,r}^{(0)}}=O(r)
\end{equation}
for the fixed finite list.

After removing an ordinary center, expand the displacement in each
probe to first order.  Equation~\eqref{eq:seed-huge-imaginary} supplies
only an \(o(r)\) additional phase in the huge-cost geometry.  Replacing
\(v_{P,r}^{(j)}\) by its starting value in the first-order displacement
costs \(O(r^2)\), by~\eqref{eq:seed-probe-direction-change}.
The first-order sum is therefore
\[
 -\ii r\sum_{a=1}^2
 \left(\sum_{k:a_k=a}\tau_k\right)P_aq(v_{P_a,r}^{(0)})=0.
\]
Here is a uniform remainder bound for this cancellation.  For a fixed
ordinary center \(d_r\), let \(\mathcal B_r(d_r)\) be the net probe
action conjugated by \(W(d_r/r)\), and put
\[
 \Delta_r=\max_j\norm{v_{P_j,r}^{(j)}-v_{P_j,r}^{(0)}},\qquad
 \alpha_r(d_r)=\max_j
       2\abs{\tau_j\im\braket{v_{P_j,r}^{(j)}}{d_r}}.
\]
Then for every designated centered tail \(\eta\),
\begin{equation}
\label{eq:seed-uniform-probe-bound}
 \norm{(\mathcal B_r(d_r)-\1)\eta}
 \leq C_K\left[
       (r^2+r\Delta_r)\norm\eta_1
           +\alpha_r(d_r)\norm\eta\right].
\end{equation}
Indeed, remove each spin phase by its operator-norm bound and replace
each displacement direction by its starting value.  The latter
replacement costs at most \(C r\Delta_r\norm\eta_1\) through the
finite product.  In the product with starting directions, hold those
directions fixed at their given \(r\)-values and introduce an auxiliary
pulse strength \(\lambda\).  Its first derivative at \(\lambda=0\)
is zero, by the clocks.  Its second derivative is a
finite sum controlled by
\(\norm{q(v)q(w)\eta}\leq C\norm v\norm w\norm\eta_1\);
the intervening displacements preserve this seminorm uniformly.
Evaluating the Taylor remainder at \(\lambda=r\) gives
\(C_Kr^2\norm\eta_1\); no differentiability of the directions in
\(r\) is assumed.
The constant \(C_K\) depends only on the fixed pulse list; all actual
insertion directions have norm one.  At a standard center
\(\alpha_r=0\).  At a huge-cost center~\eqref{eq:seed-huge-imaginary-modulus}
gives \(\alpha_r\leq C(r^2+\varepsilon_r+r^M)=o(r)\).
Together with \(\Delta_r=O(r)\), this proves the claimed uniform
\(o(r)\) estimate using the single tail seminorm \(\norm\eta_1\).
It applies unchanged to a root-spin-inserted tail.  If the input also
contains a discarded norm remainder \(\rho_r\), its contribution to
the net identity error is at most \(2\norm{\rho_r}\), by unitarity.
No field operator is applied to such a remainder in this argument.

On a special packet the change of its phase coefficient during the
sequence is \(O(r)\), by~\eqref{eq:seed-probe-direction-change} and the
boundedness of its scaled center.  The small displacements contribute
\(O(r)\) on its tail, and its starting phase coefficients converge to
the prescribed \(c_{j,a}\).  The resulting finite spin product thus
converges to~\eqref{eq:seed-zero-clocks}.  Multiplication by the packet's
amplitude \(O(r)\) makes its error \(o(r)\).
All quantities defining the probes came from the vacuum.  The estimates
after removing a center, however, use only the tail seminorms.  They
therefore apply to every bounded input family with the same center list.
\end{proof}

\subsection{The packet invariant and the singular echo}

We now apply the selective estimates to the exact initialized
decomposition and propagate it through costs, mixers, and simulated
seed gates.  The required statements include
moving input vectors: later insertion estimates must be allowed to use
the approximating circuit itself as an input.

We call the times between simulated gates \emph{standard stages}.  A
second kind of stage occurs immediately after the forward cost in a
seed echo.  The induction will maintain the following assertions, with
the same center list for every bounded input family.
\begin{enumerate}
\item At a standard stage there is a unitary \(A_r\), a finite product
of root spin rotations and spin-controlled Weyl displacements with
uniformly bounded directions, and a scalar \(\zeta_r\) of modulus one,
such that
\begin{equation}
\label{eq:seed-low-ordinary}
 \Psi_{O,r}(\eta_r)=\zeta_r A_r\eta_r+o(r).
\end{equation}
The operators \(A_r,A_r^*\) preserve all top-number seminorms uniformly.
\item The special part has a finite packet expansion up to \(o(r)\),
with amplitude \(r\), bounded tails, scaled centers
\begin{equation}
\label{eq:seed-special-centers}
 d_{zy,r}=-2bz e+yD_r,\qquad z,y\in\{-1,1\},
\end{equation}
and output spin \(y\).  Coincident centers cause no problem.  The vector
\(D_r\) is a finite sum of actual earlier insertion directions with
fixed bounded coefficients.  Each such direction is within \(O(r)\)
of a bounded family in \(\mathcal S\).
\item At a standard stage the special insertion is transported:
\begin{equation}
\label{eq:seed-transport}
 u_r^*Z\psi_{S,r}/r=w_r+o(1),
\end{equation}
where \(w_r\) is the initial vector
\eqref{eq:seed-address-direction}.  Between a special flip and its
undoing the right-hand side instead has a minus sign.
\end{enumerate}
Here and below the norm remainders need not have bounded number moments.
They are kept as separate remainder vectors and carried through later
unitaries by their norm bounds.  Only the designated tails require
moment estimates.

The initialization satisfies this invariant with \(A_r=U_{O,r}\),
\(D_r=0\), and \(\zeta_r=e^{\ii\vartheta_r}\): its ordinary amplitude
differs from one by \(O(r^2)\).  Equation~\eqref{eq:seed-transport} is
initially exact.  The following observation will be used repeatedly.

\begin{lemma}[Insertion approximation from forward action]
\label{lem:seed-low-insertions}
If~\eqref{eq:seed-exact-special-norm} and
\eqref{eq:seed-low-ordinary} hold uniformly on bounded input families,
then
\begin{equation}
\label{eq:seed-insertion-O-r}
 \norm{u_r^*Pu_r\vac-A_r^*PA_r\vac}=O(r)
 \qquad(P=Y,Z).
\end{equation}
In particular the actual insertion directions have bounded smooth
approximants.  No approximation to \(u_r^*\) on a high packet is needed.
\end{lemma}

\begin{proof}
Put \(\widehat v_{P,r}=A_r^*PA_r\vac\).  This is a bounded smooth
family.  Unitarity and \(A_r\widehat v_{P,r}=PA_r\vac\) give
\[
 \norm{u_r^*Pu_r\vac-\widehat v_{P,r}}
 =\norm{Pu_r\vac-u_r\widehat v_{P,r}}
 \leq\norm{P(u_r-\zeta_rA_r)\vac}
    +\norm{(u_r-\zeta_rA_r)\widehat v_{P,r}}.
\]
Each term is \(O(r)\) by the exact special norm and the ordinary
approximation.  The possibly nonconvergent scalar phase cancels.
\end{proof}

For this induction, assume that the nonzero projection hypotheses of
Lemma~\ref{lem:seed-selective-accuracy} hold at every required probe
start.  We call these the \emph{visibility conditions}.  The final
subsection proves that one perturbation of the finite target word
ensures all of them and then constructs the probes in order.

An unscaled cost uses the actual insertion \(v_{Z,r}\) and simply
left-multiplies \(A_r\) by \(e^{-\ii cZq(v_{Z,r})}\).
On a special packet with spin \(y\), it adds a bounded displacement to
the tail and a scalar phase, leaving its scaled center and spin
unchanged.  It commutes with \(Z\), so it preserves the transport
identity exactly.

To implement a bulk mixer \(R_\theta=e^{-\ii\theta X}\), first use
selective probes to apply \(R_\theta^{-1}\) on the special packets,
then apply the global mixer \(R_\theta\).  On ordinary packets the
first step is identity to \(o(r)\), so the low update is \(R_\theta\).
On special packets the two spin actions cancel to \(o(1)\).  The same
is true on their \(Z\)-inserted tails.  Thus the special center and
spin return to their initial values, and transport is preserved to
\(o(1)\) after division by the special amplitude \(r\).

For a longitudinal seed gate use a fixed \(t\neq0\) and the four
operations
\begin{equation}
\label{eq:seed-echo-list}
 \text{cost }t/r,\qquad
 \text{selective }(-\ii X),\qquad
 \text{cost }-t/r,\qquad
 \text{selective }(\ii X).
\end{equation}
The next lemma verifies both the singular cancellation and its error
scale.

\begin{lemma}[One echo]
\label{lem:seed-one-echo}
Under the invariant and the visibility conditions, the echo
\eqref{eq:seed-echo-list} preserves the invariant.  Up to a scalar of
modulus one its low update is
\begin{equation}
\label{eq:seed-echo-low}
 A_r\longmapsto W_Z(g_r)A_r,
 \qquad
 g_r=\frac{t}{r}(v'_r-v_r)=-2t w_r+o(1),
\end{equation}
where \(v_r\) is the actual \(Z\)-insertion before the forward cost
and \(v'_r\) is the actual insertion immediately after the selective
flip.  The special-center update can be taken to be exactly
\begin{equation}
\label{eq:seed-D-update}
 D_r\longmapsto D_r-t(v_r+v'_r).
\end{equation}
All ordinary errors are \(o(r)\) before and after the large costs.
\end{lemma}

\begin{proof}
Write \(u_{\mathrm{in},r}\) for the exact prefix before the echo,
and \(V_r\) for its selective flip. The next two prefixes are
\[
 u_{\mathrm{cost},r}=e^{-\ii tZq(v_r)/r}u_{\mathrm{in},r},
 \qquad u_{\mathrm{flip},r}=V_r u_{\mathrm{cost},r}.
\]
Their actual \(Z\)-insertions satisfy
\[
 \begin{aligned}
 u_{\mathrm{in},r}^*Zu_{\mathrm{in},r}\vac
 &=u_{\mathrm{cost},r}^*Zu_{\mathrm{cost},r}\vac=v_r,\\
 u_{\mathrm{flip},r}^*Zu_{\mathrm{flip},r}\vac&=v'_r.
 \end{aligned}
\]
The first equality holds because the forward cost commutes with the
root \(Z\). It splits the ordinary
part into spins \(y\) with centers \(-tyv_r\), and changes each
special center to \(d_{zy,r}-tyv_r\).  Centered tails remain bounded.
The standard-stage imaginary projection against \(v_r\) is unchanged
by adding \(-tyv_r\).  Thus the assumed special visibility separates
special and ordinary centers at this huge-cost stage.

Denote the exact ordinary and special components after that forward
cost by \(\Psi^{\mathrm{cost}}_{O,r}\) and
\(\psi^{\mathrm{cost}}_{S,r}\), respectively.
Lemma~\ref{lem:seed-selective-accuracy}
gives net identity to \(o(r)\) on the ordinary part, including its
\(Z\)-inserted version.  On the special part it gives \(-\ii X\)
to error \(o(r)\), again including the \(Z\)-inserted version.
Since \(XZX=-Z\), it follows that
\begin{align}
 v'_r-v_r
 &=u_{\mathrm{cost},r}^*(V_r^*ZV_r-Z)
      (\Psi^{\mathrm{cost}}_{O,r}(\vac)+\psi^{\mathrm{cost}}_{S,r})\notag\\
 &=-2u_{\mathrm{cost},r}^*Z\psi^{\mathrm{cost}}_{S,r}+o(r)
 =-2r w_r+o(r).
\label{eq:seed-flip-change}
\end{align}
The last equality uses transport, which the forward cost preserves.
The same calculation shows that transport changes sign at this flip.

On an ordinary input the reverse cost combines with the forward cost
as
\begin{equation}
\label{eq:seed-exact-echo-cancellation}
 e^{\ii tZq(v'_r)/r}e^{-\ii tZq(v_r)/r}
 =\xi_r W_Z\!\left(\frac{t}{r}(v'_r-v_r)\right),
 \qquad |\xi_r|=1.
\end{equation}
The Weyl phase \(\xi_r\) is independent of the root spin because the
two spin factors multiply to \(Z^2=\1\).  It may fail to converge.
Equation~\eqref{eq:seed-flip-change} makes the residual direction
bounded and proves~\eqref{eq:seed-echo-low}.  The ordinary \(o(r)\)
probe error is carried through the reverse cost by unitarity.  No
insertion approximation has been divided by \(r\): the cancellation
in~\eqref{eq:seed-exact-echo-cancellation} used the two actual
directions, and only then was~\eqref{eq:seed-flip-change} invoked.

A special packet has spin \(-y\) between the flips.  Its two cost
displacements therefore add as \(-tyv_r-tyv'_r\), giving
\eqref{eq:seed-D-update}.  Here \(y\) remains its label from the
standard stage before the echo; it is not relabeled when the spin
flips.  Its scalar phases are harmless.  The second
selective flip is applied at a standard ordinary geometry, since
\eqref{eq:seed-exact-echo-cancellation} has already removed the huge
ordinary centers.  It acts trivially to \(o(r)\) on ordinary tails,
restores special spin \(y\), and restores the plus sign in transport.
At this second-flip probe start the exact special norm is still \(r\),
and the ordinary low action already has the updated unitary
\(W_Z(g_r)A_r\).  Lemma~\ref{lem:seed-low-insertions} consequently
approximates the actual insertion to \(O(r)\) by this updated low
insertion, irrespective of the temporarily negative transport sign.

The special insertion remains transported because each cost commutes
with \(Z\), while each selective flip anticommutes with \(Z\) on the
special component to error \(o(r)\).  In particular, this proof needs
only forward estimates on that component.  Boundedness of the residual
direction implies uniform moment bounds for the new \(A_r,A_r^*\).
Lemma~\ref{lem:seed-low-insertions} supplies bounded approximants to
the next actual insertion.  Finally, \(v'_r-v_r=O(r)\), so the updated
\(D_r\) retains the asserted bounded-approximant property.
\end{proof}

This induction also proves its claimed uniformity on input families.
The exact decomposition~\eqref{eq:seed-exact-special-norm} is linear in
the input.  At initialization every input has the same four special
centers.  Unscaled costs and compensated mixers preserve that center
list, and~\eqref{eq:seed-D-update} updates it independently of the tail.
The probe parameters are determined on the vacuum, but
Lemma~\ref{lem:seed-selective-accuracy} acts uniformly on the whole
family of tails having these centers.  Every discarded vector has norm
\(o(r)\) and is propagated by exact unitaries.  Since the total gate
and probe lists are fixed and finite before \(r\downarrow0\), their
errors add to \(o(r)\).  This establishes the invariant with no
assertion about number moments of the discarded remainders.

\subsection{The external-mode limit and the energy pairing}

The echo has now produced the bounded direction \(-2t w_r+o(1)\).
It remains to show that this direction acts as an independent Gaussian
seed, and that the resulting identification preserves root-shift
energy.  Ordinary Gram convergence alone does not give the latter
conclusion, so we compute the root-shift pairing explicitly.

We now identify the limiting low circuit and at the same time establish
the limits of the projections used to select the probes.  Introduce at
each vertex a two-dimensional real external mode space with orthonormal
vectors \(\ell,w\), and let \(E_{\mathrm{ext}}=\Span_{\C}\{\ell,w\}\).
Its Gaussian position operators commute and are independent standard
Gaussians in the vacuum.  To specify the recursion by height, put
\[
 \cH_{\mathrm{seed},0}=\C^2\otimes\Gamma_s(E_{\mathrm{ext}}),
 \qquad
 \cH_{\mathrm{seed},h+1}
 =\C^2\otimes\Gamma_s\!\left(
    \cH_{\mathrm{seed},h}^{\circ}\oplus E_{\mathrm{ext}}\right).
\]
The vacuum is \(\ket+\otimes\Omega\), and empty descendants give the
same compatible height embeddings as in the foundations.  In the mode
space at height \(h+1\), the first summand contains the whole centered
child rooted space, including that child's seeds.  The second summand
contains the two new seeds at the current root.  Equivalently, the root
Hilbert space is the spin times the pointed Gaussian space of these two
seeds times the symmetric Fock space over centered children.

The formal low circuit starts with
\begin{equation}
\label{eq:seed-formal-initialization}
 A^{\mathrm{form}}=R_\star e^{2\ii bZq(\ell)}.
\end{equation}
A bulk mixer and an unscaled cost have their usual seeded rules.  An
echo with parameter \(t\) has the formal update
\(e^{-2\ii t Zq(w)}\). At each vertex, \(q(w)\) is the Gaussian
seed of the requested word: the same coordinate is reused at every
seed gate, while different vertices have independent copies.
The coordinate \(q(\ell)\) is independent of those seeds and enters
through initialization; its effect is removed by taking \(b\) small
in the final proof. Thus an echo with \(t=\sigma/2\) has the seed
law of the requested gate of angle \(\sigma\), not a newly sampled
seed at each pulse. The root-shift calculation below identifies the
energy, which is not determined by this marginal law alone.

At a standard stage write
\(v_P^{\mathrm{form}}=(A^{\mathrm{form}})^*P
A^{\mathrm{form}}\vac\), with the appropriate accumulated word in
\(A^{\mathrm{form}}\).  Treat these insertions and \(e=Z\vac\) as
endogenous child directions.  In particular their inner products with
the external directions \(\ell,w\) are zero.

\begin{lemma}[Finite Gram induction]
\label{lem:seed-gram-induction}
Suppose the required selective probe lists have been constructed through
a given standard stage.  The joint Gram matrix of \(e\), all actual
standard-stage insertion vectors used so far, and \(\ell_r,w_r\),
converges to that of the corresponding formal endogenous directions
and the two orthogonal external directions.  The limiting energy of the
actual insertion converges to the energy of the formal seeded insertion.
The same assertions hold at the standard ordinary geometry just before
the second flip in an echo, with the temporary transport sign retained.
\end{lemma}

\begin{proof}
Every \(A_r\) in the invariant is a finite sum of root spin matrices
times Weyl displacements.  Its directions are finite linear combinations
of earlier actual insertion directions, \(d_{0,r}\), and the bounded
echo directions \(g_r\).  Its adjoint and its insertion vectors have
the same finite-Weyl-sum form.  The number of summands is finite for the
fixed circuit; no bound uniform in circuit length is required.

For such sums all vacuum inner products reduce to
\begin{equation}
\label{eq:seed-vacuum-weyl-gram}
 \braket{W(f)\Omega}{W(g)\Omega}
 =\exp\!\left(-\tfrac12\norm{g-f}^2
                 +\ii\im\braket f g\right).
\end{equation}
Thus only finitely many direction inner products are needed at a step.
The initial directions satisfy \(d_{0,r}=2b\ell_r+o(1)\), and the
echo directions satisfy \(g_r=-2tw_r+o(1)\).
Norm-small errors in these bounded directions may be removed using
Lemma~\ref{lem:seed-packet-estimates}.  Every actual standard insertion
differs by \(O(r)\) from its bounded smooth low insertion, by
Lemma~\ref{lem:seed-low-insertions}.  Thus, when computing a new Gram
entry, an earlier actual insertion may first be replaced by its low
insertion at error \(O(r)\).  Lemma~\ref{lem:seed-high-directions}
therefore makes its pairing with \(\ell_r,w_r\) vanish.  Together with
\(\braket{\ell_r}{w_r}\to0\), this is exactly the orthogonal-sum rule
for external and endogenous directions.  Formula
\eqref{eq:seed-vacuum-weyl-gram} now proves the Gram assertion inductively.

Here is the additional calculation for the root-shift energy, so that
ordinary Gram convergence is not silently substituted for convergence
of a different operator.  A low \(Y\)- or \(Z\)-insertion is centered:
the circuit preserves the joint spin-flip parity and these matrices are
odd.  Write such an insertion as
\[
 h=\sum_{j=1}^J s_j\otimes W(f_j)\Omega,
 \qquad s_j\in\C^2.
\]
In the unseeded representation,
\(\cL h=\ket+\otimes a^*(h)\Omega\).  The one-particle coefficient
of \(W(f)\Omega\) is
\(\ii e^{-\norm f^2/2}f\).  Consequently
\begin{equation}
\label{eq:seed-energy-weyl-sum}
 \braket h{\cL h}
 =-\ii\sum_{j=1}^J
       \braket{s_j}{+}\,e^{-\norm{f_j}^2/2}\braket{f_j}h.
\end{equation}
Each pairing \(\braket{f_j}h\) is one of the same finite linear
combinations of insertion and external-direction Grams.  In the formal
seeded representation the same formula holds, with \(h\) inserted in
the endogenous summand; the external components of \(f_j\) then pair
to zero. More explicitly, the external two-mode Fock factor represents
the Gaussian variables at the root, while the child Fock factor is
built over endogenous directions. The root shift puts the root-seed
factor in its pointed vacuum and creates \(h\) only in that child
factor. The two-factor decomposition is identified with Fock space
over their orthogonal direct sum.  The high-direction argument gives
exactly this zero pairing in the actual limit.  Finally, the actual insertion differs in norm
from \(h\) by \(O(r)\).  Since \(\cL\) is an isometry,
\[
 \abs{\cE(v)-\cE(h)}
 \leq(\norm v+\norm h)\norm{v-h},
\]
which proves energy convergence.  A finite scalar bulk phase cancels
from every insertion.  The same argument applies just before an echo's
second flip, because its ordinary centers have already canceled and its
low update has already been made.
\end{proof}

The standard-stage Gram limits determine the projections needed for
ordinary probe starts.  For a probe immediately after a cost of size
\(t/r\), we also need the following transverse projection estimate.

\begin{lemma}[The huge-stage transverse projection]
\label{lem:seed-huge-transverse}
Immediately after a forward cost with fixed \(t\neq0\), the actual
\(Y\)-insertion has vanishing pairing with every bounded family in
\(\mathcal S\).  In particular, its pairing with every special scaled
center tends to zero.
\end{lemma}

\begin{proof}
For a bounded input \(\eta_r\), use the standard-stage uniform forward
approximation on both \(\vac\) and \(\eta_r\).  After the huge cost
their ordinary pieces have centers \(-tyv_r\) and root spins \(y\).
In the pairing defining
\(\braket{u_r^*Yu_r\vac}{\eta_r}\), the inserted \(Y\) reverses the
spin.  Nonzero spin pairings therefore require the two opposite
ordinary centers.  Their distance is \(2|t|\), so their contribution
vanishes by~\eqref{eq:seed-packet-separation}.  Every contribution with
a special component is \(O(r)\); the ordinary remainders are \(o(r)\).
This proves the first assertion.  Before the huge cost, the invariant
approximates every special center by a bounded smooth vector to
\(O(r)\).  In particular the intermediate directions \(v'_r\) in
\eqref{eq:seed-D-update} differ by \(O(r)\) from the corresponding
earlier standard insertions, by~\eqref{eq:seed-flip-change}.
Adding the forward displacement
\(-tyv_r\) preserves this property.  The first assertion and
Cauchy--Schwarz prove the second.
\end{proof}

\subsection{Visibility and the order of construction}

The preceding arguments are conditional only on the visibility
conditions.  We now choose the target angles and all finite probe lists
so that those conditions hold.  The order is important: first use the
formal low circuit to perturb the target angles, then construct the
probe lists successively from limits established at earlier stages,
and only afterwards choose \(r\).  Fix a finite seeded word, regard all
its angles as independent real variables \(\lambda\), and replace each
seed gate of angle \(\sigma\) formally by an echo parameter
\(t=\sigma/2\).  Include in a finite list all
standard-stage probe starts: compensated mixers and the second flips
of echoes.  Also include the standard stage preceding each forward huge
cost.  At a second-flip probe start the current echo's reverse cost has
already occurred: count that echo in the center accumulator and include
its seed displacement in the formal ordinary state.  Its packet label
\(y\) is still the label before the echo, although its current spin is
\(-y\).  Define, at each probe start,
\[
 D^{\mathrm{form}}
 =-2\sum_{j:\,\text{reverse cost already applied}}
       t_jv_{Z,j}^{\mathrm{form}},
\]
where the insertion is taken immediately before that echo.  At each
listed stage and for each \(z,y\), impose
\begin{equation}
\label{eq:seed-formal-visibility}
 a_{zy}^{\mathrm{form}}(\lambda)
 =-2\im\braket{v_Z^{\mathrm{form}}}
                   {-2bz e+yD^{\mathrm{form}}}\neq0.
\end{equation}

Each function in~\eqref{eq:seed-formal-visibility} is real analytic in
the finitely many angles.  One may verify this without an
unbounded-operator analyticity theorem: expand the finite word into
Weyl sums and use~\eqref{eq:seed-vacuum-weyl-gram} recursively.  These
expressions use finitely many sums, products, and exponentials of
previous finite Gram expressions.

None of these analytic functions is identically zero.  To see this,
set \emph{all} angles of the target word to zero, while keeping the
initialization fixed.  Every listed formal standard stage is then
\eqref{eq:seed-formal-initialization}, and
\(D^{\mathrm{form}}=0\).  A direct one-Gaussian calculation gives
\[
 \braket{v_Z^{\mathrm{form}}}e
 =\ii e^{-8b^2},\qquad
 a_{zy}^{\mathrm{form}}=4bz e^{-8b^2}\neq0.
\]
The temporary minus sign in special transport before a second flip
does not change this formal ordinary insertion or its center formula.
Thus the same test point works for every listed stage.

The complement of the finite union of these analytic zero sets, and
of the hyperplanes \(t_j=0\), is dense.  Perturb the target angles
arbitrarily little into that complement.  Energy changes arbitrarily
little because it is continuous by the same finite-Weyl-sum formulas.

We can now construct the actual pulses successively.  At a standard
probe start, Lemma~\ref{lem:seed-gram-induction} supplies the limits of
both special projection coordinates, and
\eqref{eq:seed-formal-visibility} makes their first coordinate nonzero.
Choose the finite ensemble-control list from these limiting pairs.
Lemma~\ref{lem:seed-selective-accuracy} implements it with the required
errors.  At a huge-stage first flip, the first coordinate is unchanged
by the forward displacement, while the second coordinate tends to zero
by Lemma~\ref{lem:seed-huge-transverse}.  Its limiting pairs are
therefore \((a_{zy}^{\mathrm{form}},0)\), again nonzero.  They too
admit the ensemble-control construction.  Continue the packet and Gram
inductions to the next stage.  The argument at the second flip uses the
already updated formal low circuit, with the sign of transport still
negative until that flip is undone.  In particular no pulse list is
chosen using a limit that presupposes its own implementation.

\begin{proof}[Proof of Theorem~\ref{thm:seed-simulation}]
Start with the requested finite seeded word and a tolerance.  For a
small fixed \(b>0\), prepend
\(R_\star e^{2\ii b Zq(\ell)}\) in the formal seeded calculus, with
\(\ell\) an additional independent Gaussian coordinate at every site.
As \(b\downarrow0\), its energy tends to that of the original seeded
word.  Indeed, finite bounded-direction Weyl words are strongly
continuous in this initial displacement, as are their insertion and
energy Grams.  At \(b=0\) the initial global mixer acts on the original
\(\ket+\) product state by a scalar, so it changes no expectation;
the extra Gaussian coordinate is then unused.

Choose \(b\) sufficiently small for this error, and perturb the target
angles sufficiently little that all the visibility conditions hold,
with another arbitrarily small energy error.  Fix the resulting finite
probe lists by the successive construction above.  With these lists
fixed, take \(r>0\) sufficiently small.  The packet induction and
Lemma~\ref{lem:seed-gram-induction} show that the ordinary word consisting
of the initialization, unscaled costs, compensated mixers, and seed
echoes has energy arbitrarily close to that of the perturbed formal
seeded word.  Summing the three errors proves the assertion.

Every gate used is a cost or mixer pulse, possibly with negative angle.
Toggled probes and all compensations are finite such words; consecutive
gates of the same kind can be combined, or zero-angle gates inserted
to obtain the stipulated alternating QAOA format.  After \(b\), the
perturbation, the finite probe lists, and \(r\) have been chosen, the
entire word and its angles are fixed.  Proposition~\ref{prop:tree-transfer}
therefore applies before any approximation parameter is changed.
This proves the claimed order of limits and parameter independence.
\end{proof}

\section{The ordinary word and the SK limit}\label{sec:conclusion}

The preceding constructions are now fixed in their order of use.
The following error allocation leaves every circuit finite before
applying the tree-to-SK identity.

\begin{proof}[Proof of Theorem~\ref{thm:main}]
Fix $\varepsilon>0$. Apply
Proposition~\ref{prop:classical-near-optimum} with error
$\varepsilon/4$ and helper bound $\eta_H=(\varepsilon/8)^2$.
It gives a finite bounded smooth classical calculation with output $m$
and helper probability $p_H$ satisfying
\[
 \cE(m)\ge\Pstar-\varepsilon/4,
 \qquad 2\sqrt{p_H}<\varepsilon/4.
\]
Apply Theorem~\ref{thm:classical-synthesis} with $\eta=\varepsilon/8$.
The resulting finite seeded circuit has insertion $v$ with
\[
 \cE(v)\ge\cE(m)-2\eta-2\sqrt{p_H}
             >\Pstar-3\varepsilon/4.
\]
By Theorem~\ref{thm:seed-simulation}, a finite ordinary word approximates
this seeded value to error less than $\varepsilon/4$.
Its value is therefore greater than $\Pstar-\varepsilon$.

All choices just made concern limiting tree calculations and deterministic
approximation parameters. After these choices there is a single finite word
with real angles independent of $D$, $n$, and $J$. By
Proposition~\ref{prop:tree-transfer}, its tree value equals its fixed-parameter SK
value. Write the word in the form~\eqref{eq:qaoa-state}; it gives
$v_p(\gamma,\beta)\ge\Pstar-\varepsilon$ for some finite $p$.
The upper bound and monotonicity established above finish the proof.
\end{proof}

\begin{remark}[Order of choices]\label{rem:order-of-choices}
For the primary route the logical order is: accuracy; a fixed
zero-temperature minimizer and stopping time $T<1$; a sufficiently
fine finite scalar mesh at that $T$; finite sign smoothing; bounded
column and trigonometric approximation; type exclusions;
helper contractions and finite commutator approximations;
elementary seed compilation; small initialization parameter $b$; generic
angle perturbation; finite selective-probe lists; sufficiently small
positive $r$; one fixed ordinary word; thermodynamic limit. Within any nested construction
an outer accuracy is fixed before a finite approximation of its inner
operations is chosen. No uniform depth estimate and no exchange of these
limits is used. For the independent positive-temperature route, the
initial scalar choices are instead a sufficiently large fixed finite
$\beta_0$, a cutoff $T<q_{\beta_0}$, and a sufficiently fine finite
message-passing discretization. Its degree limit is taken with those
parameters fixed. The subsequent approximation and circuit choices
are the same.
\end{remark}

\section{Proof of the regular MaxCut consequence}
\label{sec:regular-maxcut-proof}

We now prove Corollary~\ref{cor:regular-maxcut}. Theorem~\ref{thm:main}
provides one fixed word with SK value close to $\Pstar$, and
Proposition~\ref{prop:tree-transfer} identifies that value with the
large-degree positive-cost tree correlation. The proof first transfers
this correlation to every graph satisfying the stated girth condition.
It then uses local tree structure and the random-regular optimum
asymptotics to obtain the separate successive-limit comparison.

\begin{proof}[Proof of Corollary~\ref{cor:regular-maxcut}]
The normalized negative cost used by Basso et al.~\cite{BFMVZ2022} is
\[
 C_D^-=-\frac1{\sqrt D}\sum_{\{u,v\}\in E}Z_uZ_v
      =\frac2{\sqrt D}H_{\mathrm{MC}}(G)
       -\frac{|E|}{\sqrt D}I.
\]
Thus its cost gate with angle $\gamma_j$ is, up to a scalar phase, the usual
MaxCut cost gate with angle $2\gamma_j/\sqrt D$, with the same sign and
unchanged mixer angle.  Let $r_p^+(D,\gamma,\beta)$ be the distinguished-edge
correlation for the positive tree cost in
Proposition~\ref{prop:tree-transfer}, and put
\[
 \nu_p(D,\gamma,\beta)=\frac{\sqrt D}{2}r_p^+(D,\gamma,\beta).
\]
That proposition gives $\nu_p(D,\gamma,\beta)\to v_p(\gamma,\beta)$ for
each fixed word.

For the observable at an edge, backward cancellation of commuting cost
gates leaves only cost edges having an endpoint at distance at most $p-1$
from the distinguished edge's endpoint set.  An extra edge in this
relevant cone would close a cycle of length at most $2p+1$.  Under the
stated girth condition the cone is therefore the regular tree cone.
Conjugating \emph{only this cone} by $X$ on one bipartition class reverses
the cost sign and the distinguished $Z_uZ_v$, while preserving the mixers
and the initial plus state.  No bipartition of $G$ is used.  Every edge
therefore has negative-cost correlation $-r_p^+(D,\gamma,\beta)$, and
\begin{equation}\label{eq:regular-tree-cut}
 q_{p,G}(2\gamma/\sqrt D,\beta)
 =\frac12+\frac{\nu_p(D,\gamma,\beta)}{\sqrt D}.
\end{equation}
This is the finite-degree edge formula in the full version
of~\cite[Equations~(2.3)--(2.6) and (3.1)]{BFMVZ2022}, in the convention
of Proposition~\ref{prop:tree-transfer}.  Choose the word by
Theorem~\ref{thm:main} with $v_p(\gamma,\beta)\ge\Pstar-\eta/4$.
For all sufficiently large $D$, fixed-word convergence gives
$\nu_p(D,\gamma,\beta)\ge\Pstar-\eta/2$, proving
\eqref{eq:high-girth-maxcut}.

For fixed $D,p$, let $b_{p,D}(G)$ be the fraction of edges whose relevant
cone is not a regular tree cone.  A conservative radius-$(p+1)$
neighborhood test bounds the number of such edges by a constant depending
on $D,p$ times the number of cycles of length at most $2p+3$.
Those fixed-length cycle counts are tight in the uniform simple regular
model by McKay, Wormald, and Wysocka~\cite[Theorem~1]{MWW2004}.
Consequently $b_{p,D}(G_{n,D+1})\to0$ in graph probability and in mean
as $n\to\infty$.  Since each edge contributes a cut probability in $[0,1]$,
the exact value on the other edges gives
\[
 \left|q_{p,G}(2\gamma/\sqrt D,\beta)
 -\frac12-\frac{\nu_p(D,\gamma,\beta)}{\sqrt D}\right|
 \le b_{p,D}(G).
\]
This proves convergence of the conditional quantum expectation to the
tree value at fixed $D,p$; the whole random graph need not have large
girth.

Dembo, Montanari, and Sen~\cite[Theorem~1.6(b), with the convention after
Equation~(1.7)]{DMS2017} identify the optimal leading correction for
the same simple regular model.  After dividing their cut size by $nd/2$
and putting $d=D+1$, their theorem says that
$\sqrt D(\operatorname{OPT}(G_{n,D+1})-1/2)$ converges to $\Pstar$ in
the same successive graph-probability sense; here
$\sqrt{D/(D+1)}\to1$.  Their $\Pstar$ is ours: pairing their ordered
off-diagonal Gaussians and using sign symmetry gives~\eqref{eq:sk-model}
in law, while the diagonal is a mean-zero configuration-independent
shift.  For large $D$ the deterministic quantity $\Pstar-\nu_p(D,\gamma,\beta)$
is at most $\eta/2$.  If $\sqrt D\,\Delta_D>\eta$, the absolute centered-optimum
error or the absolute conditional-expectation error multiplied by
$\sqrt D$ must therefore exceed $\eta/4$.  The first probability vanishes
in the stated successive limit by the cited optimum theorem, and the
second vanishes in the inner $n$ limit by the preceding bound.  This
proves~\eqref{eq:random-regular-maxcut}.
\end{proof}

\appendix
\section{The independent positive-temperature construction}
\label{sec:positive-temperature-alternative}

This appendix constructs near-optimal bounded classical formulas
without using the zero-temperature scalar companion.  Its input is the
positive-temperature Parisi measure at a fixed finite inverse
temperature.  The Parisi diffusion first gives an almost optimal
martingale output.  A finite-step tree calculation and a strong
symmetric-power limit then place that output in the rooted Hilbert
space.  Finally, the bounded approximation and color-exclusion lemmas
of Section~\ref{sec:classical-optimization} put it in the form required
by Theorem~\ref{thm:classical-synthesis}.

Besides providing an independent route to the same synthesis input,
the appendix retains the finite-degree message calculation and the
quantitative fixed-temperature mesh bounds.  Temperature is chosen
first and held fixed in those estimates.  No passage from the
log-cosh terminal datum here to the terminal datum $|x|$ is used.

\subsection{The positive-temperature variational input}

The use of a Parisi diffusion and its gradient martingale to choose
incremental message-passing coefficients follows the classical
construction of Montanari~\cite[Section~3.1, Lemmas~3.2--3.3]{Montanari2019},
with the more general framework of El Alaoui, Montanari, and
Sellke~\cite[Assumption~1 and Theorem~2]{AMS2020}. We use the finite-tree
calculation of El Alaoui, Montanari, and Sellke~\cite{EAMS2023} and
convert its output to the synthesis formulas below.

We record explicitly the spin-glass results that enter the argument.
For an inverse temperature $\beta_0>0$, let $\mu_{\beta_0}$ be the
finite-temperature Parisi measure and put
\[
 \gamma(t)=\beta_0\mu_{\beta_0}([0,t]).
\]
In field coordinates appropriate to the normalization of this paper,
the Parisi PDE and free energy are
\begin{equation}
 \begin{aligned}
 \phi_t(t,x)
   &=-\frac12\bigl(\phi_{xx}(t,x)+\gamma(t)\phi_x(t,x)^2\bigr),\\
 \phi(1,x)&=\frac1{\beta_0}\log\bigl(2\cosh(\beta_0x)\bigr).
 \end{aligned}
 \label{eq:positive-parisi-pde}
\end{equation}
\begin{equation}
 F_{\beta_0}
   =\phi(0,0)-\frac12\int_0^1t\gamma(t)\,\dd t
    =\lim_{n\to\infty}\frac1{n\beta_0}
        \E\log\sum_{\sigma}\exp\bigl(\beta_0h_{n,J}(\sigma)\bigr).
 \label{eq:positive-parisi-value}
\end{equation}
Equation~\eqref{eq:positive-parisi-value} is the positive-temperature
Parisi formula~\cite[Theorem~1.1]{Talagrand2006}, written in the field
coordinates of \eqref{eq:positive-parisi-pde}. The covariance of
$h_{n,J}$ divided by $n$ is $(R_{12}^2-1/n)/2$, where
$R_{12}=n^{-1}\sum_i\sigma_i^1\sigma_i^2$. The $O(1/n)$ term is
allowed by that theorem, so its limiting covariance is $t^2/2$.
The partition function here is the unnormalized sum over configurations.

This route does not need the zero-temperature variational formula,
even for existence of the limit defining $P_*$. Put
$A_n=n^{-1}\E\max_\sigma h_{n,J}(\sigma)$ and
let $F_{n,\beta_0}$ denote the finite-$n$ expression on the right of
\eqref{eq:positive-parisi-value}. The elementary log-sum bounds give
\[
 F_{n,\beta_0}-\frac{\log2}{\beta_0}
 \le A_n\le F_{n,\beta_0}.
\]
For every fixed finite $\beta_0$, the Parisi formula gives convergence
of $F_{n,\beta_0}$. Hence
$\limsup_n A_n-\liminf_n A_n\le(\log2)/\beta_0$.
Letting $\beta_0$ tend to infinity proves convergence of $A_n$ and yields
\begin{equation}
 P_*\le F_{\beta_0}\le P_*+\frac{\log2}{\beta_0}.
 \label{eq:positive-free-energy-bounds}
\end{equation}
A direct zero-temperature variational formula is also known
\cite[Theorem~1]{AC2017}; the argument here needs only the
positive-temperature formula and the preceding squeeze. We use the
following precise finite-temperature facts.  In the conventions of
\cite{Lopatto2026}, they are Theorem~1.1, Lemma~2.2,
Proposition~2.4, and Lemma~2.5, respectively:
\begin{enumerate}
 \item If $\beta_0>1$, there is $q=q_{\beta_0}\in(0,1)$ such that
 $\supp\mu_{\beta_0}=[0,q]$.  The restriction below $q$ has a density
 in $C^\infty([0,q))$; the point $q$ is an atom.
 \item If $X$ solves the stochastic differential equation driven by
 a standard Brownian motion $W$,
 \begin{equation}
  \dd X_t=\gamma(t)\phi_x(t,X_t)\,\dd t+\dd W_t,
  \qquad X_0=0,
  \label{eq:positive-parisi-diffusion}
 \end{equation}
 then $\E\phi_x(t,X_t)^2=t$ for every
 $t\in\supp\mu_{\beta_0}$.
 \item All positive-order spatial derivatives of $\phi$ are bounded
 and jointly continuous on $[0,1]\times\R$.
 The function $\phi$ is even in $x$, $|\phi_x|\le1$, and
 $\phi_{xx}>0$.  These regularity statements also follow from
 \cite[Propositions~1--2 and Appendix Proposition~5, (66)]{AC2015}
 after the same change of variables.
 \item Differentiating the PDE along
 \eqref{eq:positive-parisi-diffusion} gives
 $\dd\phi_x(t,X_t)=\phi_{xx}(t,X_t)\,\dd W_t$.
\end{enumerate}
The self-consistency identity $\E\phi_x(t,X_t)^2=t$ at points
$t\in\supp\mu_{\beta_0}$, including $t=0$, also follows from the
variational argument of Jagannath and Tobasco~\cite[Proposition~1.1 and Corollary~3.10]{JT2017}.
To check the conventions, if $u$ denotes the PDE solution with terminal
condition $\log\cosh x$ and diffusion coefficient $\beta_0^2$ in
\cite{Lopatto2026}, then
$\phi(t,x)=\beta_0^{-1}\bigl(u(t,\beta_0x)+\log2\bigr)$.
In particular, the full-support input used here concerns positive
temperature; no support property of a limiting zero-temperature
minimizer is assumed. The version of \cite{Lopatto2026} used here is
version~3: its Hamiltonian is written with independent Gaussian terms
over all ordered pairs. Combining the two terms for each unordered
pair gives the normalization \eqref{eq:sk-model}; the diagonal terms
form a centered, configuration-independent Gaussian. They have zero
contribution to expected pressure and do not change the Parisi
minimizer.

Write
\begin{equation}
 m(t,x)=\phi_x(t,x),\qquad
 a(t,x)=\phi_{xx}(t,x),\qquad
 b(t,x)=\gamma(t)m(t,x).
 \label{eq:classical-coefficients}
\end{equation}
On every $[0,T]$ with $T<q$, these functions have bounded spatial
derivatives of every order.  They are also Lipschitz in time there:
the density statement makes $\gamma$ smooth on $[0,T]$, and the PDE
expresses the time derivatives of $m$ and $a$ in terms of bounded
spatial derivatives.  In particular, $a$ and $b$ are bounded and
globally Lipschitz in $x$, uniformly on $[0,T]$.
The function $a$ is even in $x$ and $b$ is odd.

\begin{lemma}[A positive-temperature value bound]
\label{lem:positive-temperature-bound}
For every $\beta_0>1$ and $0<T<q$,
\begin{equation}
 \int_0^T\E a(t,X_t)\,\dd t
 \ge P_* -\frac{1/4+\log2}{\beta_0}-(q-T).
 \label{eq:positive-temperature-bound}
\end{equation}
Moreover,
\begin{equation}
 \E a(t,X_t)^2=1\quad(0\le t\le q),
 \qquad 1-q\le\beta_0^{-1}.
 \label{eq:parisi-normalization-and-endpoint}
\end{equation}
\end{lemma}

\begin{proof}
Symmetry gives $m(0,0)=0$.  The martingale identity and It\^o isometry
therefore give
\[
 \E m(t,X_t)^2=\int_0^t\E a(s,X_s)^2\,\dd s.
\]
The left-hand side equals $t$ on $[0,q]$.  Continuity and boundedness
of $a$, together with continuity of $X$, prove the first assertion of
\eqref{eq:parisi-normalization-and-endpoint}, including both endpoints.

Since $\gamma=\beta_0$ on $[q,1]$, the Cole--Hopf formula gives
\begin{equation}
 \phi(q,x)=\frac1{\beta_0}\log\bigl(2\cosh(\beta_0x)\bigr)
                  +\frac{\beta_0(1-q)}2,
 \quad
 m(q,x)=\tanh(\beta_0x),
 \quad
 a(q,x)=\beta_0\operatorname{sech}^2(\beta_0x).
 \label{eq:parisi-endpoint-formula}
\end{equation}
Thus, by Cauchy--Schwarz and the identity just proved,
\[
 1-q=\E\operatorname{sech}^2(\beta_0X_q)
 \le\bigl(\E\operatorname{sech}^4(\beta_0X_q)\bigr)^{1/2}
 =\beta_0^{-1}.
\]

It\^o's formula applied to $X_tm(t,X_t)$ and $\phi(t,X_t)$ yields
\begin{align*}
 \int_0^q\E a(t,X_t)\,\dd t
   &=\E X_qm(q,X_q)-\int_0^qt\gamma(t)\,\dd t,\\
 \E\phi(q,X_q)
   &=\phi(0,0)+\frac12\int_0^qt\gamma(t)\,\dd t.
\end{align*}
The stochastic integrals have mean zero because $a,m,b$ are bounded
and $X$ has finite moments.  The elementary inequality
\[
 x\tanh(\beta_0x)
 \ge\frac1{\beta_0}\log\bigl(2\cosh(\beta_0x)\bigr)
              -\frac{\log2}{\beta_0}
\]
follows, for example, by differentiating the difference for $x\ge0$
and using evenness.  Substitution of
\eqref{eq:parisi-endpoint-formula} consequently gives
\begin{align*}
 \int_0^q\E a(t,X_t)\,\dd t
 &\ge\phi(0,0)-\frac12\int_0^qt\gamma(t)\,\dd t
                -\frac{\beta_0(1-q)}2-\frac{\log2}{\beta_0}\\
 &=F_{\beta_0}-\frac{\beta_0(1-q)^2}{4}
                       -\frac{\log2}{\beta_0}\\
 &\ge P_*-\frac{1/4+\log2}{\beta_0}.
\end{align*}
Finally $0\le\E a(t,X_t)\le(\E a(t,X_t)^2)^{1/2}=1$ on
$[0,q]$, so removing $[T,q]$ costs at most $q-T$.
\end{proof}

\subsection{The finite-iteration tree calculation}

Here we state the finite-step input from El Alaoui, Montanari, and
Sellke with the indices and normalizations needed below.  It is the
general calculation in \cite[Sections~2--4.1]{EAMS2023}, not that
paper's zero-temperature near-optimality theorem.

Fix $\beta_0>1$ and $0<T<q$. Choose $\delta>0$ and an integer
$L\ge1$ with $L\delta\le T$.
Put $t_r=r\delta$.  On the $(D+1)$-regular tree, initialize independent
standard Gaussian labels $g_i$ and set
$u^0_{i\to j}=u_i^0=g_i$, $A^{-1}_{i\to j}=1$.
Each outgoing column multiplies a message by a coefficient from the
preceding step.  This lag will make its newest Gaussian layer
conditionally centered, as verified below.
For $0\le r<L$ use the recursions
\begin{align}
 u^{r+1}_{i\to j}
   &=\frac1{\sqrt D}\sum_{v\in\partial i\setminus\{j\}}
                       A^{r-1}_{v\to i}u^r_{v\to i},
 &u^{r+1}_i
   &=\frac1{\sqrt{D+1}}\sum_{v\in\partial i}
                       A^{r-1}_{v\to i}u^r_{v\to i},
 \label{eq:eams-message-recursion}\\
 x^{r+1}_{i\to j}
   &=x^r_{i\to j}+\delta b(t_r,x^r_{i\to j})
                       +\sqrt\delta\,u^{r+1}_{i\to j},
 &x^0_{i\to j}&=\sqrt\delta\,g_i,
 \label{eq:eams-cavity-euler}\\
 x^{r+1}_i
   &=x^r_i+\delta b(t_r,x^r_i)+\sqrt\delta\,u^{r+1}_i,
 &x^0_i&=\sqrt\delta\,g_i.
 \label{eq:eams-node-euler}
\end{align}
The coefficients, for $0\le r\le L$, are
\begin{align}
 c_{D,r}&=\bigl(\E a(t_r,x^r_{i\to j})^2\bigr)^{-1/2},
 &A^r_{i\to j}&=c_{D,r}a(t_r,x^r_{i\to j}),
 &B^r_i&=c_{D,r}a(t_r,x^r_i),
 \label{eq:eams-finite-normalization}\\
 z_i^L&=\sqrt\delta\sum_{r=1}^LB^{r-1}_iu_i^r.
 \label{eq:eams-output}
\end{align}
The same constant, computed from the directed-edge law, occurs in
$A^r$ and $B^r$.  Its definition is recursive: $x^r_{i\to j}$ has
already been constructed when $c_{D,r}$ is chosen.  Strict positivity
of $a$ ensures every denominator is positive.

The one-step lag in $A^{r-1}u^r$ is important. Let
$\mathcal G_{i\to j}^r$ be the sigma-field of Gaussian labels at
distance at most $r$ from $i$ in the component left after deleting
$\{i,j\}$, including the label at $i$. The variables
$x^{r-1}_{i\to j}$ and $A^{r-1}_{i\to j}$ are
$\mathcal G_{i\to j}^{r-1}$-measurable. The next message is linear
in the Gaussian labels on the new outer layer. More precisely,
\cite[Lemma~3.2]{EAMS2023} gives, for $r\ge1$,
\begin{equation}
 \begin{aligned}
 u^r_{i\to j}\mid\mathcal G_{i\to j}^{r-1}
       &\sim N\bigl(0,(\tau^r_{i\to j})^2\bigr),\\
 (\tau^{r+1}_{i\to j})^2
       &=\frac1D\sum_{v\in\partial i\setminus\{j\}}
       (A^{r-1}_{v\to i})^2(\tau^r_{v\to i})^2,
       \qquad \tau^0=1.
 \end{aligned}
 \label{eq:eams-conditional-innovation}
\end{equation}
The variance recursion also holds at $r=0$. Thus $A^{r-1}u^r$
has conditional mean zero given this earlier cavity history. Its
variance may depend on that history; finite-degree independence of
successive messages is not being asserted. For $r=0$ the sent column
is simply $g_i$.

Let $U_0,U_1,\ldots$ be independent standard Gaussians and define
\begin{align}
 X^\delta_0&=\sqrt\delta\,U_0,
 &X^\delta_{r+1}
  &=X^\delta_r+\delta b(t_r,X^\delta_r)+\sqrt\delta\,U_{r+1},
 \label{eq:eams-limit-euler}\\
 c_{\delta,r}&=\bigl(\E a(t_r,X^\delta_r)^2\bigr)^{-1/2},
 &Z^\delta_L
  &=\sqrt\delta\sum_{r=1}^L
        c_{\delta,r-1}a(t_{r-1},X^\delta_{r-1})U_r.
 \label{eq:eams-limit-output}
\end{align}
At fixed $\beta_0,T,\delta,L$, the update functions and constants
$c_{\delta,r}$ are deterministic. The variables $X^\delta_r$ and
$Z^\delta_L$ are random functions of the displayed Gaussian family.

\begin{lemma}[Finite-step limit]
\label{lem:eams-finite-step}
For the preceding construction, as $D\to\infty$ at fixed $\delta,L$,
\begin{align}
 c_{D,r}&\longrightarrow c_{\delta,r},
 \qquad 0\le r\le L,
 \label{eq:eams-normalizer-limit}\\
 (u^0_i,\ldots,u^L_i)&\Longrightarrow(U_0,\ldots,U_L),
 \qquad z_i^L\Longrightarrow Z^\delta_L,
 \label{eq:eams-distribution-limit}\\
 \lim_{D\to\infty}\frac{D+1}{2\sqrt D}\E z_o^Lz_v^L
   &=\delta\sum_{r=1}^{L-1}
               c_{\delta,r}\E a(t_r,X^\delta_r),
 \qquad o\sim v.
 \label{eq:eams-corrected-value}
\end{align}
The sum in \eqref{eq:eams-corrected-value} is empty when $L=1$.
The distributional statements include convergence of every fixed
moment.  The analogous joint Gaussian limit holds for a directed edge.
Moreover $\E(Z^\delta_L)^2=L\delta$.
\end{lemma}

\begin{proof}
We use the bounded-coefficient CLT and edge calculation in
\cite[Assumptions~2--3, Theorem~3.1, Proposition~2.1, and
Lemmas~4.2--4.3]{EAMS2023}.  The verification below has two parts:
first the degree-dependent normalizers and moment bounds, then the
edge value with its finite-step index.
The functions $a(t_r,\cdot)$ are bounded and Lipschitz, the drift is
Lipschitz, and only finitely many $r$ occur.  Initially
$c_{D,0}=c_{\delta,0}$, and $x^r$ uses normalizers only through
index $r-2$, so the normalization induction is not circular.
Inductively, the Gaussian
limit of the preceding messages and bounded continuity of $a^2$ give
\[
 \E a(t_r,x^r_{i\to j})^2\longrightarrow
 \E a(t_r,X^\delta_r)^2>0.
\]
Thus $c_{D,r}$ converges, and the normalized coefficients and their
Lipschitz constants are bounded uniformly for all sufficiently large
$D$.  This verifies the boundedness and normalization hypotheses used
in the finite-step CLT, including when its coefficients depend on $D$.
One may equivalently run its induction with this triangular sequence:
convergent deterministic coefficients do not affect any induction
step. All fixed moments are uniformly bounded. Indeed, an update
sums the columns $A^{r-1}_{v\to i}u^r_{v\to i}$ from disjoint
child cavity trees. They are independent across children and centered
by \eqref{eq:eams-conditional-innovation}; the bounded multiplier
is measurable before the innovation it multiplies. In this application
their marginal centering also follows from seed parity, since $a$ is
even and $b$ is odd. The moment inequality for sums of independent
centered variables gives, for each integer $k\ge2$,
\[
 \E\left|D^{-1/2}\sum_{j=1}^DY_j\right|^k
 \le C_k\left((\E Y_1^2)^{k/2}
                  +D^{1-k/2}\E|Y_1|^k\right).
\]
Induction over the finite number of updates proves uniform bounds of
every order.  They imply uniform integrability and hence the asserted
moment convergence.  The $x$ and $z$ conclusions follow from their
continuous, polynomial-growth expressions in the messages.

For the energy, put
\[
 R_{D,\ell}=\E\left[
 B_o^{\ell-1}B_v^{\ell-1}
 A_{o\to v}^{\ell-2}A_{v\to o}^{\ell-2}
 u_{o\to v}^{\ell-1}u_{v\to o}^{\ell-1}\right].
\]
Proposition~2.1 of \cite{EAMS2023}, with its degree $k=D+1$, is
the exact identity
\begin{equation}
 \begin{aligned}
 \frac{D+1}{2\sqrt D}\E z_o^Lz_v^L
 &=\delta\sum_{\ell=2}^L
   \E\left[A^{\ell-2}_{v\to o}(u^{\ell-1}_{v\to o})^2
                           B_o^{\ell-1}B_v^{\ell-2}\right]\\
 &\quad+\frac{\delta}{2\sqrt D}\sum_{\ell=1}^L R_{D,\ell}.
 \end{aligned}
 \label{eq:eams-exact-edge-expansion}
\end{equation}
The second line is $O(D^{-1/2})$ by the fixed moment bounds.
To compare node and cavity terms, split off the missing neighbor:
\[
 u_i^r=\sqrt{\frac D{D+1}}\,u_{i\to j}^r
       +\frac{A_{j\to i}^{r-2}u_{j\to i}^{r-1}}{\sqrt{D+1}}
 \qquad (1\le r\le L).
\]
Its squared difference from $u_{i\to j}^r$ is $O(D^{-1})$.
The two Euler recursions and the spatial Lipschitz bounds therefore
give, by induction over this fixed finite number of steps,
\begin{equation}
 \E|x_i^r-x_{i\to j}^r|^2+
 \E|B_i^r-A_{i\to j}^r|^2\le \frac C D.
 \label{eq:eams-node-cavity-bound}
\end{equation}
These are the comparisons in Lemma~4.2 of \cite{EAMS2023}; in
particular, the use of the spatial Lipschitz bound for $a$ is explicit.

Here is the reduction in Lemma~4.3 of that paper in our indices.
Set $s=\ell-1$ and
$H=A_{v\to o}^{s-1}B_v^{s-1}B_o^s$. This product depends only on
the depth-$(s-1)$ cavity labels on the $v$ side and the depth-$s$
cavity labels on the $o$ side. Conditioning on these labels in
\eqref{eq:eams-conditional-innovation} therefore gives
\[
 \E\bigl[H((u_{v\to o}^s)^2-1)\bigr]
 =\E\bigl[H((\tau_{v\to o}^s)^2-1)\bigr].
\]
The variance estimate in \cite[Lemmas~3.2--3.4]{EAMS2023} is uniform
for the present normalizers. Indeed, if $K\ge1$ bounds the
coefficients at the fixed horizon and
$\Delta_r=\E[((\tau^r_{i\to j})^2-1)^2]$, its proof uses only
the variance recursion, independence of child cavities, and the exact
normalization $\E(A^r_{i\to j})^2=1$ to give
\[
 \Delta_{r+1}\le2K^4\Delta_r+\frac{2(K^4+1)}D,
 \qquad \Delta_0=0.
\]
Thus $\Delta_s=O(D^{-1})$, and the preceding conditional expectation
is $O(D^{-1/2})$ by Cauchy--Schwarz. Replacing the two node
coefficients by their cavity counterparts using
\eqref{eq:eams-node-cavity-bound} now yields
\begin{align*}
 \E\bigl[H(u_{v\to o}^s)^2\bigr]
 &=\E\bigl[A_{o\to v}^s(A_{v\to o}^{s-1})^2\bigr]
       +O(D^{-1/2})\\
 &=\E A_{o\to v}^s\,\E\bigl[(A_{v\to o}^{s-1})^2\bigr]
       +O(D^{-1/2})\\
 &=\E B_o^s+O(D^{-1/2}).
\end{align*}
The middle equality uses the exact independence of the two cavity
trees, and the last uses their coefficient normalization and one
more node-to-cavity comparison. This proves the quantitative
replacement without a convergence rate for $c_{D,r}$.
All constants here may depend on $\delta,L,\beta_0,T$, which are fixed.
Taking the degree limit and setting $r=\ell-1$
proves \eqref{eq:eams-corrected-value}.\footnote{Our index convention
is fixed by \eqref{eq:eams-output}.  The finite sums printed in
\cite[(4.5), (4.18)]{EAMS2023} use the shifted endpoints $2,\ldots,L$;
the preceding exact identity and Lemma~4.3 give $1,\ldots,L-1$.
The endpoint difference vanishes in the continuum limit below.
The exact identity we use is Proposition~2.1; the same-time display
in its Lemma~2.8 lacks the factor $1/k$ present in that proposition
and in the lemma's proof.}
Finally, the coefficients multiplying $U_r$ in
\eqref{eq:eams-limit-output} are measurable with respect to
$U_0,\ldots,U_{r-1}$.  Orthogonality of martingale increments and the
definition of $c_{\delta,r-1}$ give
$\E(Z^\delta_L)^2=\delta\sum_{r=1}^L1=L\delta$.
\end{proof}

The degree limit has identified the scalar law $Z^\delta_L$ and the
finite-tree edge limit in \eqref{eq:eams-corrected-value}.  We next
compare these two quantities with the Parisi diffusion at the fixed
temperature.  The symmetric-power argument in the following subsection
will then realize the same finite formula as a rooted vector, so that
its energy can be used in the clipping and synthesis steps.
The normalized Euler and martingale comparison below is related to
\cite[Proposition~5.3]{AMS2020}. We give the estimates for the present
fixed-temperature tree construction.

\begin{lemma}[The continuum limit at fixed temperature]
\label{lem:eams-positive-continuum}
Fix $\beta_0>1$ and $0<T<q$, and let $L=\lfloor T/\delta\rfloor$.
There is a coupling with the diffusion
\eqref{eq:positive-parisi-diffusion} and a constant $C=C(\beta_0,T)$
such that, for all sufficiently small $\delta$,
\begin{align}
 \max_{r\delta\le T}\E|X^\delta_r-X_{r\delta}|^2&\le C\delta,
 &\max_{r\delta\le T}|c_{\delta,r}-1|&\le C\sqrt\delta,
 \label{eq:euler-normalization-bound}\\
 \E|Z^\delta_L-m(T,X_T)|^2&\le C\delta,
 \label{eq:martingale-euler-bound}\\
 \left|\delta\sum_{r=1}^{L-1}c_{\delta,r}\E a(t_r,X^\delta_r)
                    -\int_0^T\E a(t,X_t)\,\dd t\right|
    &\le C\sqrt\delta.
 \label{eq:value-riemann-bound}
\end{align}
\end{lemma}

\begin{proof}
Take $U_0$ independent of $W$ and
$\sqrt\delta\,U_r=W_{r\delta}-W_{(r-1)\delta}$ for $r\ge1$.
Let $\tau_\delta(t)=\delta\lfloor t/\delta\rfloor$ and use the
continuous Euler interpolation
\[
 \overline X^\delta_t
 =\sqrt\delta\,U_0+
       \int_0^tb(\tau_\delta(s),X^\delta_{\lfloor s/\delta\rfloor})\,\dd s
       +W_t.
\]
The drift is bounded and Lipschitz in space and time on $[0,T]$.
Hence
$\E|\overline X^\delta_s-X^\delta_{\lfloor s/\delta\rfloor}|^2
\le C\delta$.
Subtracting the integral equations for $X$ and
$\overline X^\delta$, using Cauchy--Schwarz, and then Gronwall gives
\[
 \E\sup_{s\le T}|\overline X^\delta_s-X_s|^2\le C\delta.
\]
The initial contribution is $\E|\sqrt\delta U_0|^2=\delta$.
This proves the first bound of \eqref{eq:euler-normalization-bound}.

Because $a$ is bounded and spatially Lipschitz, and
$\E a(t_r,X_{t_r})^2=1$,
\[
 \max_{r\delta\le T}
 \left|\E a(t_r,X^\delta_r)^2-1\right|\le C\sqrt\delta.
\]
For small $\delta$ all these expectations lie in $[1/2,3/2]$.
The map $x\mapsto x^{-1/2}$ is Lipschitz on that interval, proving
the second bound of \eqref{eq:euler-normalization-bound}.

Equation~\eqref{eq:eams-limit-output} is the stochastic integral
\[
 Z^\delta_L=\int_0^{L\delta}
 c_{\delta,\lfloor t/\delta\rfloor}
 a(\tau_\delta(t),X^\delta_{\lfloor t/\delta\rfloor})\,\dd W_t.
\]
The integrand differs from $a(t,X_t)$ in squared mean by at most
$C\delta$, uniformly in $t\le L\delta$.  This follows from
\eqref{eq:euler-normalization-bound}, the spatial and temporal
Lipschitz bounds for $a$, and
$\E|X_t-X_{\tau_\delta(t)}|^2\le C\delta$.
It\^o isometry, including the last interval of length at most
$\delta$, gives \eqref{eq:martingale-euler-bound}, since
$m(T,X_T)=\int_0^Ta(t,X_t)\,\dd W_t$.
The same estimates in first mean compare the sum in
\eqref{eq:value-riemann-bound} with its Riemann integral.
The omitted first and last intervals have total length at most
$3\delta$ and bounded integrands.  This proves the final estimate.
\end{proof}

\subsection{Strong convergence of classical descendant formulas}

We next place this classical construction in the Hilbert space used
by Proposition~\ref{prop:energy}.  This step requires norm convergence,
not only a CLT in distribution.

We use the Gaussian lift $\mathsf G$ of
Section~\ref{subsec:classical-gaussian-sector} and the finite descendant
formulas of Definition~\ref{def:finite-descendant-formula}.
In the directed
EAMS calculation the sent columns are $A^{r-1}u^r$, with $u^0=g$
for the first propagation.  The node output uses these same columns.
Replacing $D+1$ root children by $D$, or the final normalization
$\sqrt{D+1}$ by $\sqrt D$, does not change the limiting formula.

\begin{lemma}[The symmetric-power realization]
\label{lem:classical-strong-limit}
Consider a fixed finite descendant formula whose pointwise functions
are smooth with all derivatives of polynomial growth.  Assume its
leaf seeds have moments of every order and every aggregated child
column has mean zero under its finite-degree child law for each $D$.
Its finite-degree root vector converges in
norm, under the symmetric occupation embeddings, to the formula
obtained by replacing every normalized sum by $\mathsf G$.
The same conclusion holds for a convergent finite collection of
deterministic coefficients in the formula.  The finite-degree root
vectors have a uniformly bounded first child-occupation moment.
Consequently, for centered equivariant outputs $f_D$ and their limit
$f$,
\begin{equation}
 \lim_{D\to\infty}\frac{D+1}{2\sqrt D}\E f_D(o)f_D(v)
       =\cE(f),\qquad o\sim v.
 \label{eq:classical-energy-transfer}
\end{equation}
\end{lemma}

\begin{proof}
The proof first treats characteristic vectors, then arbitrary
pointwise maps, and finally the occupation bound needed for energy.
We use induction on formula depth.  All intermediate
moments are uniformly bounded: normalized centered sums satisfy the
moment bound used in Lemma~\ref{lem:eams-finite-step}, and
polynomial-growth functions preserve the existence of sufficiently
high uniform moments.

Suppose the real centered child columns
$f_{D,1},\ldots,f_{D,k}$ already converge in norm to
$f_1,\ldots,f_k$ in the common child Hilbert space.  Fix
$t\in\R^k$ and put $F_D=\sum_jt_jf_{D,j}$,
$F=\sum_jt_jf_j$. The limits are centered because the embeddings
preserve the vacuum and the convergence is in norm.
A joint characteristic vector for the parent sums
is the symmetric tensor power of
\[
 h_D=\exp(\ii F_D/\sqrt D)\vac
          =\alpha_D\vac+\zeta_D,
 \qquad \zeta_D\perp\vac.
\]
Taylor expansion in $L^2$, justified by the uniform fourth moment,
and centering give
\[
 \alpha_D=1-\frac{\norm{F_D}^2}{2D}+O(D^{-3/2}),
 \qquad \sqrt D\,\zeta_D\longrightarrow\ii F.
\]
The $j$-particle component of the embedded tensor power is
\[
 \binom Dj^{1/2}\alpha_D^{D-j}\zeta_D^{\otimes j}
 \longrightarrow
 e^{-\norm F^2/2}\frac{(\ii F)^{\otimes j}}{\sqrt{j!}}.
\]
These are the components of $e^{\ii q(F)}\vac$.
The expected occupation of the left-hand vector is
$D\norm{\zeta_D}^2=O(1)$, so its occupation tails are uniformly
small.  Thus convergence of each fixed grade implies norm convergence
of the complete vector.  A root with $M=D+1$ children has the same
proof: use $M$ factors and the normalization $M^{-1/2}$, with the
same convergent child columns.  Alternatively, retaining $D^{-1/2}$
at that root gives the same limit because $M/D\to1$.  The occupation
bound in either convention is $M\norm{\zeta_D}^2=O(1)$.

Finite linear combinations of these exponentials, multiplied by
bounded measurable functions of the independent root seed, are dense
in the limiting $L^2$ space.  Indeed the root seed and aggregate
fields have a product law.  Products of dense families in the two
factors are dense, and exponentials are dense in the field factor:
an orthogonal $L^2$ vector would give a finite signed measure with
zero Fourier transform and would therefore vanish.
There is a useful precision concerning measurable seed factors.
Let $S$ denote the root seed, $Y_D$ the finite list of aggregate
fields, and $Y$ its limit.  The law of $S$ is unchanged with $D$ and
is independent of $Y_D$ and $Y$.  For
$P(s,y)=\sum_jg_j(s)e^{\ii t_j\cdot y}$ with bounded measurable
$g_j$, and a bounded target $F(s,y)$ continuous in $y$ for every
fixed $s$, the function $|F(s,y)-P(s,y)|^2$ is bounded continuous in
$y$ for each fixed $s$.  Weak convergence of $Y_D$ therefore gives
convergence of its field integral pointwise in $s$.  The bound is
uniform in $s$, so dominated convergence in the unchanged root-seed
law proves
\[
 \E|F(S,Y_D)-P(S,Y_D)|^2
       \longrightarrow\E|F(S,Y)-P(S,Y)|^2.
\]
Thus measurability of $g_j$ does not require a discontinuous test
function in an appeal to joint weak convergence.  For a
polynomial-growth target, first multiply by smooth cutoffs outside
large boxes in the seed and field variables.  The cutoffs preserve
continuity in the field variables, and uniform higher moments make
the discarded squared norms uniformly small.  This proves the induction
step.  It also proves
continuity in deterministic coefficients, by applying the same
argument to a convergent sequence of them.

For completeness, the occupation bound is not a consequence of the
CLT alone.  At a parent, let $P_j$ average over the $j$th child, and
write its output as $H_D(s,S_{D,1},\ldots,S_{D,k})$, where
$S_{D,l}=D^{-1/2}\sum_j f_{D,l}^{(j)}$.
The occupation operator in the symmetric embedding is
$N_D=\sum_{j=1}^D(I-P_j)$.  With an independent replacement of child
$j$, conditional variance and the mean-value theorem give
\[
 \norm{(I-P_j)H_D}^2
 \le\E|H_D-H_D^{(j)}|^2\le\frac{C}{D}.
\]
Here the derivative of $H_D$ has polynomial growth and the sums,
the replaced child columns, and their products have the uniform
moments established above; H\"older's inequality therefore gives a
constant independent of $j,D$.  Summing proves
$\braket{H_D}{N_DH_D}\le C$.  If needed, the same argument with
distinct child replacements and mixed derivatives bounds the
factorial moments of $N_D$ of every fixed order.

The energy identity is now Proposition~\ref{prop:energy} applied to
these classical vectors.  More explicitly, with $M=D+1$ root
children, the projection onto an excitation in a distinguished child
has norm at most $C/\sqrt M$ by symmetry and the occupation estimate.
The term with excitations across the distinguished edge on both sides
is therefore $O(M^{-1})$.  The two terms with only one crossing
select a single child; after multiplication by $\sqrt M$ they
converge to $\braket{f}{\cL f}$ and its conjugate.  The term with
neither crossing is zero by centering.  The factor
$M/(2\sqrt D)$ in \eqref{eq:classical-energy-transfer} thus gives
$\re\braket{f}{\cL f}$, since $M/D\to1$.
\end{proof}

Apply this lemma to the EAMS formulas at fixed $\delta,L$.
All sent columns are odd under simultaneous reversal of the seed
labels, since $a$ is even and $b$ is odd. They are therefore centered
under each finite-degree child law, as required by the lemma.
The smooth polynomial-growth hypotheses hold because a transmitted
column is a bounded smooth coefficient times a message
variable.  Equation~\eqref{eq:eams-normalizer-limit} and the
coefficient-continuity assertion allow us to replace $c_{D,r}$ by the
fixed limiting constants $c_{\delta,r}$ in the limiting formula.
Let its output vector be $z^\delta$.  We obtain
\begin{equation}
 \cE(z^\delta)=\delta\sum_{r=1}^{L-1}
                    c_{\delta,r}\E a(t_r,X^\delta_r),
 \quad
 \norm{z^\delta}^2=L\delta,
 \quad
 z^\delta\ \hbox{has law }Z^\delta_L.
 \label{eq:classical-eams-vector}
\end{equation}
All coefficients in this vector are independent of degree.

\subsection{Clipping and finite bounded formulas}

Let $\pi(x)=\max(-1,\min(x,1))$.  Since $|m(T,X_T)|\le1$,
the law in \eqref{eq:classical-eams-vector} and the metric projection
property give
\[
 \norm{z^\delta-\pi(z^\delta)}^2
 =\E|Z^\delta_L-\pi(Z^\delta_L)|^2
 \le\E|Z^\delta_L-m(T,X_T)|^2.
\]
Only the marginal law of $Z^\delta_L$ enters the equality. The
diffusion coupling bounds this scalar clipping error; the energy of
$z^\delta$ has already been identified separately in the rooted space
by \eqref{eq:classical-eams-vector}.
Consequently, by Lemma~\ref{lem:eams-positive-continuum},
$\norm{z^\delta-\pi(z^\delta)}\le C\sqrt\delta$.
We apply the shared energy-continuity estimate
\eqref{eq:energy-l2-continuity} with $f=z^\delta$ and
$g=\pi(z^\delta)$; both norms are at most one.  Combining
\eqref{eq:positive-temperature-bound},
\eqref{eq:value-riemann-bound}, and
\eqref{eq:classical-eams-vector} proves the following: for every
$\varepsilon>0$, first choosing a sufficiently large finite
$\beta_0$, then $T<q$ sufficiently close to $q$, and finally a
sufficiently small positive $\delta$, gives a centered odd classical
vector
\begin{equation}
 m^*=\pi(z^\delta),\qquad |m^*|\le1,
 \qquad \cE(m^*)\ge P_*-\varepsilon.
 \label{eq:clipped-classical-vector}
\end{equation}
At this point the depth $L$ and the entire descendant calculation are
finite and fixed.

The next proposition completes this route using only the common
bounded-approximation and exclusion lemmas.  It gives the same
conclusion as Proposition~\ref{prop:classical-near-optimum}, with an
independent proof from the clipped vector just constructed.

\begin{proposition}[The independent positive-temperature construction]
\label{prop:classical-near-optimum-positive-temperature}
For every $\varepsilon>0$ and every $\eta_H>0$, there are a finite
color partition with helper probability $0<p_H<\eta_H$, a finite
number of helper types, and finite bounded odd angle formulas
$\theta_c$ on the nonhelper colors such that the following hold.
Each transmitted column is bounded and odd, every pointwise map
may be chosen as a finite trigonometric sum in its odd inputs, and
every aggregation excludes its current color, all path-ancestor
colors, and all helpers.  Conditional on the finite labels, all
continuous derivatives of the pointwise maps are bounded.  The
output
\begin{equation}
 m^{\mathrm{col}}
  =\sum_{c=1}^M\1_{\{C=c\}}\sin(2\theta_c)
 \label{eq:positive-final-classical-formula}
\end{equation}
is centered, satisfies $|m^{\mathrm{col}}|\le1$, vanishes on helper
roots, and obeys
\begin{equation}
 \cE(m^{\mathrm{col}})\ge P_*-\varepsilon.
 \label{eq:positive-final-classical-value}
\end{equation}
After the uncolored finite formula has been fixed, the total helper
probability may be taken arbitrarily small, with all individual
helper probabilities positive.  All depths, coefficients, label laws,
and functions in the conclusion are fixed independently of degree,
system size, and disorder.
\end{proposition}

\begin{proof}
Use \eqref{eq:clipped-classical-vector} with error
$\varepsilon/3$.  Its underlying formula $z^\delta$ has odd
transmitted columns and jointly odd smooth local maps with
polynomial-growth derivatives, as checked after
Lemma~\ref{lem:classical-strong-limit}.  It therefore meets the
hypotheses of Lemma~\ref{lem:classical-trigonometric-density}.
Apply that lemma and then Lemma~\ref{lem:classical-color-exclusion}
with the choices in the proof of
Proposition~\ref{prop:classical-near-optimum}: the first output error
is less than $\varepsilon/12$ in norm, the second less than
$\varepsilon/6$, and one positive-probability helper is reserved per
propagation depth with total probability below $\eta_H$.
The two further energy losses are less than $\varepsilon/6$ and
$\varepsilon/3$, respectively, by
\eqref{eq:energy-l2-continuity}.  Their sum with the initial loss is
less than $\varepsilon$.  Those two shared lemmas also give the
boundedness, trigonometric form, exclusions, and parity asserted here.
The entire formula and all helper probabilities are fixed after the
degree limit of the finite calculation.
\end{proof}

\begin{thebibliography}{99}
\addcontentsline{toc}{section}{\refname}
\small
\raggedright

\bibitem{AC2015}
A.~Auffinger and W.-K.~Chen,
\emph{On properties of Parisi measures},
Probability Theory and Related Fields \textbf{161} (2015), 817--850.
\href{https://doi.org/10.1007/s00440-014-0563-y}{doi:10.1007/s00440-014-0563-y};
\href{https://arxiv.org/abs/1303.3573v1}{arXiv:1303.3573v1}.

\bibitem{AC2017}
A.~Auffinger and W.-K.~Chen,
\emph{Parisi formula for the ground state energy in the mixed $p$-spin model},
Annals of Probability \textbf{45} (2017), no.~6B, 4617--4631.
\href{https://doi.org/10.1214/16-AOP1173}{doi:10.1214/16-AOP1173};
\href{https://arxiv.org/abs/1606.05335v2}{arXiv:1606.05335v2}.

\bibitem{ACZ2020}
A.~Auffinger, W.-K.~Chen, and Q.~Zeng,
\emph{The SK Model Is Infinite Step Replica Symmetry Breaking at Zero
Temperature},
Communications on Pure and Applied Mathematics \textbf{73} (2020), no.~5,
921--943.
\href{https://doi.org/10.1002/cpa.21886}{doi:10.1002/cpa.21886};
\href{https://arxiv.org/abs/1703.06872v2}{arXiv:1703.06872v2}.

\bibitem{BFMVZ2022}
J.~Basso, E.~Farhi, K.~Marwaha, B.~Villalonga, and L.~Zhou,
\emph{The Quantum Approximate Optimization Algorithm at High Depth for
MaxCut on Large-Girth Regular Graphs and the Sherrington--Kirkpatrick Model},
17th Conference on the Theory of Quantum Computation, Communication and
Cryptography, LIPIcs \textbf{232} (2022), 7:1--7:21.
\href{https://doi.org/10.4230/LIPIcs.TQC.2022.7}{doi:10.4230/LIPIcs.TQC.2022.7};
full version \href{https://arxiv.org/abs/2110.14206v3}{arXiv:2110.14206v3}.

\bibitem{BKL2026}
S.~Boulebnane, A.~Khan, M.~Liu, J.~Larson, D.~Herman, R.~Shaydulin,
and M.~Pistoia,
\emph{Spin--Boson Mapping of the Quantum Approximate Optimization Algorithm},
Physical Review Letters \textbf{136} (2026), 240601.
\href{https://doi.org/10.1103/2w94-rymn}{doi:10.1103/2w94-rymn};
\href{https://arxiv.org/abs/2505.07929v2}{arXiv:2505.07929v2}.

\bibitem{Chen2026}
H.-B.~Chen,
\emph{FRSB in the SK spin glass: convergence to full-interval support
at zero temperature},
\href{https://arxiv.org/abs/2607.18032v1}{arXiv:2607.18032v1} (2026).

\bibitem{DMS2017}
A.~Dembo, A.~Montanari, and S.~Sen,
\emph{Extremal Cuts of Sparse Random Graphs},
Annals of Probability \textbf{45} (2017), no.~2, 1190--1217.
\href{https://doi.org/10.1214/15-AOP1084}{doi:10.1214/15-AOP1084};
\href{https://arxiv.org/abs/1503.03923v2}{arXiv:1503.03923v2}.

\bibitem{AMS2020}
A.~El Alaoui, A.~Montanari, and M.~Sellke,
\emph{Optimization of mean-field spin glasses},
Annals of Probability \textbf{49} (2021), no.~6, 2922--2960.
\href{https://doi.org/10.1214/21-AOP1519}{doi:10.1214/21-AOP1519};
full version \href{https://arxiv.org/abs/2001.00904v1}{arXiv:2001.00904v1} (2020).

\bibitem{EAMS2023}
A.~El Alaoui, A.~Montanari, and M.~Sellke,
\emph{Local algorithms for Maximum Cut and Minimum Bisection on locally
treelike regular graphs of large degree},
Random Structures \& Algorithms \textbf{63} (2023), no.~3, 689--715.
\href{https://doi.org/10.1002/rsa.21149}{doi:10.1002/rsa.21149};
\href{https://arxiv.org/abs/2111.06813v2}{arXiv:2111.06813v2}.

\bibitem{FGG2014}
E.~Farhi, J.~Goldstone, and S.~Gutmann,
\emph{A Quantum Approximate Optimization Algorithm},
\href{https://arxiv.org/abs/1411.4028v1}{arXiv:1411.4028v1} (2014).

\bibitem{FGGZ2022}
E.~Farhi, J.~Goldstone, S.~Gutmann, and L.~Zhou,
\emph{The Quantum Approximate Optimization Algorithm and the
Sherrington--Kirkpatrick Model at Infinite Size},
Quantum \textbf{6} (2022), 759.
\href{https://doi.org/10.22331/q-2022-07-07-759}{doi:10.22331/q-2022-07-07-759};
\href{https://arxiv.org/abs/1910.08187v4}{arXiv:1910.08187v4}.

\bibitem{Guerra2003}
F.~Guerra,
\emph{Broken Replica Symmetry Bounds in the Mean Field Spin Glass Model},
Communications in Mathematical Physics \textbf{233} (2003), no.~1, 1--12.
\href{https://doi.org/10.1007/s00220-002-0773-5}{doi:10.1007/s00220-002-0773-5};
\href{https://arxiv.org/abs/cond-mat/0205123v1}{arXiv:cond-mat/0205123v1}.

\bibitem{JT2017}
A.~Jagannath and I.~Tobasco,
\emph{Some properties of the phase diagram for mixed $p$-spin glasses},
Probability Theory and Related Fields \textbf{167} (2017), 615--672.
\href{https://doi.org/10.1007/s00440-015-0691-z}{doi:10.1007/s00440-015-0691-z};
\href{https://arxiv.org/abs/1504.02731v2}{arXiv:1504.02731v2}.

\bibitem{JonesKnill1999}
J.~A.~Jones and E.~Knill,
\emph{Efficient Refocussing of One Spin and Two Spin Interactions for NMR
Quantum Computation},
Journal of Magnetic Resonance \textbf{141} (1999), 322--325.
\href{https://doi.org/10.1006/jmre.1999.1890}{doi:10.1006/jmre.1999.1890};
\href{https://arxiv.org/abs/quant-ph/9905008v1}{arXiv:quant-ph/9905008v1}.

\bibitem{LiKhaneja2005}
J.-S.~Li and N.~Khaneja,
\emph{Noncommuting vector fields, polynomial approximations and control of
inhomogeneous quantum ensembles},
\href{https://arxiv.org/abs/quant-ph/0510012v1}{arXiv:quant-ph/0510012v1} (2005).

\bibitem{Lopatto2026}
P.~Lopatto,
\emph{Full replica symmetry breaking in the Sherrington--Kirkpatrick model},
\href{https://arxiv.org/abs/2607.11756v3}{arXiv:2607.11756v3} (2026).

\bibitem{MWW2004}
B.~D.~McKay, N.~C.~Wormald, and B.~Wysocka,
\emph{Short Cycles in Random Regular Graphs},
Electronic Journal of Combinatorics \textbf{11} (2004), no.~1, R66.
\href{https://doi.org/10.37236/1819}{doi:10.37236/1819}.

\bibitem{Montanari2019}
A.~Montanari,
\emph{Optimization of the Sherrington--Kirkpatrick Hamiltonian},
SIAM Journal on Computing \textbf{54} (2025), no.~4, FOCS19-1--FOCS19-38.
\href{https://doi.org/10.1137/20M132016X}{doi:10.1137/20M132016X};
full version \href{https://arxiv.org/abs/1812.10897v2}{arXiv:1812.10897v2} (2019).

\bibitem{SKSupport2026}
OpenAI,
\emph{Full support of the zero-temperature Sherrington--Kirkpatrick
order parameter},
OpenAI Math Release preprint
\href{https://github.com/openai/math/blob/main/preprints/Full-support-of-the-zero-temperature-Sherrington-Kirkpatrick-order-parameter-September-27-2026/main.pdf}{OAI:Full-support-of-the-zero-temperature-Sherrington-Kirkpatrick-order-parameter-September-27-2026},
2026.

\bibitem{Parisi1979}
G.~Parisi,
\emph{Infinite Number of Order Parameters for Spin-Glasses},
Physical Review Letters \textbf{43} (1979), no.~23, 1754--1756.
\href{https://doi.org/10.1103/PhysRevLett.43.1754}{doi:10.1103/PhysRevLett.43.1754}.

\bibitem{SM2026}
D.~Sels and F.~Morone,
\emph{Absence of quantum advantage for approximate spin glass optimization},
\href{https://arxiv.org/abs/2607.08708v1}{arXiv:2607.08708v1} (2026).

\bibitem{SK1975}
D.~Sherrington and S.~Kirkpatrick,
\emph{Solvable Model of a Spin-Glass},
Physical Review Letters \textbf{35} (1975), no.~26, 1792--1796.
\href{https://doi.org/10.1103/PhysRevLett.35.1792}{doi:10.1103/PhysRevLett.35.1792}.

\bibitem{Talagrand2006}
M.~Talagrand,
\emph{The Parisi formula},
Annals of Mathematics \textbf{163} (2006), no.~1, 221--263.
\href{https://doi.org/10.4007/annals.2006.163.221}{doi:10.4007/annals.2006.163.221}.

\end{thebibliography}

\end{document}
